\documentclass[11pt]{article}
\usepackage[T1]{fontenc}
\usepackage{lmodern,amsmath,amssymb,amsthm,amscd,mathrsfs,mathtools,booktabs,longtable,array,microtype}
\usepackage[margin=1.05in]{geometry}
\usepackage[hidelinks]{hyperref}
\newtheorem{theorem}{Theorem}[section]
\newtheorem{proposition}[theorem]{Proposition}
\newtheorem{lemma}[theorem]{Lemma}
\newtheorem{corollary}[theorem]{Corollary}
\theoremstyle{definition}
\newtheorem{definition}[theorem]{Definition}
\newtheorem{example}[theorem]{Example}
\theoremstyle{remark}
\newtheorem{remark}[theorem]{Remark}
\newcommand{\Q}{\mathbf Q}
\newcommand{\Z}{\mathbf Z}
\newcommand{\R}{\mathbf R}
\newcommand{\F}{\mathbf F}
\newcommand{\PP}{\mathbf P}
\newcommand{\OO}{\mathcal O}
\newcommand{\Jac}{\operatorname{Jac}}
\newcommand{\Pic}{\operatorname{Pic}}
\newcommand{\Sel}{\operatorname{Sel}}
\newcommand{\Nm}{\operatorname{N}}
\newcommand{\ord}{\operatorname{ord}}
\newcommand{\rank}{\operatorname{rank}}
\newcommand{\Gal}{\operatorname{Gal}}
\renewcommand{\sf}{\operatorname{sf}}
\newcommand{\leg}[2]{\left(\frac{#1}{#2}\right)}
\newcommand{\eps}{\varepsilon}
\hypersetup{pdftitle={Arithmetic constraints and covering curves for rational distances in a square},pdfauthor={Yuan Si}}
\title{Arithmetic constraints and covering curves\newline for rational distances in a square}
\author{Yuan Si\\\small University of Waterloo\\\small\texttt{y3si@uwaterloo.ca}}
\date{Version 1.1.0}
\setlength{\emergencystretch}{1em}
\begin{document}
\maketitle
\begin{abstract}
We study points in the strict interior of a square whose four vertex
 distances are rational. Primitive integral and Gaussian models yield exact
 denominator identities, prime-power bounds, and reciprocity restrictions
 on the squarefree classes assigned to opposite corners. We construct
 ordered elliptic sources and covering curves of genera five, nine and
 seventeen, with explicit maps retaining the four distances. A common
 squareclass descent gives infinite families of genus-nine curves without
 rational points, while a compensated source descent excludes an entire
 squarefree kernel. Explicit quaternion evaluations give a parity
 obstruction and an infinite family of complete genus-five failures of the
 Hasse principle. We describe nine elliptic
 factors of the genus-nine Jacobian and derive rank constraints. For the
 fixed opposite-distance difference $1/5$, we determine all twenty-eight
 bitangents of a genus-three quotient, its odd-theta algebra and the theta
 class defining the original double cover. A height comparison determines
 all rational points of this quotient and excludes the complete one-fifth
 source. We also give normalized symmetric return maps, pencils of smooth
 plane quartics with explicit elliptic factors, a uniform dyadic
 independence criterion, and exact signed-sum valuations at denominator
 primes. In particular, a primitive radius tuple of content one requires
 at least three denominator sign patterns. Finally, we prove that rational
 boundary points and their real components produce locally soluble lifts
 at every fixed finite Abel--Jacobi level.
\end{abstract}
\clearpage
\tableofcontents
\clearpage
\section{Introduction}
\label{sec:introduction}

The rational-distance problem for a square asks for a point whose distances
from all four vertices are rational. We consider the strict interior of the
unit square and retain the four distances as part of the arithmetic data.
Writing the vertices in the order $00,10,01,11$, the equations are
\[
 \begin{aligned}
 r_0^2&=x^2+y^2,&r_1^2&=(1-x)^2+y^2,\\
 r_2^2&=x^2+(1-y)^2,&r_3^2&=(1-x)^2+(1-y)^2,
 \end{aligned}
 \qquad 0<x,y<1,
\]
with $r_i\in\Q_{>0}$. Differences of the squared distances recover $x$
and $y$, so this is a rational-point problem with a specified real chamber.

The geometry of rational-distance varieties has been studied from several
directions. Berry relates the square to a Kummer surface and a branched
double cover, and studies associated curves of genus five
\cite{berry1990}; his work on triangles develops the corresponding
three-distance geometry \cite{berry1992}. Bremner and Ulas give
constructions for rational-distance problems attached to several geometric
objects \cite{bremnerulas2016}. Corvaja, Turchet and Zannier study distance
varieties for general configurations of rational points, including their
surfaces and product-quotient descriptions
\cite{corvajaturchetzannier2025}. The present paper develops the arithmetic
of explicit covering curves, the squareclasses carried by one ordered
source, and congruences that hold at every denominator height.

The first part concerns the primitive integral equations. Let $L$ be the
least common denominator of $x,y$, and write $A_i=Lr_i$. We prove that
$24\mid L$, that the common radius content $G=\gcd(A_0,A_1,A_2,A_3)$
satisfies $G^2\mid L$, and that the six reduced denominators of the pairwise
radius sums and differences satisfy exact gcd and product identities
(Theorem~\ref{thm:int:pair-denominators}). Signed radius identities compare
each prime-power divisor of $L$ with the size of the square. They imply,
in particular, that $L$ has a prime divisor other than $2$ or $3$
(Corollary~\ref{cor:int:six-smooth}).

Gaussian factorization assigns four canonical, pairwise coprime
squarefree integers $m_0,\ldots,m_3$ to the corners. Their product is
the squarefree part of
\[
 (r_0r_3+r_1r_2)^2-1.
\]
All primes in these integers are $1$ or $7$ modulo $8$. A reciprocity
identity for each pair of corners retains the common denominator and the
sign of the Gaussian norm. For the two opposite pairs it gives
\[
 m_0m_3\equiv1\text{ or }15\pmod{16},\qquad
 m_1m_2\equiv1\text{ or }15\pmod{16}
\]
(Theorem~\ref{thm:profile-main}). We also classify the residue pattern
arising when $v_p(L)=2v_p(G)$. An elliptic residue curve identifies the
exceptional primes $5$ and $13$; imposing the additional product-square
condition gives the refined exceptional set $\{5,13,17\}$
(Theorems~\ref{thm:crit-basic} and~\ref{thm:crit-refined}).

The second part gives explicit rational correspondences. An ordered pair
$Q=(a,b)$, $R=(c,e)$ on
\[
 E_n:\quad Y^2=X(X-n)(X+n)
\]
encodes the original point through two ordered product squareclasses,
the two ordinates, and the single further square
\[
 \frac{a+c}{2(c+n)}.
\]
Theorem~\ref{thm:ell-exact-source} gives the forward and inverse maps,
including all signs and strict inequalities. Fixing a signed sum or
difference of two distances instead gives a smooth curve of genus five.
The resulting fixed-pair finiteness theorem follows from the explicit
branch calculation and Faltings' theorem
(Theorem~\ref{thm:ell-fixed-pair-finite}).

The product parameter $t=r_0r_3+r_1r_2$ leads to a genus-nine curve $D_t$
with four ordered square roots over a common coordinate $\ell$. Its
valuation conditions constrain one global class $[\ell]$. A generally
nonsymmetric Legendre-symbol matrix then reduces the possible classes to
two, and a local identity at $37$ excludes both under an explicit residue
condition. This proves an infinite family of all-height no-point results
(Theorem~\ref{thm:common-squareclass} and
Corollary~\ref{cor:common-infinite}). The associated unramified double
covers have genus seventeen. Nine explicit elliptic quotients describe
$\Jac(D_t)$ and yield rank constraints on parameters admitting rational
points (Theorems~\ref{thm:nine-isogeny}
and~\ref{thm:nine-rank-strata}).

Further descent must compare local data arising from the same global
source. We express this requirement through a character map and a
secondary obstruction in its cokernel. Keeping the complete compensation
coset gives the criterion of Theorem~\ref{thm:cs-matrix}. A short
arithmetic certificate applies it, together with a higher quotient, to
exclude the entire kernel $n=15569$
(Corollary~\ref{cor:cs-15569-layer}). A separate construction attaches
a proper quaternion class to a complete genus-five source. Its local
evaluations give an explicit parity criterion through Hilbert reciprocity
(Theorem~\ref{thm:hp-uniform}). An explicit conic parametrization yields
an infinite family of complete genus-five failures of the Hasse principle
(Theorem~\ref{thm:hp-infinite}).

The theta section treats the fixed opposite-distance difference $1/5$
in detail. Its genus-five source has an unramified genus-three quotient
which is a smooth plane quartic. We determine all twenty-eight
bitangents, their contact divisors and Galois action, obtaining the
odd-theta algebra
\[
 \Q^8\times\Q(\sqrt2)^4\times\Q(\sqrt{51})^4
 \times\Q(\sqrt2,\sqrt{51}).
\]
The calculation supplies explicit theta fractions of the type used in
generalized descent \cite{bruinpoonenstoll2016,misme2026}. A polynomial
identity identifies the original double cover with a specified difference
of theta characteristics. We then prove a local-solubility theorem for
every finite Abel--Jacobi multiplication cover through a rational boundary
point. Its application preserves the original double cover and a strict
real point (Corollary~\ref{cor:theta-surviving-covers}).

The symmetric coordinates of Section~\ref{sec:symmetric-models} give a
normalized return in two rational parameters and two pencils of smooth
plane quartics. Their Jacobians split, up to isogeny, into three elliptic
factors. The multiplicity-two factor changes with the common distance
squareclass. A dyadic comparison of the ordered elliptic sources then
shows that every strict rational-distance point forces two independent
points on its congruent-number curve
(Corollary~\ref{cor:indp-original}).

Section~\ref{sec:denominator-patterns} records the signs of the four
radii at every prime dividing $L/G^2$. The valuations of all eight signed
sums determine square or half-square denominators for
$r_0r_3+r_1r_2$ and $r_0r_3-r_1r_2$. Combining these valuations with real
size bounds shows that, when $G=1$, at least three sign patterns occur.
It also gives a sharper dyadic bound and a split-prime requirement when
only one corner colour is nontrivial.

Finally, Section~\ref{sec:one-fifth-points} completes the rational-point
calculation for the one-fifth source. Two maps to the same elliptic curve,
an explicit canonical-height comparison and a finite table of nonsquare
tests give exactly four rational points on its plane-quartic quotient
and eight boundary points on the genus-five curve. Thus the entire
opposite-distance difference $\pm1/5$ is excluded at every rational height.

For the models encoding the original four-distance configurations, we give
explicit reconstruction formulas and the inequalities defining the strict
real chamber. Finite residue calculations
are given as explicit identities, congruences or complete small tables.
The elliptic group needed in the one-fifth calculation is supplied by an
explicit rational two-descent and full-saturation certificate. We use the
classical torsion theorem of Mazur, the finiteness theorem of
Faltings, and standard facts about elliptic curves, covers and Brauer groups
in the statements where they enter. The appendix describes the compact
four-radius surface and the Brauer evaluations associated with its real
components.

\paragraph{Companion material.}
The paper sources and a partial Lean formalization are archived at
\url{https://doi.org/10.5281/zenodo.23638744}.
The archive identifies the verified statements, pinned build environment,
and exact arithmetic certificates. The Lean companion verifies 59 of the
original 69 numbered statements. The additional statements of this
version have written proofs; their formalization status is listed
separately in the accompanying coverage table.

\section{Integral and Gaussian forms of the four-distance equations}
\label{sec:integral}

We label the vertices of the unit square by
\[
 P_0=(0,0),\qquad P_1=(1,0),\qquad
 P_2=(0,1),\qquad P_3=(1,1).
\]
A \emph{strict rational-distance point} is a point $(x,y)$ with
$0<x,y<1$ whose distances $r_i$ to all four $P_i$ belong to
$\Q_{>0}$. Thus
\begin{equation}\label{eq:int:four}
\begin{aligned}
 r_0^2&=x^2+y^2,&r_1^2&=(1-x)^2+y^2,\\
 r_2^2&=x^2+(1-y)^2,&r_3^2&=(1-x)^2+(1-y)^2.
\end{aligned}
\end{equation}
The labels in these equations will be retained throughout.

\subsection{Recovering the point from its four radii}

Subtracting adjacent equations gives
\begin{equation}\label{eq:int:recover}
 2x=1+r_0^2-r_1^2,\qquad
 2y=1+r_0^2-r_2^2.
\end{equation}
In particular the coordinates are rational. The following elimination
keeps both the point and its four positive radii.

\begin{proposition}\label{prop:int:radius-model}
Positive rational numbers $(r_0,r_1,r_2,r_3)$ arise from a unique strict
rational-distance point if and only if
\begin{align}
 r_0^2+r_3^2&=r_1^2+r_2^2,\label{eq:int:flag}\\
 (1+r_0^2-r_1^2)^2+(1+r_0^2-r_2^2)^2&=4r_0^2,\label{eq:int:radiusquartic}\\
 -1<r_0^2-r_1^2&<1,\qquad
 -1<r_0^2-r_2^2<1.\label{eq:int:strict-radii}
\end{align}
The point is given by \eqref{eq:int:recover}.
\end{proposition}

\begin{proof}
Equations \eqref{eq:int:flag} and \eqref{eq:int:radiusquartic} follow
by adding and subtracting \eqref{eq:int:four}. Conversely, define $x,y$
by \eqref{eq:int:recover}. The inequalities give $0<x,y<1$, and
\eqref{eq:int:radiusquartic} gives $x^2+y^2=r_0^2$. Subtracting $2x-1$
and $2y-1$ gives the equations for $r_1,r_2$; the remaining equation
then follows from \eqref{eq:int:flag}. Positivity fixes the four square
roots, and \eqref{eq:int:recover} gives uniqueness.
\end{proof}

Let $L$ be the least positive common denominator of $x,y$, and put
\begin{equation}\label{eq:int:lattice}
 X=Lx,\qquad Y=Ly,\qquad A_i=Lr_i.
\end{equation}
Then $X,Y,L$ are integers with
\[
 0<X,Y<L,\qquad \gcd(X,Y,L)=1.
\]
Each $A_i$ is an integer: a rational number whose square is an integer
is itself an integer. Equations \eqref{eq:int:four} become
\begin{equation}\label{eq:int:integer-four}
\begin{aligned}
 A_0^2&=X^2+Y^2,&A_1^2&=(L-X)^2+Y^2,\\
 A_2^2&=X^2+(L-Y)^2,&A_3^2&=(L-X)^2+(L-Y)^2.
\end{aligned}
\end{equation}
Conversely, positive integers satisfying these equations and
$0<X,Y<L$ give a strict rational-distance point by division by $L$;
the additional gcd condition makes $L$ its least coordinate denominator.

\subsection{Parity, symmetry lines, and the primitive denominator}

We first record the elementary descent needed to deal with the symmetry
lines of the square.

\begin{lemma}[Square-area descent]\label{lem:int:fermat}
A rational right triangle cannot have positive rational-square area.
\end{lemma}

\begin{proof}
Clearing denominators would give an integer right triangle of square
area. Dividing by the gcd of its sides preserves this property: the
new area is an integer and a rational square, hence an integer square.
Choose a primitive such triangle with least hypotenuse. Its legs and
hypotenuse have the form
\[
 2rs,\quad r^2-s^2,\quad r^2+s^2,
 \qquad r>s>0,\quad \gcd(r,s)=1,
\]
where $r,s$ have opposite parity. To obtain this parametrization, write
the odd leg as $a$ and the hypotenuse as $h$. The coprime positive
integers $(h+a)/2,(h-a)/2$ have square product and are therefore
squares.

The area $rs(r-s)(r+s)$ is a square. Its four factors are pairwise
coprime, so
\[
 r=u^2,\qquad s=v^2,\qquad r+s=w^2,\qquad r-s=z^2.
\]
The numbers $w,z$ are odd. Since $w^2-z^2=2v^2$ and odd squares are
$1$ modulo $8$, the integer $v$ is even. The positive integers
\[
 a'=(w+z)/2,\qquad b'=(w-z)/2
\]
satisfy
\[
 (a')^2+(b')^2=u^2,\qquad a'b'/2=(v/2)^2.
\]
They give a right triangle of square area with hypotenuse $u$ strictly
smaller than $r^2+s^2$. Primitive reduction, if required, decreases
the hypotenuse further. This contradicts its minimal choice.
\end{proof}

\begin{lemma}\label{lem:int:symmetry}
A strict rational-distance point lies on neither a midline nor a
diagonal of the square. Its four radii are pairwise distinct.
\end{lemma}

\begin{proof}
On $y=1/2$, put $u=2x-1$, $p=2r_0$, and $q=2r_1$. Then
\[
 p^2=u^2+2u+2,\qquad q^2=u^2-2u+2,
 \qquad (pq)^2=u^4+4.
\]
If $u\ne0$, the rational numbers $|u|,2/|u|,pq/|u|$ are the sides
of a right triangle of area one, contrary to
Lemma~\ref{lem:int:fermat}. If $u=0$, the squared distance to a vertex
is $1/2$, which is not a rational square. The other midline follows
by exchanging coordinates. On $x=y$ the distance to $P_0$ is
$\sqrt2\,x$; on $x+y=1$ the distance to $P_1$ is
$\sqrt2\,(1-x)$. These are irrational because the coordinates are
rational and the indicated factors are positive. Finally, the
perpendicular bisector of every pair of square vertices is one of
these four lines.
\end{proof}

\begin{proposition}[Primitive congruences]\label{int:primitive}
For every strict rational-distance point,
\[
 24\mid L.
\]
After a single exchange of the coordinates, $X$ is odd and $4\mid Y$.
All four $A_i$ are odd and prime to $3$. Every prime divisor of
\begin{equation}\label{eq:int:diagonal-product}
 (X-Y)(X+Y-L)
\end{equation}
is congruent to $1$ or $7$ modulo $8$.
\end{proposition}

\begin{proof}
If $L$ were odd, one could choose an odd element from each of
$\{X,X-L\}$ and $\{Y,Y-L\}$. The corresponding equation in
\eqref{eq:int:integer-four} would have right side $2$ modulo $4$.
Thus $L$ is even. Primitivity excludes both $X,Y$ even, and their
squared sum excludes both odd. Exchange coordinates so that $X$ is
odd. All $A_i$ are odd; reduction modulo $8$ then gives $4\mid Y$
and $4\mid L-Y$, hence $4\mid L$.

If $3\nmid L$, each of $\{X,X-L\}$ and $\{Y,Y-L\}$ contains a
unit modulo $3$. The associated squared distance would be $2$ modulo
$3$. Hence $3\mid L$. Now exactly one of $X,Y$ is divisible by $3$:
both units would again give a nonsquare sum, while both nonunits would
contradict primitivity. Consequently every $A_i$ is a $3$-adic unit.

The integer in \eqref{eq:int:diagonal-product} is nonzero by
Lemma~\ref{lem:int:symmetry}, and both factors are odd. Let an odd
prime $p$ divide $X-Y$. If $X,Y$ are units modulo $p$, then $A_0^2$
reduces to twice a nonzero square. Otherwise both are zero modulo $p$;
primitivity gives $p\nmid L$, and $A_3^2$ reduces to $2L^2$.
In either case $\leg2p=1$. If instead $p\mid X+Y-L$, then
$A_1^2\equiv2Y^2$ and $A_2^2\equiv2X^2\pmod p$; at least one of
$X,Y$ is a unit, again by primitivity. Thus $\leg2p=1$ in this case
as well. The supplementary law for $2$ gives $p\equiv1,7\pmod8$.

A nonzero integer supported on these primes, with either sign, is
$1$ or $7$ modulo $8$. Were $L\equiv4\pmod8$, our orientation would
instead give
\[
 (X-Y)(X+Y-L)=X^2-Y^2-LX+LY\equiv5\pmod8.
\]
This contradiction proves $8\mid L$, and therefore $24\mid L$.
\end{proof}

\subsection{Common content of the integer radii}

The common divisor of the four radii has a precise relation to the
coordinate denominator.

\begin{proposition}\label{int:content}
Set
\[
 G=\gcd(A_0,A_1,A_2,A_3),\qquad a_i=A_i/G.
\]
Then $G^2\mid L$. Every prime divisor of $G$ is $1$ modulo $4$,
and $\gcd(G,6)=1$. Writing
\begin{equation}\label{eq:int:content-normalization}
 H=L/G^2,
\end{equation}
one has
\[
 24\mid H,\qquad r_i=\frac{a_i}{GH},\qquad
 \gcd(a_0,a_1,a_2,a_3)=1.
\]
For every prime $p\mid H$, all four $a_i$ are $p$-adic units.
\end{proposition}

\begin{proof}
The congruences already proved give $\gcd(G,6)=1$. Suppose an odd
prime $p$ divides $G$ but not $L$. Subtracting the origin square from
the two adjacent squares gives $X\equiv Y\equiv L/2\pmod p$.
Then $A_0^2\equiv L^2/2\ne0\pmod p$, a contradiction. Thus
$p\mid L$. Primitivity and $X^2+Y^2\equiv0\pmod p$ now force both
$X,Y$ to be units and $(Y/X)^2\equiv-1\pmod p$. In particular
$p\equiv1\pmod4$.

Write $s=v_p(G)$ and $e=v_p(L)$. The identity
\[
 A_1^2-A_0^2=L(L-2X)
\]
has right side of exact valuation $e$, since $p\nmid X$. Its left
side is divisible by $p^{2s}$; hence $e\ge2s$. This proves
$G^2\mid L$ at every prime. Division of the adjacent squared
differences by $G^2$ shows that all four $a_i^2$ have the same residue
modulo each prime divisor of $H$. If one were divisible by that prime,
all would be, contrary to their collective gcd. The remaining
assertions follow from \eqref{eq:int:content-normalization} and
$24\mid L$.
\end{proof}

There is also a uniform relation between the denominators of a sum
and a difference of two radii. Here the denominator of a nonzero
rational number always means its positive reduced denominator.

\begin{theorem}[Complementary denominators]\label{thm:int:pair-denominators}
For any distinct $i,j$, let $d_+$ and $d_-$ be the denominators of
$r_i+r_j$ and $r_i-r_j$. Then
\begin{equation}\label{eq:int:denominator-product}
 \gcd(d_+,d_-)=G,\qquad d_+d_-=L/2.
\end{equation}
\end{theorem}

\begin{proof}
Both denominators divide $L$. Fix an odd $p\mid L$, and write
$e=v_p(L)$, $s=v_p(G)$. If $s=0$, the four $A_i$ are units modulo
$p$, because their squares have the same residue and at least one is
a unit. Every difference $A_i^2-A_j^2$ is divisible by $p^e$.
Exactly one of $A_i+A_j,A_i-A_j$ is divisible by $p$, so their
reduced denominator depths are $0$ and $e$.

Suppose $s>0$. The proof of Proposition~\ref{int:content} shows
$Y/X\equiv\pm\sqrt{-1}\pmod p$. Thus
$L-2X$, $L-2Y$, $X-Y$, and $X+Y-L$ are all units modulo $p$.
The six squared differences are, up to sign,
\begin{equation}\label{eq:int:six-differences}
 L(L-2X),\quad L(L-2Y),\quad
 2L(X+Y-L),\quad 2L(X-Y),
\end{equation}
with repetitions in the parallel directions. Each has exact depth
$e$. If $e>2s$, the four $a_i$ are units, and their sum and difference
have depths $0$ and $e-2s$. The depths of the corresponding
denominators are consequently $e-s$ and $s$. If $e=2s$, the squared
difference of the normalized pair is a unit, so both its sum and
difference are units; both denominator depths are $s$. In all cases
the sum of the two depths is $e$ and their minimum is $s$.

At $2$, write $e=v_2(L)\ge3$ and keep $X$ odd, $4\mid Y$. All four
$A_i$ are odd. The horizontal and opposite squared differences in
\eqref{eq:int:six-differences} have depth $e+1$; the vertical ones
have depth at least $e+3$. Of a sum and a difference of two odd
integers, precisely one has depth one. The other therefore has depth
at least $e$. Their reduced denominator depths are $e-1$ and $0$.
Combining the prime-by-prime statements proves
\eqref{eq:int:denominator-product}.
\end{proof}

\subsection{Gaussian roots and the four corner classes}

Put $c_0=0$, $c_1=1$, $c_2=i$, $c_3=1+i$, and define the integer
Gaussian displacements
\[
 U_j=X+iY-Lc_j.
\]
For $j=(\epsilon,\eta)$ put
\[
 d_j=\gcd(|X-\epsilon L|,|Y-\eta L|),\qquad
 m_j=\sf(d_j),
\]
where $\sf$ denotes the positive squarefree part of a positive integer.
The coordinate orientation is the one in Proposition~\ref{int:primitive}.

\begin{theorem}[Canonical corner classes]\label{int:gaussian}
The integers $d_j$ are odd, pairwise coprime, and coprime to $L$.
Every prime dividing any $d_j$ is $1$ or $7$ modulo $8$. There are
nonzero $q_j\in\Z[i]$ such that
\begin{equation}\label{eq:int:gaussian-roots}
 U_j=m_jq_j^2,\qquad \Nm(q_j)=A_j/m_j>0.
\end{equation}
In particular $m_0,m_1,m_2,m_3$ are positive odd squarefree integers
with disjoint supports, and, for $z=x+iy$,
\begin{equation}\label{eq:int:gaussian-classes}
 [z-c_j]=[Lm_j]\quad\text{in }\Q(i)^*/\Q(i)^{*2}.
\end{equation}
Conversely, if a rational strict point has a primitive representation
with $X$ odd and admits \eqref{eq:int:gaussian-roots} with positive
integers $m_j$ and $q_j\in\Z[i]$, all four of its distances are rational.
\end{theorem}

\begin{proof}
Every first coordinate $X-\epsilon L$ is odd, so all $d_j$ are odd.
A prime dividing $d_j$ and $L$ would divide $X,Y,L$. A prime dividing
two different $d_j$ would, by subtracting their coordinate pairs,
divide $L$ as well. This proves the asserted coprimalities.
If $p\mid d_j$, the squared distance from the opposite vertex is
$2L^2$ modulo $p$. Since $p\nmid L$, this implies $\leg2p=1$.

Divide the two legs and hypotenuse of the $j$th Pythagorean triangle
by $d_j$. This gives coprime integers $a,b$ with $a$ odd, $b$ even,
and $a^2+b^2$ an integer square. The Gaussian integers $a+ib,a-ib$
are coprime: a common odd Gaussian prime would divide both $a,b$,
while $1+i$ does not divide their odd norm. Unique factorization in
$\Z[i]$ therefore makes $a+ib$ a unit times a square. Here unique
factorization follows from the norm-Euclidean algorithm: rounding both
coordinates of a complex quotient gives a remainder whose norm is
at most half that of the divisor. The unit is
not $\pm i$, because a Gaussian square of odd norm has odd real
part and even imaginary part. The remaining units $\pm1$ are
Gaussian squares. Writing $d_j=m_jt_j^2$ and absorbing $t_j$ gives
$U_j=m_jq_j^2$. Taking the positive square root of its norm proves
$A_j=m_j\Nm(q_j)$. Finally $U_j/L$ has squareclass $m_j/L$, which
equals $Lm_j$. For the converse, taking norms in the four displayed
root equations gives the four distances
$m_j\Nm(q_j)/L\in\Q_{>0}$.
\end{proof}

The product
\begin{equation}\label{eq:int:kernel}
 n=m_0m_1m_2m_3
\end{equation}
is thus a positive odd squarefree integer, prime to $L$, supported on
primes $1$ and $7$ modulo $8$. Equation
\eqref{eq:int:gaussian-classes} records both its corner allocation
and the single common factor $L$.

\begin{proposition}[The common squarefree kernel]\label{prop:int:kernel-product}
Put $t=r_0r_3+r_1r_2$. Then $t>1$ and
\[
 [t^2-1]=[n]\quad\text{in }\Q^*/\Q^{*2}.
\]
Consequently the Gaussian corner kernel \eqref{eq:int:kernel} equals
$\sf(t^2-1)$.
\end{proposition}
\begin{proof}
Write $z=x+iy$, and set
\[
 A=z(z-1-i),\qquad B=(z-1)(z-i),\qquad
 a=r_0r_3,\qquad b=r_1r_2.
\]
Then $B-A=i$, $A\overline A=a^2$ and $B\overline B=b^2$.
The triangle inequality gives $a+b\geq1$. Equality forces $A$ and $B$
to lie on the imaginary axis, and hence
\[
 \operatorname{Re}(A)=(x-y)(x+y-1)=0.
\]
Lemma~\ref{lem:int:symmetry} excludes equality, so $t>1$.
Expanding with the two norm identities gives
\[
 (bA+aB)^2
 =AB\bigl(A\overline B+B\overline A+2ab\bigr)
 =AB\bigl((a+b)^2-|A-B|^2\bigr)
 =(t^2-1)AB.
\]
Multiplication of the four Gaussian identities
\eqref{eq:int:gaussian-roots} gives
$AB=n(q_0q_1q_2q_3/L^2)^2$. Thus $(t^2-1)/n$ is a square
in $\Q(i)$. A positive rational number which is a square in
$\Q(i)$ is a square in $\Q$: if $(u+iv)^2$ is positive rational,
then $uv=0$ and positivity forces $v=0$. This proves the claim.
\end{proof}

The common factor also occurs in the unit restrictions at ramified
split primes. For $p\equiv1\pmod8$, choose either square root
$\iota$ of $-1$ in $\F_p$ and define
\[
 \epsilon(p)=\leg{1+\iota}{p}.
\]
This is independent of the choice: replacing $\iota$ by $-\iota$
multiplies $1+\iota$ by $-\iota$, a square in $\F_p^*$.
Equivalently,
\begin{equation}\label{eq:int:quartic-sign}
 \epsilon(p)=2^{(p-1)/4}(-1)^{(p-1)/8}\in\{1,-1\},
\end{equation}
where the power is evaluated in $\F_p$ and the result identified
with its sign.

\begin{proposition}\label{prop:int:ramified-units}
If $p\equiv1\pmod8$ divides $m_j$, then
\[
 \leg{Lm_k}{p}=1
\]
for either adjacent corner $k$, and
\[
 \leg{Lm_{j^{\mathrm{op}}}}p=\epsilon(p)
\]
for the opposite corner $j^{\mathrm{op}}$.
\end{proposition}

\begin{proof}
Use the reduction homomorphism $\Z[i]\longrightarrow\F_p$ given by
$i\mapsto\iota$. From \eqref{eq:int:gaussian-roots},
$X+iY\equiv Lc_j\pmod p$, and therefore for $k\ne j$
\[
 m_kq_k^2\equiv L(c_j-c_k)\pmod p.
\]
All factors on the right are nonzero. Adjacent differences are
$\pm1$ or $\pm i$, all squares modulo $p$. Opposite differences are
these units times $1+i$. Taking quadratic characters proves the
assertions. For \eqref{eq:int:quartic-sign}, use
$(1+\iota)^2=2\iota$ and raise to $(p-1)/4$.
\end{proof}

\subsection{Signed radii and denominator growth}

The integral equations also connect divisibility to the real size of
the square. Set
\[
 d=2+\sqrt2.
\]
For every strict point,
\begin{equation}\label{eq:int:total-bound}
 A_0+A_1+A_2+A_3<dL.
\end{equation}
Indeed, the sum of the distances to the four vertices is a convex
function on the square and equals $d$ at each vertex. A strict point
is a convex combination of all four vertices with positive weights.
The triangle inequality for the distance to $P_0$ is strict in this
combination, because the two adjacent nonzero vertex vectors are not
parallel. This proves \eqref{eq:int:total-bound}.

Choose arbitrary signs and write $a,b,c,e$ for the signed versions
of $A_0,A_1,A_2,A_3$. Put
\[
 S=a+b+c+e,\quad T=a-b-c+e,\quad
 P=ae-bc,\quad Q=ae+bc.
\]
The four-distance equations give the identities
\begin{equation}\label{eq:int:signed-identities}
 a^2+e^2=b^2+c^2,\qquad
 PQ=-L^2(L-2X)(L-2Y),\qquad ST=2P.
\end{equation}
The second follows by expanding $A_0^2A_3^2-A_1^2A_2^2$;
the last follows from the first. Lemma~\ref{lem:int:symmetry} shows
that the product on the right is nonzero. Thus $S,T,P,Q$ are all
nonzero, for every sign choice. Moreover
\begin{equation}\label{eq:int:signed-bound}
 |T|\le A_0+A_1+A_2+A_3<dL.
\end{equation}

\begin{theorem}[Prime-power comparisons]\label{thm:int:prime-comparisons}
Let $K=2^{v_2(L)}$. Then
\begin{equation}\label{eq:int:two-bound}
 K^2<\frac{2+\sqrt2}{4}\,L.
\end{equation}
If $p$ is an odd prime with $p\mid L$ and $p\mid XY$, and
$P_p=p^{v_p(L)}$, then
\begin{equation}\label{eq:int:odd-bound}
 P_p^2<\frac{2+\sqrt2}{2p}\,L.
\end{equation}
In particular
\begin{equation}\label{eq:int:three-bound}
 3^{2v_3(L)}<\frac{2+\sqrt2}{6}\,L.
\end{equation}
\end{theorem}

\begin{proof}
For the dyadic assertion use $X$ odd and $4\mid Y$. Write
$h=v_2(L)\ge3$ and $j=v_2(L-2Y)\ge3$; the latter is finite by
Lemma~\ref{lem:int:symmetry}. Choose the signs of $a,b,c,e$ so that
all four are $1$ modulo $4$. Every pair sum then has depth one.
The horizontal and vertical squared differences give
\[
 v_2(a-b)=v_2(c-e)=h,\qquad
 v_2(a-c)=v_2(b-e)=h+j-1.
\]
Consequently $a,b,c,e$ are congruent modulo $2^h$. It follows that
\[
 S\equiv4a\pmod{2^h},\qquad
 Q\equiv2a^2\pmod{2^h},
\]
so $v_2(S)=2$ and $v_2(Q)=1$. Since $v_2(L-2X)=1$, the identities
\eqref{eq:int:signed-identities} yield
\[
 v_2(P)=2h+j,\qquad v_2(T)=2h+j-1\ge2h+2.
\]
The nonzero integer $T$ therefore has $|T|\ge4K^2$.
Equation~\eqref{eq:int:signed-bound} proves \eqref{eq:int:two-bound}.

For the odd-prime assertion, exchange the coordinates independently
if needed so that $p\nmid X$ and $p\mid Y$. This is possible by
primitivity and the hypothesis $p\mid XY$. The four $A_i$ are
$p$-adic units and their squares have the same residue. Choose their
signs so that $a,b,c,e$ have the same nonzero residue modulo $p$.
Then $S\equiv4a$ and $Q\equiv2a^2\pmod p$ are units.
Set $h=v_p(L)$ and $j=v_p(L-2Y)\ge1$, again finite. The same
identities now give the exact depth
\[
 v_p(T)=2h+j\ge2h+1.
\]
Every $A_i$ is odd, so the signed sum $T$ is even. Hence
\[
 2pP_p^2\le2p^{2h+j}\le |T|<dL,
\]
which proves \eqref{eq:int:odd-bound}. At $p=3$, the required
divisibility $3\mid XY$ was proved in
Proposition~\ref{int:primitive}, so \eqref{eq:int:three-bound}
follows.
\end{proof}

Writing $L=Km$ with $m$ odd, \eqref{eq:int:two-bound} takes the
linear form
\[
 K<\frac{2+\sqrt2}{4}\,m.
\]
For odd $p$ as in the theorem, writing $L=P_ps$ similarly gives
$P_p<(2+\sqrt2)s/(2p)$. These comparisons apply to the complete
prime powers in the denominator.

\begin{corollary}\label{cor:int:six-smooth}
The least coordinate denominator of a strict rational-distance point
has a prime divisor at least $5$.
\end{corollary}

\begin{proof}
If $L$ were supported on $2$ and $3$, write $L=KP$ with
$K=2^{v_2(L)}$ and $P=3^{v_3(L)}$. The preceding theorem gives
\[
 K^2<dL/4,\qquad P^2<dL/6.
\]
Multiplication and $KP=L$ yield $1<d^2/24$. But
$d^2=6+4\sqrt2<24$, a contradiction.
\end{proof}

The same signed identities give a precise statement in the other
$5$-adic direction, where all four radii have nontrivial content.

\begin{proposition}[The unit-coordinate direction at $5$]
\label{prop:int:five-depth}
Suppose $5\mid L$ and $5\nmid XY$. Write $L=5^h m$ with $5\nmid m$.
There is an integer $s\ge1$ such that
\[
 v_5(A_i)=s\quad(0\le i\le3),\qquad h>2s.
\]
Some signed sum $T$ of the four $A_i$ has exact depth $2h-3s$,
and
\begin{equation}\label{eq:int:five-bound}
 5^{h-3s}<\frac{2+\sqrt2}{2}\,m.
\end{equation}
\end{proposition}

\begin{proof}
Nonzero squares modulo $5$ are $1$ and $4$. Since $X,Y$ are units
and $X^2+Y^2$ is a square, this sum must be zero modulo $5$.
Thus $Y/X\equiv\pm2\pmod5$, and all four $A_i$ are divisible by $5$.
Let $s=\min_i v_5(A_i)$. The adjacent squared differences have
exact depth $h$, so $h\ge2s$.

If $h=2s$, division of the four squared-distance equations by
$5^{2s}$ gives residues
\[
 c,\quad c+u,\quad c+v,\quad c+u+v,
 \qquad u,v\ne0,\quad v/u=\pm2
\]
in $\F_5$. These four residues are distinct, whereas $\F_5$ has
only three square residues. Therefore $h>2s$. Every squared
difference has depth at least $h$, and hence all four radii have
depth exactly $s$.

Divide the four radii by $5^s$ and choose signs so that the resulting
units $a,b,c,e$ have the same residue modulo $5$. Their sums
$S=a+b+c+e$ and $Q=ae+bc$ are units. The determinant identity becomes
\[
 (ae-bc)(ae+bc)=-\frac{L^2(L-2X)(L-2Y)}{5^{4s}}.
\]
Both $L-2X,L-2Y$ are units, so $ae-bc$ and
$a-b-c+e$ have exact depth $2h-4s$. Multiplying the latter by
$5^s$ gives an actual signed sum of the integer radii with exact
depth $2h-3s$. It is nonzero and even. Its absolute value is at
least $2\cdot5^{2h-3s}$ and less than $dL$, which proves
\eqref{eq:int:five-bound}.
\end{proof}

\section{Reciprocity for the four corner colours}\label{sec:profiles}

The squarefree parts of the four corner contents carry a finer reciprocity
constraint than their product.  We shall prove that the product belonging
to either pair of opposite corners is congruent to $1$ or $15$ modulo
$16$.  The proof also describes the prime divisors of all six sums of two
integer distances.  Its ingredients are the Gaussian square roots of the
individual displacements and ordinary quadratic reciprocity.

\subsection{Integer distances and Gaussian roots}

Fix a strictly interior point with four positive rational distances, and use the corner
order
\[
 c_0=0,\qquad c_1=1,\qquad c_2=\mathrm i,\qquad c_3=1+\mathrm i.
\]
Write
\begin{equation}\label{eq:profile-lattice}
 x=X/L,\qquad y=Y/L,\qquad L>0,\qquad \gcd(X,Y,L)=1,
 \qquad A_j=Lr_j\in\Z_{>0}.
\end{equation}
Interchanging $x$ and $y$ if necessary, take $X$ odd.  The integral
normalization of Section~\ref{sec:integral} gives $4\mid Y$ and
$8\mid L$.  This interchange only exchanges the labels $1$ and $2$,
so it preserves the two opposite pairs as unordered pairs.

Put
\begin{equation}\label{eq:profile-colours}
 U_j=X+\mathrm iY-Lc_j,\qquad
 g_j=\gcd\bigl(|\operatorname{Re}U_j|,|\operatorname{Im}U_j|\bigr),
 \qquad m_j=\sf(g_j).
\end{equation}
The $m_j$ are positive, odd, squarefree, pairwise coprime, and coprime
to $L$.  Every prime dividing them is $1$ or $7$ modulo $8$.
These are canonical positive integers attached to the primitive
coordinate lattice.

We shall use Gaussian integers $w_j$ satisfying
\begin{equation}\label{eq:profile-integer-roots}
 U_j=m_jw_j^2,\qquad
 N_{\Q(\mathrm i)/\Q}(w_j)=\frac{A_j}{m_j}>0.
\end{equation}
For completeness, the elementary construction is as follows.  The
triple
\[
 \left(\frac{\operatorname{Re}U_j}{g_j},
       \frac{\operatorname{Im}U_j}{g_j},\frac{A_j}{g_j}\right)
\]
is a primitive Pythagorean triple whose first leg is odd and whose
second leg is even.  If its legs and hypotenuse are denoted by $a,b,k$,
the positive integers $(k+a)/2$ and $(k-a)/2$ are relatively prime,
and their product is $(b/2)^2$.  Thus they are squares, say $s^2,t^2$,
and a choice of the sign of $t$ gives $a+\mathrm ib=(s+\mathrm it)^2$.
Multiplying this root by the positive integer $\sqrt{g_j/m_j}$ gives
$w_j$.  Its norm is positive and its norm squared is $A_j^2/m_j^2$,
which proves the second identity in \eqref{eq:profile-integer-roots}.

In particular, $A_j/m_j$ is an odd sum of two integer squares, whence
\begin{equation}\label{eq:profile-radius-residue}
 A_j\equiv m_j\pmod4.
\end{equation}
The signs in the norm identity will determine the characters at primes
congruent to $7$ modulo $8$.

The diagonal pairs have squared separation $2$, and adjacent pairs have
squared separation $1$.  For $i\ne j$ define
\begin{equation}\label{eq:profile-pair-data}
 s_{ij}=m_im_j,\qquad d_{ij}=|c_i-c_j|^2\in\{1,2\},\qquad
 V_{ij}=\frac{A_i+A_j}{2}\in\Z_{>0}.
\end{equation}

\begin{lemma}[The six pair identities]\label{lem:profile-pair-identity}
For every pair $i\ne j$,
\begin{equation}\label{eq:profile-pair-identity}
 (A_i+A_j)^2-
 4s_{ij}\operatorname{Re}(w_i\overline{w_j})^2=d_{ij}L^2.
\end{equation}
Consequently, with $\Sigma_{ij}=r_i+r_j$ and
$T_{ij}=2\operatorname{Re}(w_i\overline{w_j})/L$, one has
\[
 \Sigma_{ij}^2-s_{ij}T_{ij}^2=d_{ij}.
\]
\end{lemma}

\begin{proof}
Set $Z=w_i\overline{w_j}$ and $s=s_{ij}$.  Equation
\eqref{eq:profile-integer-roots} gives
\[
 N(Z)=\frac{A_iA_j}{s},\qquad
 Z^2=\frac{U_i\overline{U_j}}{s}.
\]
The identity $2\operatorname{Re}(Z)^2=N(Z)+\operatorname{Re}(Z^2)$
therefore shows that the left side of
\eqref{eq:profile-pair-identity} is
\[
 A_i^2+A_j^2-2\operatorname{Re}(U_i\overline{U_j})
 =|U_i-U_j|^2=d_{ij}L^2.
\]
\end{proof}

\subsection{Prime divisors of all six half-sums}

\begin{theorem}[Six-pair support]\label{thm:profile-six-pair-support}
Let $p$ be an odd prime dividing $V_{ij}$.  Then
\begin{equation}\label{eq:profile-six-pair-support}
 p\nmid s_{ij},\qquad \leg{-s_{ij}}p=1.
\end{equation}
Equivalently, every odd prime divisor of $V_{ij}$ is unramified and
split in $\Q(\sqrt{-s_{ij}})$.
\end{theorem}

\begin{proof}
Write $s=s_{ij}$, $d=d_{ij}$, and $V=V_{ij}$.
If $p\mid m_i$, the point reduces to corner $i$ modulo $p$:
$U_i\equiv0$ and $A_i\equiv0$.  Since $p\nmid L$,
\[
 A_j^2\equiv dL^2\not\equiv0\pmod p.
\]
Thus $p\nmid V$.  The same argument applies to a prime dividing $m_j$,
and proves the first assertion.

Suppose next that $p\mid L$.  If $p\equiv1\pmod4$, choose a square
root of $-1$ modulo $p$ for which $X+\mathrm iY$ is nonzero.  Such a
choice exists: vanishing for both choices would force $p\mid X,Y,L$.
In this one embedding all the $U_k$ have the same nonzero residue,
and hence all $w_k$ have nonzero residues.  Reducing
\[
 \left(\frac{w_i}{w_j}\right)^2
   =\frac{m_j}{m_i}\frac{U_i}{U_j}
\]
gives $\leg{s}p=1$.  Since $\leg{-1}p=1$, the required conclusion
follows.

If $p\equiv3\pmod4$, the form $X^2+Y^2$ is anisotropic modulo $p$.
Primitivity therefore makes its value nonzero, and every $A_k$ and
$w_k$ is a unit at $p$.  Work in $\F_{p^2}=\F_p(\mathrm i)$, where
conjugation is Frobenius.  For $h=w_i/w_j$ and $a=m_j/m_i\in\F_p^*$,
we have $h^2=a$, and consequently
\[
 N_{\F_{p^2}/\F_p}(h)=h^{p+1}=a^{(p+1)/2}
     =a\leg ap.
\]
The actual norm identity in \eqref{eq:profile-integer-roots} gives at
the same time
\[
 N(h)=\frac{m_j}{m_i}\frac{A_i}{A_j}.
\]
It follows that $A_i/A_j=\leg{s}p$ modulo $p$.  The assumption
$p\mid V$ makes this ratio $-1$, so $\leg{s}p=-1$ and
$\leg{-s}p=1$.

It remains to treat $p\nmid L$.  Reducing
\eqref{eq:profile-pair-identity}, and using $p\mid V$ and $p\nmid s$,
gives
\[
 -4s\operatorname{Re}(w_i\overline{w_j})^2=dL^2.
\]
All factors on the right are nonzero modulo $p$.  For an adjacent
pair, $d=1$, this proves $\leg{-s}p=1$.  For an opposite pair the same
argument proves the assertion whenever $\leg2p=1$.

Finally, an opposite sum cannot vanish modulo a prime with
$\leg2p=-1$.  For the pair $(0,3)$, the equality $A_0=-A_3$ modulo $p$
implies $x+y=1$ in $\F_p$.  The other two distance squares then become
$2x^2$ and $2y^2$.  At least one is nonzero, which is impossible when
$2$ is a nonsquare.  For the pair $(1,2)$, the same reasoning gives
$x=y$ and the other two squares $2x^2$ and $2(x-1)^2$, again with at
least one nonzero.  This exhausts the cases.
\end{proof}

\subsection{The characters at a corner prime}

Define a character on the subgroup $\{1,7,9,15\}$ of
$(\Z/16\Z)^*$ by
\begin{equation}\label{eq:profile-sigma}
 \sigma(1)=\sigma(15)=1,\qquad
 \sigma(7)=\sigma(9)=-1.
\end{equation}
We use the same notation for positive integers supported on primes
congruent to $1$ or $7$ modulo $8$.  Multiplication of residues shows
that this is multiplicative.  For a prime in the support,
\begin{equation}\label{eq:profile-prime-sigma}
 \sigma(p)=
 \begin{cases}
 (-1)^{(p-1)/8},&p\equiv1\pmod8,\\
 (-1)^{(p+1)/8},&p\equiv7\pmod8.
 \end{cases}
\end{equation}

\begin{lemma}[Ramified character]\label{lem:profile-ramified-character}
For $p\mid m_i$ and $j\ne i$,
\begin{equation}\label{eq:profile-ramified-character}
 \leg{V_{ij}}p=\leg{m_j}p\,\sigma(p)^{d_{ij}-1}.
\end{equation}
\end{lemma}

\begin{proof}
At the corner prime $p$ we have $A_i\equiv0$ and
\[
 U_j\equiv L(c_i-c_j)\pmod p.
\]
Since $\leg2p=1$, the characters of $V_{ij}$ and $A_j$ agree.
The difference $c_i-c_j$ is a Gaussian unit for adjacent corners,
and a Gaussian unit times $1+\mathrm i$ for opposite corners.

First suppose $p\equiv7\pmod8$.  In $\F_{p^2}$, put
$u=c_i-c_j$.  The two identities for $w_j$ give
\[
 \frac{A_j}{m_j}
  =N(w_j)
  =\left(\frac{Lu}{m_j}\right)^{(p+1)/2},
\]
and hence
\begin{equation}\label{eq:profile-frobenius-radius}
 A_j=L\leg{Lm_j}p\,u^{(p+1)/2}.
\end{equation}
The exponent $(p+1)/2$ is divisible by $4$.  It therefore kills every
Gaussian unit.  In the adjacent case its final factor is $1$.
In the opposite case that factor is
\[
 z_p=(1+\mathrm i)^{(p+1)/2}
     =(-1)^{(p+1)/8}\,2^{(p+1)/4}\in\F_p,
 \qquad z_p^2=2.
\]
Because $\leg{-1}p=-1$ and $\leg2p=1$, its character is
$\leg{z_p}p=\sigma(p)$.  Taking characters in
\eqref{eq:profile-frobenius-radius}, a sign $\pm1$ has its own value
as character, so
\[
 \leg{A_j}p
 =\leg Lp\leg{Lm_j}p\,\sigma(p)^{d_{ij}-1}
 =\leg{m_j}p\,\sigma(p)^{d_{ij}-1}.
\]

Now suppose $p\equiv1\pmod8$ and work in a Gaussian embedding into
$\F_p$.  All four Gaussian units are squares.  For an adjacent pair,
$w_j^2=Lu/m_j$ therefore gives $\leg{Lm_j}p=1$, while
$A_j^2=L^2$ gives $\leg{A_j}p=\leg Lp=\leg{m_j}p$.

For an opposite pair write
$\lambda_p=\leg{1+\mathrm i}p$.  This value is unchanged on replacing
$\mathrm i$ by $-\mathrm i$, because their ratio is a Gaussian unit.
The root equation gives $\leg{Lm_j}p=\lambda_p$.  Choose
$J\in\F_p$ with $J^2=\mathrm i$.  It has order $8$, so
\[
 \leg Jp=(-1)^{(p-1)/8}=\sigma(p).
\]
Since $((1+\mathrm i)/J)^2=2$, either square root of $2$ has character
$\lambda_p\sigma(p)$; changing its sign does not change that
character.  The equation $A_j^2=2L^2$ now yields
\[
 \leg{A_j}p
  =\leg Lp\lambda_p\sigma(p)
  =\leg{m_j}p\sigma(p).
\]
This proves the formula for both congruence classes and both separations.
\end{proof}

\subsection{A pair reciprocity identity}

For a positive integer $a$, let
\begin{equation}\label{eq:profile-G}
 \mathcal G(a)=\prod_{p\equiv3\ (4)}(-1)^{v_p(a)}.
\end{equation}
Thus $\mathcal G(a)$ is $1$ or $-1$ according as the odd part of $a$
is $1$ or $3$ modulo $4$.  Let $e_j\in\{0,1\}$ be defined by
\begin{equation}\label{eq:profile-e}
 m_j\equiv(-1)^{e_j}\pmod4.
\end{equation}
Equivalently, $e_j$ is the parity of the number of prime factors of
$m_j$ which are $7$ modulo $8$.

\begin{theorem}[Pair reciprocity]\label{thm:profile-master}
For every pair $i\ne j$, let $e\in\{0,1\}$ satisfy
$e\equiv e_i+e_j\pmod2$.  Then
\begin{equation}\label{eq:profile-master}
 \mathcal G(V_{ij})^{1+e}
   =(-1)^{e_ie_j}\sigma(m_im_j)^{d_{ij}-1}.
\end{equation}
\end{theorem}

\begin{proof}
Put $s=m_im_j$ and $V=V_{ij}$.  All Jacobi symbols below have coprime
positive arguments, by Theorem~\ref{thm:profile-six-pair-support}.
Multiplication of \eqref{eq:profile-ramified-character} over the primes
of $m_i$ and $m_j$ gives
\[
 \leg Vs
  =\leg{m_j}{m_i}\leg{m_i}{m_j}\sigma(s)^{d_{ij}-1}
  =(-1)^{e_ie_j}\sigma(s)^{d_{ij}-1}.
\]
Here the second equality is quadratic reciprocity; a Jacobi symbol with
denominator $1$ is interpreted as $1$.

Let $V_{\mathrm{odd}}$ be the odd part of $V$.  The six-pair support
theorem gives, for every odd $p\mid V$,
\[
 \leg sp=\leg{-1}p.
\]
Multiplication with valuations therefore gives
\[
 \leg{s}{V_{\mathrm{odd}}}=\mathcal G(V).
\]
Also $\leg2s=1$, since every prime factor of $s$ is $1$ or $7$
modulo $8$.  A second application of quadratic reciprocity yields
\[
 \leg Vs
 =\leg{V_{\mathrm{odd}}}s
 =\leg{s}{V_{\mathrm{odd}}}\mathcal G(V)^e
 =\mathcal G(V)^{1+e}.
\]
Equating the two expressions proves the result.
\end{proof}

\subsection{The two opposite products}

\begin{theorem}[Opposite-colour congruences]\label{thm:profile-main}
For every strictly interior point with four positive rational distances,
the canonical corner colours in \eqref{eq:profile-colours} satisfy
\begin{equation}\label{eq:profile-main}
 \boxed{\quad m_0m_3\equiv1\text{ or }15\pmod{16},\qquad
               m_1m_2\equiv1\text{ or }15\pmod{16}.\quad}
\end{equation}
More precisely, an opposite product with an even number of prime factors
congruent to $7$ modulo $8$ is $1$ modulo $16$; one with an odd number
is $15$ modulo $16$.
\end{theorem}

\begin{proof}
Fix an opposite pair $i,j$ and put $s=m_im_j$ and $V=V_{ij}$.
Here $d_{ij}=2$.  If $e_i\ne e_j$, the left side of
\eqref{eq:profile-master} is $\mathcal G(V)^2=1$, and
$e_ie_j=0$.  Hence $\sigma(s)=1$.  Since $s\equiv7\pmod8$ in
this case, the conclusion is $s\equiv15\pmod{16}$.

Suppose instead that $e_i=e_j$.  By
\eqref{eq:profile-radius-residue}, the integers $A_i,A_j$ have the
same residue $(-1)^{e_i}$ modulo $4$.  Thus $V$ is odd and their
half-difference $D=(A_i-A_j)/2$ is even.  The squared-distance
identities give the exact formulas
\[
 A_0^2-A_3^2=2L(X+Y-L),\qquad
 A_2^2-A_1^2=2L(X-Y).
\]
The factors $X+Y-L$ and $X-Y$ are odd.  Consequently, in either
opposite pair,
\[
 v_2(VD)=v_2(L)-1\ge2.
\]
Since $V$ is odd, $4\mid D$.  It follows that
\[
 V\equiv A_i\equiv(-1)^{e_i}\pmod4,
 \qquad \mathcal G(V)=(-1)^{e_i}.
\]
Now $e=0$ and $e_ie_j=e_i$, so the master identity reads
\[
 (-1)^{e_i}=(-1)^{e_i}\sigma(s).
\]
Therefore $\sigma(s)=1$.  As $s\equiv1\pmod8$, this gives
$s\equiv1\pmod{16}$.
\end{proof}

Write $b=m_0m_3$, $c=m_1m_2$ and $n=bc$.  The allowable residue
patterns take the particularly simple form
\begin{equation}\label{eq:profile-residue-table}
\begin{array}{c|c}
 n\pmod{16}&(b,c)\pmod{16}\\ \hline
 1&(1,1),\ (15,15)\\
 15&(1,15),\ (15,1).
\end{array}
\end{equation}

\begin{corollary}\label{cor:profile-split}
If every prime dividing $m_0m_1m_2m_3$ is $1$ modulo $8$, then
\[
 m_0m_3\equiv m_1m_2\equiv1\pmod{16}.
\]
\end{corollary}

\begin{proof}
Every $e_j$ is zero, so both opposite products fall in the even case
of Theorem~\ref{thm:profile-main}.
\end{proof}

For adjacent pairs the exponent $d_{ij}-1$ is zero, and the same
reciprocity theorem gives
\[
 \mathcal G(V_{ij})^{1+e}=(-1)^{e_ie_j}.
\]
In particular, if two adjacent colours have the same residue modulo
$4$, the odd part of their half-sum has that same residue modulo $4$.
Together with Theorem~\ref{thm:profile-six-pair-support}, this supplies
simultaneous arithmetic information for every pair of corners.

The support conclusion also has a primitive-radius formulation.  Let
$G=\gcd(A_0,A_1,A_2,A_3)$ and write $M_j=A_j/G$.  Since all $A_j$ are
odd, $G$ and the $M_j$ are odd, and
\[
 \frac{A_i+A_j}{2}=G\frac{M_i+M_j}{2}.
\]
Every odd prime dividing $(M_i+M_j)/2$ consequently satisfies
\eqref{eq:profile-six-pair-support}.  This formulation uses the same
canonical colours \eqref{eq:profile-colours} and removes the common
content of the four integer radii.

\section{Critical radius content and an elliptic residue curve}\label{sec:critical-content}

The common content of the four integral radii produces a finite-field
curve independent of the original denominator. Its elliptic structure
gives a classification at every content prime. A second elliptic curve
records the square class of the opposite-product parameter.

\subsection{Critical reduction}

Write the primitive integral equations as
\begin{equation}\label{crit:integral}
\begin{aligned}
 A_0^2&=X^2+Y^2,& A_1^2&=(X-L)^2+Y^2,\\
 A_2^2&=X^2+(Y-L)^2,& A_3^2&=(X-L)^2+(Y-L)^2,
\end{aligned}
\end{equation}
with corner order \(00,10,01,11\), positive radii, \(0<X,Y<L\),
and \(\gcd(X,Y,L)=1\). Thus \(r_j=A_j/L\). Put
\[
 G=\gcd_j A_j,\qquad H=L/G^2,\qquad a_j=A_j/G.
\]
Recall from Proposition~\ref{int:content} that \(G^2\mid L\) and
\(p\mid G\Rightarrow p\equiv1\pmod4\). The local argument also
fixes the valuations needed here. If an odd \(p\mid G\) did not
divide \(L\), subtraction would give \(X=Y=L/2\pmod p\), contrary
to \(X^2+Y^2=0\). Hence \(p\mid L\), and primitivity makes
\(X,Y\) units. Consequently \((Y/X)^2=-1\pmod p\) and
\[
 v_p(A_1^2-A_0^2)=v_p\bigl(L(L-2X)\bigr)=v_p(L)\ge2v_p(G).
\]
The prime 2 cannot divide \(G\), since the equations modulo four
would then force \(X,Y,L\) all even.

\begin{definition}
A prime \(p\mid G\) is \emph{critical} if \(v_p(L)=2v_p(G)\).
\end{definition}

At a critical prime set \(s=v_p(G)>0\) and, in \(\F_p\),
\[
 \iota=Y/X,\qquad\iota^2=-1,\qquad b=-2HX\ne0,\qquad c=a_0^2.
\]
The four normalized radii satisfy
\begin{equation}\label{crit:pattern}
 (a_0^2,a_1^2,a_2^2,a_3^2)
 =\bigl(c,c+b,c+\iota b,c+(1+\iota)b\bigr).
\end{equation}
The shifts are distinct, so at most one normalized radius vanishes.
At a noncritical content prime, \(p\mid H\); all \(a_j^2\) are then
equal modulo \(p\), and \(\gcd_j a_j=1\) makes every \(a_j\) a unit.

\subsection{The fixed elliptic curve}

Over a field \(k\) of characteristic different from two containing
\(\iota^2=-1\), define
\begin{equation}\label{crit:curve}
 \mathcal C_\iota:\quad
 \begin{cases}
 V_2^2-V_0^2=\iota(V_1^2-V_0^2),\\
 V_3^2-V_0^2=(1+\iota)(V_1^2-V_0^2)
 \end{cases}
 \quad\subset\PP^3_k.
\end{equation}
Write \(b=V_1^2-V_0^2\), and let \(U_\iota\) be the open subset
\(bV_0V_1V_2V_3\ne0\).

\begin{proposition}\label{prop:crit-elliptic}
The curve \(\mathcal C_\iota\) is smooth, geometrically connected,
and \(k\)-isomorphic to the smooth projective curve
\(E:y^2=x^3-x\). Its eight points with \(b=0\) are
\([1:\epsilon_1:\epsilon_2:\epsilon_3]\), with \(\epsilon_j=\pm1\).
\end{proposition}

\begin{proof}
The parameter \(T=V_0^2/b\) expresses its function field as
\[
 k(T)\left(\sqrt{\frac{T+1}{T}},\sqrt{\frac{T+\iota}{T}},
                  \sqrt{\frac{T+1+\iota}{T}}\right).
\]
The three square classes are geometrically independent, by their
respective odd valuations at \(-1,-\iota,-1-\iota\).
The degree-eight cover has four branch values
\(0,-1,-\iota,-1-\iota\), each with four points of ramification
index two, and is unramified at infinity. Riemann--Hurwitz gives
\(2g-2=-16+16=0\).
The diagonal pencil of the two quadrics has coefficients
\[
 (\iota-1)\lambda+\iota\mu,\quad
 -\iota\lambda-(1+\iota)\mu,\quad\lambda,\quad\mu.
\]
Their zero positions are distinct. A singular point would therefore
have at most one nonzero coordinate, and no such point satisfies the
two equations. Thus the displayed intersection is smooth.

Even coordinate-sign changes, modulo simultaneous sign change, give
a degree-four quotient
\begin{equation}\label{crit:quartic}
 W^2=T(T+1)(T+\iota)(T+1+\iota),\qquad
 W=V_0V_1V_2V_3/b^2.
\end{equation}
A nonidentity element has a projective fixed point only if two
coordinates vanish, which is impossible. After choosing a rational
origin, each such involution is translation by nonzero two-torsion.
Indeed an automorphism \(P\mapsto\alpha(P)+Q\) with
\(\alpha\ne1\) has a fixed point because \(1-\alpha\) is surjective.
The quotient group is therefore the entire two-torsion translation
group. Factoring multiplication by two shows that the quotient is
\(k\)-isomorphic to \(\mathcal C_\iota\).

In \eqref{crit:quartic}, set \(R=(1+\iota)T+\iota\), \(S=2W\).
Then \(S^2=1-R^4\), and
\[
 X_E=\frac{2(1+S)}{R^2},\quad Y_E=\frac{4(1+S)}{R^3},\qquad
 R=\frac{2X_E}{Y_E},\quad S=\frac{X_E^2-4}{X_E^2+4}
\]
identify its function field with \(Y_E^2=X_E^3+4X_E\).
Scaling by \(x=X_E/(1+\iota)^2\), \(y=Y_E/(1+\iota)^3\), and
using \((1+\iota)^4=-4\), gives \(y^2=x^3-x\).
The identifications extend to the smooth projective curves.
\end{proof}

For \(p\equiv1\pmod4\), write \(\chi_p\) for the quadratic
character, with \(\chi_p(0)=0\), and put
\[
 \delta_p=\begin{cases}
 1,&\chi_p(\iota)=\chi_p(1+\iota)=1,\\0,&\text{otherwise}.
 \end{cases}
\]
The zero-coordinate fibres are disjoint; each has four points when
\(\delta_p=1\), and none otherwise. For example \(V_0=0,V_1=1\)
gives \(V_2^2=\iota,V_3^2=1+\iota\). Dividing by a nonzero
coordinate in each other fibre gives the same criterion. Therefore
\begin{equation}\label{crit:count}
 \#U_\iota(\F_p)=\#E(\F_p)-8-16\delta_p.
\end{equation}
Each admissible \(T=c/b\) gives eight projective points. Equivalently,
\((\#E(\F_p)-8-16\delta_p)/8\) counts the \(T\) for which
\(T,T+1,T+\iota,T+1+\iota\) are nonzero and have a common quadratic
character. The square class of \(b\) supplies either common character.

\subsection{The complete prime classification}

\begin{theorem}\label{thm:crit-basic}
For a prime \(p\equiv1\pmod4\), the following are equivalent:
\begin{enumerate}
\item \(p\notin\{5,13\}\).
\item A pattern \eqref{crit:pattern} exists with \(b\ne0\) and all
      four entries nonzero squares in \(\F_p\).
\item For every \(s\ge1\), there are infinitely many primitive
      triples \((X,Y,L)\), with \(0<X,Y<L\) and \(v_p(L)=2s\),
      such that all four radicands in \eqref{crit:integral} have
      roots in \(\Q_p\) of valuation exactly \(s\).
\end{enumerate}
At 5 and 13 no critical pattern exists even with zero entries.
In particular, for a rational four-distance point,
\[
 p\in\{5,13\},\quad p\mid G\quad\Longrightarrow\quad p\mid H.
\]
\end{theorem}

\begin{proof}
Hasse's bound \(\#E(\F_p)\ge p+1-2\sqrt p\)
\cite{silverman2009} exceeds 24 for \(p\ge37\).
Equation~\eqref{crit:count} then gives a point in \(U_\iota\).
The remaining existence cases are
\[
\begin{array}{c|c|c|c|c}
p&\iota&(V_0,V_1,V_2,V_3)&b&c\\\hline
17&4&(1,2,8,4)&3&1\\
29&12&(2,8,12,1)&2&4.
\end{array}
\]
Interchanging \(V_1,V_2\) replaces \(\iota\) by
\(\iota^{-1}=-\iota\), covering the other root of minus one.

For the exclusions, division by a nonzero square reduces \(b\) to
one of its two square classes. The square sets including zero are
\[
 S_5=\{0,1,4\},\qquad S_{13}=\{0,1,3,4,9,10,12\}.
\]
All possibilities after the first two square tests are listed below:
\[
\begin{array}{c|c|c|l}
p&\iota&b&c\ ;\ (c+\iota b,c+(1+\iota)b)\\\hline
5&2&1&0;(2,3)\quad4;(1,2)\\
5&2&2&4;(3,0)\\
13&5&1&0;(5,6)\quad3;(8,9)\quad9;(1,2)\quad12;(4,5)\\
13&5&2&1;(11,0)\quad10;(7,9)\quad12;(9,11).
\end{array}
\]
Every pair contains a nonsquare.

To realize an allowed signed unit pattern \((V_0,V_1,V_2,V_3)\),
put \(c=V_0^2,b=V_1^2-V_0^2\), and fix \(s\ge1\).
Choose positive integers \(X,Y\) with
\begin{equation}\label{crit:lift-coordinates}
 X\equiv-b/2\pmod p,\qquad
 Y\equiv\widetilde\iota X+p^{2s}z\pmod{p^{2s+1}},\qquad
 z\equiv\frac{c}{2\iota X}\pmod p,
\end{equation}
where Hensel's lemma supplies
\(\widetilde\iota^2\equiv-1\pmod{p^{2s+1}}\).
For sufficiently large positive integers \(k\), put
\begin{equation}\label{crit:lift-denominator}
 M=1+pXYk,\qquad L=p^{2s}M.
\end{equation}
Then \(L>X,Y\), \(v_p(L)=2s\), and \(\gcd(X,Y,L)=1\), since
\(p\nmid XY\) and \(M\equiv1\pmod{XY}\).
Expanding \eqref{crit:lift-coordinates} gives
\((X^2+Y^2)/p^{2s}\equiv c\pmod p\).
Subtracting \(2LX\) or \(2LY\), and adding \(L^2\), gives the
other three reductions in \eqref{crit:pattern}.
Hensel's lemma now supplies roots \(A_j=p^sv_j\), with
\(v_j\in\Z_p^\times\) and \(v_j\equiv V_j\pmod p\), of the four
original radicands. Varying \(k\) gives infinitely many coordinate
pairs. Conversely, roots with these valuations reduce to a critical
pattern, to which the exclusions at 5 and 13 apply.
\end{proof}

\begin{corollary}\label{cor:crit-crt}
Let \(\mathcal P\) be a finite set of primes congruent to one modulo
four, excluding 5 and 13, and prescribe \(s_p\ge1\).
There are infinitely many primitive triples \((X,Y,L)\), with
\(0<X,Y<L\), for which \(v_p(L)=2s_p\) and all four original
radicands have roots of valuation \(s_p\) in \(\Q_p\), simultaneously
for every \(p\in\mathcal P\).
\end{corollary}

\begin{proof}
Choose \(X,Y\) by the Chinese remainder theorem using
\eqref{crit:lift-coordinates} at the selected primes. If
\(D=\gcd(X,Y)\), no selected prime divides \(D\). The congruences
\[
 L\equiv p^{2s_p}\pmod{p^{2s_p+1}}\quad(p\in\mathcal P),
 \qquad L\equiv1\pmod D
\]
are compatible. Arbitrarily large solutions give the strict
inequalities and primitivity. The prescribed unit-square reductions
persist, so Hensel's lemma supplies the roots.
\end{proof}

\subsection{The opposite product and the prime 17}

Define
\begin{equation}\label{crit:product}
 t=r_0r_3+r_1r_2=\frac{Q}{G^2H^2},\qquad
 Q=a_0a_3+a_1a_2,\quad Q'=a_0a_3-a_1a_2.
\end{equation}
At a critical prime, multiplication in \(\F_p\) gives
\begin{equation}\label{crit:product-unit}
 QQ'=c\bigl(c+(1+\iota)b\bigr)-(c+b)(c+\iota b)
     =-\iota b^2\ne0.
\end{equation}
Thus both products are units, even with a zero radius reduction.
For \(s=v_p(G)\), the denominator unit in \eqref{crit:product} is a
square, and hence
\begin{equation}\label{crit:product-class}
 v_p(t)=-2s,\qquad [p^{2s}t]=[Q]
 \text{ in }\Q_p^\times/\Q_p^{\times2},\qquad
 \frac{t+1}{t}\in\Q_p^{\times2}.
\end{equation}
The last assertion follows from
\((t+1)/t\in1+p^{2s}\Z_p\).

\begin{theorem}\label{thm:crit-17}
If 17 is a critical content prime of a rational four-distance point,
then none of \(t,2t,t+1,2(t+1)\) is a square in
\(\Q_{17}\). Therefore, if any one is a rational square,
\(17\mid G\) implies \(17\mid H\).
At a critical 17, with \(s=v_{17}(G)\), one also has
\[
 a_0^2+a_3^2\equiv a_1^2+a_2^2\equiv0\pmod{17},\qquad
 17^{2s+1}\mid(2X-L)^2+(2Y-L)^2.
\]
\end{theorem}

\begin{proof}
Take \(\iota=4\). The square set is
\(S_{17}=\{0,1,2,4,8,9,13,15,16\}\), and the scaling classes of
\(b\) are represented by 1 and 3. Successive square tests give
\[
\begin{array}{c|c|c|c}
b&c,c+b\in S_{17}&\text{also }c+4b\in S_{17}
 &\text{also }c+5b\in S_{17}\\\hline
1&\{0,1,8,15,16\}&\{0,15\}&\varnothing\\
3&\{1,13,15,16\}&\{1,13\}&\{1\}.
\end{array}
\]
Thus all coordinates are nonzero, and division by \(a_0\) gives
\([1:\pm2:\pm8:\pm4]\). Hence
\(Q/a_0^2\in\{3,5,12,14\}\), all nonsquares modulo 17.
Interchanging the middle coordinates handles \(\iota=-4\).
Since 2 and \((t+1)/t\) are local squares,
\eqref{crit:product-class} gives the four exclusions.

The same table yields \(c/b=-(1+\iota)/2\), so
\(c=H(X+Y)\) and \(2c+(1+\iota)b=0\).
The identity
\begin{equation}\label{crit:center-identity}
 \frac{(2X-L)^2+(2Y-L)^2}{G^2}
 =4a_0^2-4H(X+Y)+2G^2H^2
\end{equation}
proves the divisibility. At a noncritical 17, its right side is a unit,
since \(17\mid H\) and \(a_0\) is a unit, so the corresponding
valuation is exactly \(2s\).
\end{proof}

\subsection{The normalized product-square cover}

On \(\mathcal C_\iota\) put
\(f=V_0V_3+V_1V_2\), \(f'=V_0V_3-V_1V_2\), so that
\(ff'=-\iota b^2\). Let \(\widetilde{\mathcal C}_\iota\) be the
smooth projective curve obtained by adjoining
\begin{equation}\label{crit:double}
 S^2=2f/V_0^2.
\end{equation}
On \(V_j\ne0\), the same extension is \(S_j^2=2f/V_j^2\).

\begin{proposition}\label{prop:crit-double}
The map \(\widetilde{\mathcal C}_\iota\to\mathcal C_\iota\) is a
geometrically connected unramified double cover of genus-one curves
with rational points. For \(p\equiv1\pmod4\), its number of
rational points above \(b=0\) is
\[
 B_p=\begin{cases}16,&p\equiv1\pmod8,\\8,&p\equiv5\pmod8.
 \end{cases}
\]
Moreover,
\begin{equation}\label{crit:double-count}
 \#\{V\in\mathcal C_\iota(\F_p):b(V)\ne0,\ \chi_p(2f(V))=1\}
 =\frac{\#E(\F_p)-B_p}{2}.
\end{equation}
The character is unchanged by rescaling projective coordinates.
\end{proposition}

\begin{proof}
At \(b=0\), the four zeros of \(f\) are
\([1:\epsilon_1:\epsilon_2:-\epsilon_1\epsilon_2]\).
There \(f'\ne0\), and \(b/V_0^2\) is a uniformizer, because the
map \(T=V_0^2/b\) is unramified at infinity. Thus \(ff'=-\iota b^2\)
makes every zero of \(f/V_0^2\) have order two. On \(V_0=0\), one has \(f=V_1V_2\ne0\), so its poles
also have even order, proving that the cover is
geometrically unramified.

If the extension were geometrically trivial, its square root times
\(V_0\) would be a regular global section of \(\OO_{\mathcal C_\iota}(1)\)
vanishing precisely at the four displayed points. This degree-four
line bundle has four sections by Riemann--Roch; the four independent
coordinate sections span them. But the coordinate matrix of these
points has determinant \(\pm16\), so no linear form vanishes at all
four. The cover is connected. Riemann--Hurwitz gives genus one, and
\(S=\pm2\) above \([1:1:1:1]\) gives rational points.
With compatible origins the map is a degree-two isogeny.

At the four points with \(b=0,f\ne0\), the values
\(2f/V_0^2=\pm4\) give eight rational preimages. At each remaining
point, the normalized local equation has residue
\[
 \left(\frac{S}{b/V_0^2}\right)^2
 =\frac{-2\iota V_0^2}{f'}=\pm\iota.
\]
It gives two rational points precisely when \(\iota\) is square,
equivalently \(p\equiv1\pmod8\). This proves the boundary count.
Isogenous elliptic curves over a finite field have equal point
counts. Apply this to the double cover and
Proposition~\ref{prop:crit-elliptic}. For \(b\ne0\), one has
\(f\ne0\), and each square value has exactly two preimages, proving
\eqref{crit:double-count}.
\end{proof}

\begin{theorem}\label{thm:crit-refined}
For a prime \(p\equiv1\pmod4\), the following are equivalent:
\begin{enumerate}
\item \(p\notin\{5,13,17\}\).
\item There is a signed critical pattern with every \(V_j\ne0\)
      and \(2(V_0V_3+V_1V_2)\in\F_p^{\times2}\).
\item For every \(s\ge1\), there are infinitely many primitive
      strictly interior rational coordinate pairs \((X/L,Y/L)\),
      with \(v_p(L)=2s\), and roots \(A_j\in\Q_p\) of
      \eqref{crit:integral}, each of valuation \(s\), such that
      \(t=(A_0A_3+A_1A_2)/L^2\) satisfies
      \(2t,2(t+1)\in\Q_p^{\times2}\).
\end{enumerate}
\end{theorem}

\begin{proof}
If \(p\equiv5\pmod8\), zero-coordinate base points are impossible
because \(\iota\) is nonsquare. For every such prime \(p\ge29\),
Hasse's lower bound on the product cover exceeds 16 and therefore its
boundary count, giving the required unit point.

If \(p\equiv1\pmod8\), the base has at most sixteen geometric
zero-coordinate points, giving at most 32 points on the double cover.
Together with its sixteen points above \(b=0\), at most 48 points
are excluded. For \(p\ge73\), Hasse's lower bound exceeds 48.
Below 73 the primes in this class are 17 and 41. At 41 take
\(\iota=9\) and \((V_0,V_1,V_2,V_3)=(6,18,2,13)\).
The squared coordinates are \((36,37,4,5)\), so \(b=1\),
\(Q=32\), and \(2Q=23=8^2\pmod{41}\).
Swapping \(V_1,V_2\) preserves \(f\) and handles the other
choice of \(\iota\). The exclusions follow from
Theorems~\ref{thm:crit-basic} and
\ref{thm:crit-17}, including every root sign.

For local realization repeat
\eqref{crit:lift-coordinates}--\eqref{crit:lift-denominator} with
the prescribed signed pattern. Hensel's lemma gives
\(A_j=p^sv_j\), \(v_j\equiv V_j\pmod p\). The product is
\[
 t=\frac{v_0v_3+v_1v_2}{p^{2s}M^2}.
\]
Twice the numerator is a unit square, so \(2t\) is square.
The ratio \((t+1)/t\in1+p^{2s}\Z_p\) supplies the second
root. Conversely, reduction gives an admissible signed pattern.
\end{proof}

\subsection{Excess radius valuations}

\begin{corollary}\label{cor:crit-excess}
Suppose at least one of \(t,2t,t+1,2(t+1)\) is a
rational square. If \(p\mid G\) and \(v_p(A_j)>v_p(G)\) for
some corner, then
\[
 v_p(L)=2v_p(G),\qquad p\equiv1\pmod{16},\qquad
 2^{(p-1)/4}\equiv1\pmod p.
\]
\end{corollary}

\begin{proof}
At a noncritical prime all normalized radii are units. Hence \(p\)
is critical. A square symmetry moves the unique zero to corner zero
and preserves the opposite-product sum. Dividing by \(a_1\) gives
\[
 a_0=0,\quad a_1=1,\quad a_2^2=\iota,\quad
 a_3^2=1+\iota,\quad Q=a_2.
\]
Thus \(p\equiv1\pmod8\), and 2 is square in \(\Q_p\).
By \eqref{crit:product-class}, any of the four square hypotheses
forces \(\chi_p(Q)=1\). But
\[
 \chi_p(a_2)=\iota^{(p-1)/4}=(-1)^{(p-1)/8},
\]
so \(p\equiv1\pmod{16}\). Finally \((1+\iota)^2=2\iota\) gives
\[
 1=\chi_p(1+\iota)
 =2^{(p-1)/4}(-1)^{(p-1)/8}=2^{(p-1)/4}.
\]
\end{proof}

\section{Ordered elliptic sources and fixed-distance fibres}
\label{sec:elliptic-sources}

We label the vertices $00,10,01,11$ by $0,1,2,3$, respectively.
The distance surface is given by
\begin{equation}\label{eq:ell-original}
 r_0^2=x^2+y^2,\quad r_1^2=(x-1)^2+y^2,\quad
 r_2^2=x^2+(y-1)^2,\quad r_3^2=(x-1)^2+(y-1)^2.
\end{equation}
Throughout, a strict rational point has $0<x,y<1$ and
$r_0,r_1,r_2,r_3\in\Q_{>0}$.  This section gives two descriptions of
the complete source of such points: an ordered pair on a congruent-number
elliptic curve, with one final square condition, and a genus-five curve
obtained by fixing a signed pair of distances.

\subsection{Diagonal sums and a simultaneous square model}

Set
\[
 s=r_0+r_3,\qquad d=r_0-r_3,\qquad
 \tau=r_1+r_2,\qquad e=r_2-r_1.
\]
Subtracting the equations in \eqref{eq:ell-original} gives
\[
 sd=2(x+y-1),\qquad \tau e=2(x-y).
\]
The four equations are equivalently encoded by
\begin{equation}\label{eq:ell-diagonal}
 s^2+d^2=\tau^2+e^2,\qquad
 (s^2-2)(2-d^2)=\tau^2e^2,
\end{equation}
together with the return formulas
\begin{equation}\label{eq:ell-diagonal-return}
 \begin{aligned}
 (r_0,r_1,r_2,r_3)&=\tfrac12(s+d,\tau-e,\tau+e,s-d),\\
 x&=\tfrac12+\tfrac14(sd+\tau e),\qquad
 y=\tfrac12+\tfrac14(sd-\tau e).
 \end{aligned}
\end{equation}
For example, eight times the difference between the first two sides of
$x^2+y^2=r_0^2$ is
$-(s^2-2)(2-d^2)+\tau^2e^2$; the other three equations follow by
subtraction and the first equation of \eqref{eq:ell-diagonal}.

\begin{lemma}\label{lem:ell-diagonal-exclusion}
Neither diagonal of the square contains a strict rational four-distance
point.  The square symmetries let us choose the chamber
\begin{equation}\label{eq:ell-chamber}
 x+y>1,\qquad x>y.
\end{equation}
\end{lemma}
\begin{proof}
On $x=y$
one has $r_0=\sqrt2x$, and on $x+y=1$ one has
$r_1=\sqrt2(1-x)$.  The square symmetries act on
$(x+y-1,x-y)$ by signed permutations.
\end{proof}
In \eqref{eq:ell-chamber} we have
\[
 s>d>0,\qquad \tau>e>0,\qquad s,\tau>\sqrt2.
\]
The strict triangle inequalities give $d^2,e^2<2$.  In fact
\begin{equation}\label{eq:ell-J-negative}
 d^2+e^2\leq
 2\bigl((x+y-1)^2+(x-y)^2\bigr)
 =4\bigl((x-\tfrac12)^2+(y-\tfrac12)^2\bigr)<2.
\end{equation}

Define
\begin{equation}\label{eq:ell-pure-coordinates}
 p=\frac{2-e^2}{d^2},\quad
 \beta=\frac{s^2}{2-e^2},\quad
 U=\frac{s}{d},\quad V=\frac{2e}{d},\quad
 W=\frac{2\tau}{d},\quad T=\frac2d.
\end{equation}
Here $p,\beta>1$, and \eqref{eq:ell-diagonal} becomes
\begin{equation}\label{eq:ell-pure}
 \begin{aligned}
 U^2&=p\beta,\qquad
 V^2=2(p-1)(\beta-1),\\
 W^2&=2(p+1)(\beta+1),\qquad
 T^2=(p-1)\beta+p+1.
 \end{aligned}
\end{equation}
Conversely these four equations give
\begin{equation}\label{eq:ell-pure-return}
 \begin{aligned}
 (r_0,r_1,r_2,r_3)
 &=\left(\frac{U+1}{T},\frac{W-V}{2T},
          \frac{W+V}{2T},\frac{U-1}{T}\right),\\
 x&=\frac12+\frac{U+WV/4}{T^2},\qquad
 y=\frac12+\frac{U-WV/4}{T^2}.
 \end{aligned}
\end{equation}
Thus all four square roots remain attached to the same $p,\beta$.

\subsection{The complete ordered elliptic source}

For a positive squarefree integer $n$, let
\[
 E_n:\quad Y^2=X(X-n)(X+n).
\]
An ordered source consists of two points
$Q=(a,b)$ and $R=(c,\ell)$ on $E_n$, and three rational roots
$u,v,h$ satisfying
\begin{equation}\label{eq:ell-ordered-roots}
 u^2=-\frac{ac}{n},\qquad
 v^2=-\frac{(a-n)(c-n)}{n},\qquad
 h^2=\frac{a+c}{2(c+n)}.
\end{equation}
The third product root is fixed by the ordinates:
\begin{equation}\label{eq:ell-w-ordinate}
 w=\frac{b\ell}{nuv},\qquad w^2=(a+n)(c+n).
\end{equation}
In particular the first two equations specify the ordered square classes
\[
 [ac]=[-n],\qquad [(a-n)(c-n)]=[-n]
 \quad\text{in }\Q^*/\Q^{*2}.
\]
The condition on $h$ is a further square condition on this same pair.

\begin{theorem}[Exact ordered-source correspondence]
\label{thm:ell-exact-source}
Strict rational points in the chamber \eqref{eq:ell-chamber} are in
bijection with the union, over positive squarefree $n$, of the rational
solutions of \eqref{eq:ell-ordered-roots} for which
\begin{equation}\label{eq:ell-source-signs}
 -n<a<0<c-n,\qquad b,\ell,u,v,h>0,
\end{equation}
and
\begin{equation}\label{eq:ell-strict-metric}
 \frac{2u}{a+c}\left(-\frac ba+\frac\ell c\right)<2.
\end{equation}
For this correspondence the final square is
\begin{equation}\label{eq:ell-final-square}
 f(Q,R):=\frac{a+c}{2(c+n)}
       =\frac{1}{(r_0+r_3)^2}.
\end{equation}
\end{theorem}

\begin{proof}
We first give rational maps that preserve the actual roots and ordinates.
From an ordered source set
\begin{equation}\label{eq:ell-forward-pair}
 \begin{aligned}
 p&=-\frac na,& k&=\frac b{a^2},&
 \beta&=\frac{c+n}{c-n},\\
 U&=\frac{(n-a)w}{bv},&
 V&=-\frac{2b}{av},&
 W&=\frac{2(a-n)u}{av},& T&=2Uh.
 \end{aligned}
\end{equation}
The curve equations and \eqref{eq:ell-ordered-roots} imply
\eqref{eq:ell-pure} and
\begin{equation}\label{eq:ell-k}
 nk^2=p(p^2-1).
\end{equation}
For instance, $k^2=b^2/a^4=(a^2-n^2)/a^3$ gives
\eqref{eq:ell-k}; the other identities follow by cancelling the factors
$a,a\pm n,c,c\pm n$.
Conversely, from \eqref{eq:ell-pure} and \eqref{eq:ell-k} put
\begin{equation}\label{eq:ell-inverse-pair}
 \begin{aligned}
 a&=-\frac np,& b&=ka^2,&
 c&=n\frac{\beta+1}{\beta-1},&
 \ell&=\frac nk\frac{UVW}{(\beta-1)^2},\\
 u&=\frac{(p-1)W}{kV},&
 v&=\frac{2(p^2-1)}{kV},&
 h&=\frac T{2U}.&&
 \end{aligned}
\end{equation}
These formulas give, separately,
\[
 \begin{aligned}
 b^2&=\frac{n^3(p^2-1)}{p^3}=a(a-n)(a+n),\\
 \ell^2&=\frac{4n^3\beta(\beta+1)}{(\beta-1)^3}
          =c(c-n)(c+n),\\
 u^2&=\frac{n(\beta+1)}{p(\beta-1)}=-\frac{ac}{n},\\
 v^2&=\frac{2n(p+1)}{p(\beta-1)}=-\frac{(a-n)(c-n)}n,\\
 h^2&=\frac{(p-1)\beta+p+1}{4p\beta}
       =\frac{a+c}{2(c+n)}.
 \end{aligned}
\]
Inserting \eqref{eq:ell-inverse-pair} into
\eqref{eq:ell-forward-pair} returns $U,V,W,T$ with their signs.
Conversely it returns $b$ and $\ell$ as well as $a,c,u,v,h$.

To obtain $n$ and $k$ from a strict original point, set
\[
 t=r_0r_3+r_1r_2=\frac{s^2-e^2}{2},\quad
 J=d^2+e^2-2,\quad K=2-d^2,
\]
and choose the unique positive squarefree integer $n$ and positive
$q\in\Q$ with $t^2-1=nq^2$.  This is possible because
\eqref{eq:ell-diagonal} gives
\begin{equation}\label{eq:ell-theta-id}
 t^2-1=\left(\frac{de}{J}\right)^2K(2-e^2)>0.
\end{equation}
For completeness, writing $D=d^2$ and $E=e^2$, the two equations of
\eqref{eq:ell-diagonal} give
\[
 s^2=-\frac{(2-E)(2+E-D)}{D+E-2},\qquad
 \tau^2=-\frac{(2-D)(2+D-E)}{D+E-2}.
\]
These expressions prove \eqref{eq:ell-theta-id} and the identity
\[
 \frac{p(p^2-1)}{t^2-1}
   =\left(\frac{\tau J^2}{d^4eK}\right)^2.
\]
Consequently
\[
 k=\frac{q\tau J^2}{d^4eK}>0
\]
satisfies \eqref{eq:ell-k}.  Equations
\eqref{eq:ell-pure-coordinates} and \eqref{eq:ell-inverse-pair}
now give the required source.  All their denominators are nonzero
by \eqref{eq:ell-J-negative}, $d,e>0$, and $p,\beta>1$.

On the source, direct use of \eqref{eq:ell-forward-pair} gives
\begin{equation}\label{eq:ell-metric-identities}
 q=\frac{2u}{a+c},\qquad
 t=\frac{c-a}{c+a},\qquad
 \frac{WV}{T^2}=-\frac{qb}{a},\qquad
 \frac{4U}{T^2}=\frac{q\ell}{c}.
\end{equation}
Thus \eqref{eq:ell-strict-metric} is precisely
$(4U+WV)/T^2<2$.  The signs in \eqref{eq:ell-source-signs}
give $U>1$, $W>V>0$ and $T>0$.  All four radii in
\eqref{eq:ell-pure-return} are positive.  That formula shows that
the metric inequality is equivalent to $x<1$ in the chamber;
it also gives $x>1/2$ and $|y-1/2|<1/2$.
Finally $r_0+r_3=2U/T=1/h$ proves
\eqref{eq:ell-final-square}.  The displayed inverse formulas and the
uniqueness of $n,q>0$ prove bijectivity.
\end{proof}

The formulas above also describe the signed algebraic covers.  Let $S$
denote the surface \eqref{eq:ell-original}, and let $X_n$ be the
normalization of $E_n\times E_n$ in the field obtained by adjoining
$u,v,h$ from \eqref{eq:ell-ordered-roots}.
\begin{proposition}\label{prop:ell-original-cover}
There is an equality of function fields
\[
 \Q(X_n)=\Q(S)(q),\qquad q^2=\frac{t^2-1}{n},
 \qquad t=r_0r_3+r_1r_2.
\]
The extension has degree two.
\end{proposition}
\begin{proof}
The four radicands in \eqref{eq:ell-original} have independent
square classes over $\overline\Q(x,y)$: each factors into two lines,
and each of these eight lines occurs for only one corner.  Hence
$S$ is geometrically integral.
The preceding maps are rational identities before signs are chosen.
Their inverse recovers $q$ from the simultaneous-square coordinates by
\[
 q=\frac{kVW}{(p^2-1)T^2}.
\]
The two root classes defining $u,v$ are geometrically independent:
their half-divisor classes on $E_n\times E_n$ restrict on either
factor to the independent classes
$[(0,0)-O]$ and $[(n,0)-O]$ in $E_n[2]$.
Their degree-four cover is therefore connected and unramified.
The zero divisor $a+c=0$ has simple generic multiplicity, which
remains odd after this unramified pullback.  Adjoining $h$ is
therefore ramified there, so this last quadratic extension is also
connected.  These statements hold over $\overline\Q$.
The resulting involution
\[
 (a,b,c,\ell,u,v,h)\longmapsto(a,-b,c,-\ell,-u,-v,h)
\]
fixes $U,V,W,T$, hence fixes $S$, and sends $q$ to $-q$.
It is nontrivial, proving the degree assertion.  The normalization of
the graph gives the corresponding birational identification of proper
models.
\end{proof}

\subsection{Fixing an opposite pair of distances}

Let $c_0=\mu/\nu$ be a fixed positive rational number, with coprime
positive integers $\mu,\nu$.  Suppose
\[
 c_0=r_3+\varepsilon r_0,\qquad
 d_0=r_3-\varepsilon r_0,\qquad \varepsilon\in\{1,-1\};
\]
for a difference we orient the pair so that $r_3>r_0$.
The strict real bounds give $0<c_0<2$.
With
\[
 A=2\nu^2-\mu^2,\quad B=4\nu^2-\mu^2,\quad
 P=\nu d_0,\quad Z=\nu^2\bigl((1-2x)-(1-2y)\bigr),
\]
and $Q=2\nu r_1$, $R=2\nu r_2$, the full equations are
\begin{equation}\label{eq:ell-fixed-opposite}
 Z^2=A(P^2-2\nu^2),\qquad
 Q^2=P^2+\mu^2+2Z,\qquad R^2=P^2+\mu^2-2Z.
\end{equation}
Indeed, if $u_0=1-2x$ and $v_0=1-2y$, then
\[
 c_0d_0=u_0+v_0,\qquad
 c_0^2+d_0^2=2+u_0^2+v_0^2,
\]
so $(u_0-v_0)^2=(2-c_0^2)(d_0^2-2)$.
The inverse is
\begin{equation}\label{eq:ell-fixed-opposite-return}
 \begin{aligned}
 x&=\frac12-\frac{\mu P+Z}{4\nu^2},&
 y&=\frac12-\frac{\mu P-Z}{4\nu^2},\\
 r_3&=\frac{\mu+P}{2\nu},&
 r_0&=\frac{\varepsilon(\mu-P)}{2\nu},&
 r_1&=\frac{|Q|}{2\nu},&r_2&=\frac{|R|}{2\nu}.
 \end{aligned}
\end{equation}
Its strict locus is specified by positivity of these radii and
\begin{equation}\label{eq:ell-fixed-strict}
 |\mu P+Z|<2\nu^2,\qquad |\mu P-Z|<2\nu^2.
\end{equation}
Substitution into \eqref{eq:ell-original} proves the inverse.

\begin{proposition}\label{prop:ell-opposite-genus}
For every rational $0<c_0<2$, the smooth projective curve determined
by \eqref{eq:ell-fixed-opposite} is geometrically connected of genus five.
\end{proposition}
\begin{proof}
Neither $A$ nor $B$ is zero.  Over the $Z$-line the three radicands,
after multiplication by the same nonzero constant $A$, are
\[
 f_0=Z^2+2A\nu^2,\quad
 f_+=(Z+A)^2+2A\mu^2,\quad
 f_-=(Z-A)^2+2A\mu^2.
\]
All three have two simple roots.  A common root of $f_0,f_+$
would be $Z=-\mu^2/2$, where $f_0=B^2/4\ne0$;
the pair $f_0,f_-$ is identical after $Z\mapsto-Z$.
A common root of $f_+,f_-$ would be $Z=0$, where
$f_+=A(2\nu^2+\mu^2)\ne0$.
Thus there are six distinct branch values.  The three square classes
are independent over $\overline\Q(Z)$, since each has simple zeros
belonging to no other radicand.  The cover has degree eight, with
four points of ramification index two above each branch value;
infinity is unramified.  Riemann--Hurwitz gives
$2g-2=-16+6\cdot4=8$.
\end{proof}

Multiplication of the last two equations in
\eqref{eq:ell-fixed-opposite} produces an elliptic quotient:
\begin{equation}\label{eq:ell-fixed-elliptic}
 F_{\mu,\nu}:\quad
 \eta^2=\xi^3+\alpha\xi^2+B^2\xi,
 \qquad \alpha=2(3\mu^2-4\nu^2),
 \qquad (\xi,\eta)=(P^2,PQR).
\end{equation}
Its cubic discriminant is nonzero because
$\alpha^2-4B^2=-32\mu^2A\ne0$.
Every strict image satisfies
\begin{equation}\label{eq:ell-quotient-interval}
 0<\xi<B.
\end{equation}
For the upper bound use $c_0^2+d_0^2<4$.
For the lower bound, $P=0$ for a sum would place the original point
on $x+y=1$, while for a difference $P=\nu(r_3+r_0)>0$.
The quotient always has the four torsion points
\begin{equation}\label{eq:ell-order-four}
 O,\quad(0,0),\quad(B,2\mu B),\quad(B,-2\mu B).
\end{equation}
The duplication formula
\begin{equation}\label{eq:ell-duplication}
 \xi(2S)=\frac{(\xi(S)^2-B^2)^2}{4\eta(S)^2}
\end{equation}
shows that the last two double to $(0,0)$.

\begin{theorem}[Uniform torsion for odd pair parameters]
\label{thm:ell-torsion}
If $\mu,\nu$ are odd, the rational torsion subgroup of
$F_{\mu,\nu}$ is cyclic of order four, with points
\eqref{eq:ell-order-four}.  Any strict four-distance point with
opposite-pair value $\mu/\nu$ gives a point of infinite order on
$F_{\mu,\nu}$.
\end{theorem}
\begin{proof}
Put $S=(0,0)$ and $T=(B,2\mu B)$. The quadratic factor
$\xi^2+\alpha\xi+B^2$ has discriminant
$32\mu^2(\mu^2-2\nu^2)$, of exact $2$-adic valuation five.
Consequently $S$ is the unique nonzero rational point killed by two.
The duplication formula gives $2T=S$, so $T$ has exact order four.
We have
\[
 B\equiv3\pmod8,\qquad \alpha\equiv14\pmod{16},\qquad
 B^2\equiv9\pmod{16}.
\]

We first exclude odd torsion. For the integral Weierstrass equation over
$\Z_2$, let $G_1$ be the subgroup of $F_{\mu,\nu}(\Q_2)$ reducing
to infinity. Its parameter $z=-\xi/\eta$ lies in $2\Z_2$.
The formal group law has integral coefficients and linear part
$z_1+z_2$. Indeed, in coordinates $z=-\xi/\eta$ and $w=-1/\eta$
the equation has a unit derivative in $w$ at $(0,0)$; successive
coefficient comparison solves for $w$ as an integral power series in
$z$. Substitution into the addition law gives
\[
 [m]z=mz+z^2g_m(z),\qquad g_m(z)\in\Z_2[[z]].
\]
If $m$ is odd and $z\ne0$, the two terms on the right have different
valuations, with the first strictly smaller. Thus $G_1$ has no
nonzero odd torsion. This argument uses the smooth point at infinity
and applies to the present integral equation at $2$.

For an arbitrary rational point $P$, put $Q=2P$. If $Q$ is finite,
its abscissa $\xi$ is a rational square by
\eqref{eq:ell-duplication}. If $\xi=0$, then $Q=S$ and $4P=O$.
Suppose $\xi\ne0$, and write $v=\ord_2(\xi)$.
The case $v=0$ is impossible. The curve equation makes the ordinate
integral, and a unit square is $1$ or $9$ modulo $16$, so
\[
 \xi(\xi^2+\alpha\xi+B^2)\equiv8\pmod{16},
\]
whereas a square in $\Z_2$ cannot have this residue.

If $v>0$, the term $B^2\xi$ has uniquely least valuation in the
curve equation. Hence $2\ord_2(\eta)=v$, and
\eqref{eq:ell-duplication} gives
$\ord_2(\xi(2Q))=-v-2\le-4$.
If $v<0$, the cubic term has uniquely least valuation, so
$2\ord_2(\eta)=3v$ and
$\ord_2(\xi(2Q))=v-2\le-4$.
The numerator and denominator of the duplication formula are nonzero
in both cases. A point with a pole in its abscissa belongs to $G_1$.
We have therefore proved that $4P\in G_1$ for every rational point
$P$, including the cases $4P=O$.
If $P$ has odd order, then $4P=O$ by the preceding formal-group
argument, and coprimality with four gives $P=O$.

Every rational torsion point has $2$-power order. If $4P=O$, then
$2P$ is $O$ or $S$. In the first case $P$ is $O$ or $S$; in the
second, $2(P-T)=O$, so $P$ is $T$ or $-T$.
There is no point of order eight: a half of $T$ or $-T$ would make
$B$ a rational square by \eqref{eq:ell-duplication}, contrary to
$B\equiv3\pmod8$. A point of larger $2$-power order would have a
multiple of order eight. Thus the torsion subgroup is precisely
$\{O,S,T,-T\}$.
Finally, every strict image satisfies
\eqref{eq:ell-quotient-interval}, whereas the finite torsion abscissas
are $0$ and $B$. This proves the infinite-order assertion.
\end{proof}

\begin{corollary}\label{cor:ell-rank-zero}
If $\mu,\nu$ are odd and $\rank F_{\mu,\nu}(\Q)=0$, then
$\mu/\nu$ is neither an opposite-pair sum nor a nonzero absolute
difference at a strict rational four-distance point.
\end{corollary}
\begin{proof}
Such a point would give a rational point of $F_{\mu,\nu}$ in the interval
\eqref{eq:ell-quotient-interval}. Rank zero makes every rational point
torsion, whereas Theorem~\ref{thm:ell-torsion} places all torsion points
outside that interval.
\end{proof}

\subsection{Adjacent pairs and fixed-pair finiteness}

For an adjacent pair choose signed radii $R_0,R_1$ with
$R_0+R_1=c_0>0$, where $R_i^2=r_i^2$.  Write $\lambda=c_0^2$.
A strict adjacent sum is greater than one and an absolute difference
is less than one, so a nonzero permitted value has
$\lambda\ne0,1$.
With $U_0=2x-1$ and $V_0=2y$, the first two distance equations
give the conic
\begin{equation}\label{eq:ell-adjacent-conic}
 (\lambda-1)U_0^2+\lambda V_0^2
       -\lambda(\lambda-1)H_0^2=0.
\end{equation}
It is smooth and contains $(c_0,0,1)$.  In the chart $H_0=1$,
\[
 R_0=\frac{c_0+U_0/c_0}{2},\qquad
 R_1=\frac{c_0-U_0/c_0}{2}.
\]
The remaining roots are given by
\[
 f_2=\frac{(U_0+H_0)^2+(V_0-2H_0)^2}{H_0^2},\qquad
 f_3=\frac{(U_0-H_0)^2+(V_0-2H_0)^2}{H_0^2}.
\]
Adjoin both square roots to the conic function field.  On its affine
chart the return is
\[
 x=\frac{U_0+1}{2},\quad y=\frac{V_0}{2},\quad
 (r_0,r_1,r_2,r_3)
 =\left(|R_0|,|R_1|,\frac{|\sqrt{f_2}|}{2},
                           \frac{|\sqrt{f_3}|}{2}\right).
\]

\begin{proposition}\label{prop:ell-adjacent-genus}
For every rational $c_0>0$ with $c_0\ne1$, this normalized cover
is geometrically connected of genus five.
\end{proposition}
\begin{proof}
Over $\overline\Q$, the numerator of $f_2$ factors into the two lines
$U_0+H_0\pm i(V_0-2H_0)=0$.
For the line with sign $\delta=\pm1$, evaluation of the dual conic is
\[
 \frac{1}{\lambda-1}-\frac1\lambda
 -\frac{(1-2\delta i)^2}{\lambda(\lambda-1)}
 =\frac{4+4\delta i}{\lambda(\lambda-1)}\ne0.
\]
Both intersections are therefore transverse.  The two lines meet at
$(-1,2,1)$, which would belong to the conic only if
$\lambda^2-6\lambda+1=0$; this is impossible for rational $\lambda$.
Thus $f_2$ has four distinct simple zeros.  The same calculation
with $U_0\mapsto-U_0$ proves this for $f_3$.

The two zero sets are disjoint.  Their numerator difference is
$4U_0H_0$.  At $H_0=0$ a common zero would satisfy
$U_0^2+V_0^2=0$, which together with the conic forces
$U_0=V_0=0$.  In the chart $H_0=1$, a common zero has
$U_0=0$ and $V_0=2\pm i$, whence the conic requires
$\lambda=4\pm4i$, again impossible.
The two points of the conic at infinity are distinct and neither
numerator vanishes there; both poles of each $f_i$ have order two.
The radicands are consequently independent square classes, and
the degree-four cover has eight branch values, each contributing
two to the ramification divisor.  Hence
$2g-2=4(-2)+8\cdot2=8$.
\end{proof}

\begin{theorem}[Fixed-pair finiteness]\label{thm:ell-fixed-pair-finite}
Fix two distinct vertices, signs $\varepsilon_i,\varepsilon_j\in\{1,-1\}$,
and a nonzero rational number $c_0$.
There are only finitely many strict rational four-distance points satisfying
$\varepsilon_i r_i+\varepsilon_j r_j=c_0$.
\end{theorem}
\begin{proof}
Changing both signs and applying a square symmetry reduces the pair
to one of the two constructions above.  Their return formulas show
that every strict point is the image of a rational point on the
corresponding smooth projective genus-five curve.  Faltings'
theorem \cite{faltings1983} makes this rational-point set finite.
\end{proof}

\subsection{The fibre with opposite difference one fifth}
\label{subsec:ell-one-fifth}

For $c_0=1/5$, take $\mu=1$, $\nu=5$, $\varepsilon=-1$ and
$Z=7z$ in \eqref{eq:ell-fixed-opposite}.  The curve becomes
\begin{equation}\label{eq:ell-one-fifth}
 C:\quad P^2=z^2+50,\qquad
 Q^2=z^2+14z+51,\qquad R^2=z^2-14z+51.
\end{equation}
Its strict positive branch has $P,Q,R>0$ and
\begin{equation}\label{eq:ell-one-fifth-return}
 \begin{aligned}
 x&=\tfrac12-\frac{P+7z}{100},&
 y&=\tfrac12-\frac{P-7z}{100},\\
 (r_0,r_1,r_2,r_3)&=\tfrac1{10}(P-1,Q,R,P+1),&
 |P+7z|&<50,&|P-7z|&<50.
 \end{aligned}
\end{equation}
The projective model in coordinates $[P:Q:R:z:T]$ is obtained by
homogenizing the three equations.  Above $T=0$ it has the eight
rational points $[\pm1:\pm1:\pm1:1:0]$.  These points are smooth:
in the chart $z=1$ the derivatives in $P,Q,R$ form an invertible
diagonal matrix, and $T$ is a local parameter.

\begin{proposition}[Elliptic decomposition of the fixed fibre]
\label{prop:ell-one-fifth-jacobian}
There is a $\Q$-isogeny
\[
 \Jac(C)\sim_{\Q} E_1\times F\times E_3\times E_4^2,
\]
where
\begin{align*}
 E_1&:\ y^2=x^3-4606x^2+6245001x,\\
 F&:\ y^2=x^3-194x^2+9801x,\\
 E_3&:\ y^2=x^3-9902x^2+25492401x,\\
 E_4&:\ y^2=x^3-202x^2+9801x.
\end{align*}
\end{proposition}
\begin{proof}
Put $f_0=z^2+50$, $f_\pm=z^2\pm14z+51$.
The seven nontrivial characters of the sign-change group
$(\Z/2\Z)^3$ give quotients $w^2=\prod_{j\in I}f_j$.
Their genera are $0,0,0,1,1,1,2$, according as $|I|$ is one,
two, or three.  The pullbacks of their differential spaces are
the respective character spaces in $H^0(C,\Omega_C^1)$, and
their dimensions sum to five.  The product of the associated
Jacobian pullback maps is therefore an isogeny over $\Q$.

The quotient $w^2=f_+f_-=z^4-94z^2+2601$ maps with degree two
to $y^2=x^3-94x^2+2601x$ by $(x,y)=(z^2,zw)$.
The scaling $(x,y)\mapsto(49x,343y)$ gives $E_1$.
The two quotients $w^2=f_0f_\pm$ are isomorphic by $z\mapsto-z$.
For the plus sign the substitutions
\[
 W_0=w+z^2+7z+26,\qquad
 x=2W_0+49,\qquad y=(4z+14)W_0+336
\]
give $E_4$.  Their rational inverse on $W_0\ne0$ is
\[
 W_0=(x-49)/2,\qquad
 z=\frac{y-14W_0-336}{4W_0},\qquad
 w=W_0-z^2-7z-26.
\]

The remaining genus-two curve is
\[
 D:\quad w^2=(z^2+50)(z^4-94z^2+2601).
\]
Its involutions $(z,w)\mapsto(-z,w)$ and
$(z,w)\mapsto(-z,-w)$ have elliptic quotients.
The first is $(x,y)=(z^2+50,w)$ on $F$.
For the second put $\zeta=z^2$ and $v=zw$, so that
$v^2=\zeta(\zeta+50)(\zeta^2-94\zeta+2601)$.
More generally, for
$v^2=\zeta(\zeta+c_1)(\zeta^2+a_1\zeta+b_1)$, the substitutions
\[
 X=\frac{c_1b_1}{\zeta}+b_1,\qquad
 Y=\frac{c_1b_1v}{\zeta^2}
\]
give
\[
 Y^2=X^3+(c_1a_1-2b_1)X^2
          +b_1(b_1-c_1a_1+c_1^2)X.
\]
Taking $(c_1,a_1,b_1)=(50,-94,2601)$ yields $E_3$.
The two one-dimensional differential eigenspaces on $D$ prove
$\Jac(D)\sim_\Q F\times E_3$, completing the decomposition.
\end{proof}

\begin{example}\label{ex:ell-nonlift}
The elliptic quotient $F$ contains
\[
 (\xi,\eta)=\left(
 \frac{1476225}{18496},-\frac{589188735}{2515456}\right).
\]
This follows by direct substitution in its cubic equation.
Its first coordinate is $(1215/136)^2$ and satisfies $50<\xi<99$.
A lift to \eqref{eq:ell-one-fifth} would require
\[
 Z^2=49(\xi-50)=\left(\frac{35}{136}\right)^2 22057.
\]
Since $148^2<22057<149^2$, the required rational root does not exist.
Thus the interval \eqref{eq:ell-quotient-interval} and a square
elliptic coordinate do not characterize the image of $C(\Q)$;
the three simultaneous equations in \eqref{eq:ell-one-fifth}
give the complete lift.
\end{example}

\section{A common squareclass descent}\label{sec:common-descent}

The sum of the products of opposite distances gives a useful one-parameter
family of covering curves.  The four square roots in this family determine
one point of the square.  In particular, the squareclass of their common
coordinate can be compared at different places.

\subsection{The four-root curve and its return map}

Let $t\in\Q\setminus\{0,1,-1\}$, and put
\[
 s=t-1,\qquad T=t+1,\qquad d=sT=t^2-1.
\]
Let $\mathcal U_t$ be the affine curve
\begin{equation}\label{eq:common-four-roots}
 P^2=\ell+T,\qquad C^2=\ell-s,\qquad
 M^2=\frac{s(T-\ell)}\ell,\qquad B^2=\frac{T(\ell+s)}\ell,
 \qquad \ell PCMB\ne0.
\end{equation}
We denote its smooth projective completion by $D_t$.  The order of the
four roots will be fixed throughout this section.

For a point whose distances from $(0,0),(1,0),(0,1),(1,1)$ are
$r_0,r_1,r_2,r_3$, respectively, the coordinates are
\begin{equation}\label{eq:common-from-distances}
\begin{gathered}
 t=r_0r_3+r_1r_2,\qquad P=r_0+r_3,\qquad M=r_3-r_0,\\
 B=r_1+r_2,\qquad C=r_2-r_1,\qquad \ell=P^2-t-1.
\end{gathered}
\end{equation}
To verify the forward construction and its range, put $u=1-2x$,
$v=1-2y$, $A=r_0r_3$, and $A'=r_1r_2$.
Expansion of the four squared-distance identities gives
\[
 A^2-(A')^2=-uv,\qquad
 (A+A')^2+(A-A')^2=1+(u^2+v^2)^2/4.
\]
Thus $t^2=(A+A')^2$ is the larger root of
\[
 F_t(Z)=Z^2-\bigl(1+(u^2+v^2)^2/4\bigr)Z+u^2v^2.
\]
The diagonal exclusion in Lemma~\ref{lem:ell-diagonal-exclusion} gives
$u^2\ne v^2$, and hence $F_t(1)=-(u^2-v^2)^2/4<0$.
For $a=u^2,b=v^2\in[0,1)$,
\[
 16F_t(5/4)=5-5a^2+6ab-5b^2>0.
\]
Indeed this polynomial is strictly concave on $[0,1]^2$, with vertex
values $5,0,0,1$, and its two zero vertices are excluded.
The axis of $F_t$ is below $1$, so $1<t<\sqrt5/2$.
Furthermore
\[
 PM=u+v,\quad BC=v-u,\quad
 P^2-C^2=B^2-M^2=2t,\quad P^2+M^2=u^2+v^2+2.
\]
These identities give $(P^2-T)(M^2+s)=d$ by expansion.
The first two equations of \eqref{eq:common-four-roots} follow from the
definition of $\ell$; the last two then follow from this product
identity and $B^2-M^2=2t$.  The diagonal exclusion also gives
$C,M\ne0$.  Since $s>0$ and $d>0$, the product identity first gives
$\ell>0$, and the equations for $C^2,M^2$ then give $s<\ell<T$.

Conversely, a tuple in \eqref{eq:common-four-roots} determines
\begin{equation}\label{eq:common-return}
\begin{gathered}
 u=\frac{PM-BC}{2},\qquad v=\frac{PM+BC}{2},\qquad
 x=\frac{1-u}{2},\qquad y=\frac{1-v}{2},\\
 r_0=\frac{P-M}{2},\quad r_3=\frac{P+M}{2},\quad
 r_1=\frac{B-C}{2},\quad r_2=\frac{B+C}{2}.
\end{gathered}
\end{equation}
These formulas give the four squared-distance identities over every field
of characteristic zero.  Indeed, the root equations imply
\[
 B^2-M^2=P^2-C^2=2t,\qquad
 P^2M^2+B^2C^2=2(P^2+M^2-2).
\]
Consequently $u^2+v^2=P^2+M^2-2=B^2+C^2-2$.  Since
$u+v=PM$ and $u-v=-BC$, substitution into
$\bigl((1-u)^2+(1-v)^2\bigr)/4$ and its three translates gives
$r_0^2,r_1^2,r_2^2,r_3^2$.  Moreover,
\[
 r_0r_3+r_1r_2=\frac{P^2-M^2+B^2-C^2}{4}=t.
\]
For $t>1$, the strict real return chamber is
\begin{equation}\label{eq:common-chamber}
 s<\ell<T,\qquad P,B>0,\qquad |PM|+|BC|<2.
\end{equation}
In this chamber $P>|M|$ and $B>|C|$, so all four returned distances are
positive.  The last inequality is equivalent to $|u|,|v|<1$, and hence to
$0<x,y<1$.  These identities also show directly that the transformation
\begin{equation}\label{eq:common-reflection}
 (\ell,P,C,M,B)\longmapsto(d/\ell,B,M,C,P)
\end{equation}
is the reflection $x\mapsto1-x$ on the square.

\begin{proposition}\label{prop:common-genus-boundary}
The curve $D_t$ is geometrically connected of genus nine.  The map
$\ell:D_t\to\PP^1$ has degree sixteen and is branched precisely at
\[
 0,\quad\infty,\quad s,\quad T,\quad-s,\quad-T.
\]
If $t>1$, none of the points above these six values is rational.
\end{proposition}
\begin{proof}
Over $\overline\Q(\ell)$, the four radicands in
\eqref{eq:common-four-roots} have independent squareclasses: their
valuations at $-T,s,T,-s$, respectively, each isolate one of the four
classes.  The extension therefore has degree sixteen.  At each of the
six displayed branch values the inertia group has order two.  For
example, the valuation vectors at $0$ and $\infty$, in the order
$(P,C,M,B)$, are $(0,0,1,1)$ and $(1,1,0,0)$ modulo two.  Thus
Riemann--Hurwitz gives
\[
 2g(D_t)-2=-32+6\cdot8=16.
\]
The following table records a regular function on the normalization and
its value at each omitted fibre:
\[
\begin{array}{c|rrrrrr}
 \ell&s&T&-s&-T&0&\infty\\\hline
 \text{regular square}&M^2&C^2&P^2&B^2&C^2&M^2\\
 \text{value}&2&2&2&2&-s&-s.
\end{array}
\]
The first four require $\sqrt2\in\Q$, and the last two require a square
root of the negative number $-s$.  This excludes every rational point
on the six fibres.
\end{proof}

\subsection{Prime support and the ramified unit conditions}

Write $[a]_K$ for the class of $a\in K^\times$ modulo squares; omit the
subscript when $K=\Q$.  For a nonzero rational number $d$, set
\[
 S(d)=\{p:\ p\text{ is odd and }v_p(d)>0\text{ is odd}\}.
\]

\begin{lemma}[Support of the common coordinate]\label{lem:common-support}
If $(\ell,P,C,M,B)\in\mathcal U_t(\Q)$, then
\[
 [\ell]\in\langle[-1],[2],[p]:p\in S(d)\rangle.
\]
For $t>1$, the sign is positive and the generator $[-1]$ can be omitted.
\end{lemma}
\begin{proof}
Fix an odd prime $p$ and put $m=v_p(t)$ and $n=v_p(\ell)$.
Suppose first that $m<0$, so $v_p(s)=v_p(T)=m$.
If $n>m$, the first and third equations of
\eqref{eq:common-four-roots} have valuations $m$ and $2m-n$;
both must be even, so $n$ is even.  If $n<m$, the fourth and first
have valuations $m$ and $n$, giving the same conclusion.  If $n=m$,
the reductions of $P^2/t$ and $C^2/t$ differ by $2$.
At least one is a unit; its square has valuation $m=n$, which is even.

Now let $m\ge0$.  If both $s$ and $T$ are units, then $n<0$ is even
by the first equation and $n>0$ is even by the third, whose valuation
is $-n$.
If $e=v_p(s)>0$, then $T$ is a unit.  Negative $n$ is even by the
first equation, and $0<n<e$ is even by the second.  For $n>e$,
the second and third equations have valuations $e$ and $e-n$, so
both $e$ and $n$ are even.  Hence odd $n$ can occur only at
$n=e$ with $e$ odd.  In that case $e=v_p(d)>0$.
If instead $e=v_p(T)>0$, the first equation handles $n<0$ and
$0<n<e$, while the first and fourth have valuations $e,e-n$ for
$n>e$.  Again odd $n$ forces $n=e=v_p(d)>0$ odd.

This accounts for every odd prime, including equal-valuation cases.
Squareclasses in $\Q$ are determined by sign and valuation parity.
Finally $C^2=\ell-s$ implies $\ell>s>0$ when $t>1$.
\end{proof}

\begin{lemma}[Ramified classes]\label{lem:common-ramified}
Let $p$ be odd and suppose $v_p(d)>0$ is odd.  If
$\mathcal U_t(\Q_p)$ is nonempty, then $\leg2p=1$ and every point satisfies
\begin{equation}\label{eq:common-ramified-classes}
 [\ell]_{\Q_p}\in\{1,[d]_{\Q_p}\}.
\end{equation}
If $\leg2p=1$, both classes in \eqref{eq:common-ramified-classes}
occur.
\end{lemma}
\begin{proof}
First suppose $e=v_p(s)>0$ is odd.  If $v_p(\ell)>0$, then
$P^2\equiv2\pmod p$.  Otherwise
$B^2=T(1+s/\ell)\equiv2\pmod p$.  Thus $2$ is a square modulo $p$.
A negative $n=v_p(\ell)$ makes $M^2=-s(1-T/\ell)$ have odd valuation $e$.
For $0<n<e$, the valuations $n$ of $C^2$ and $e-n$ of $M^2$
cannot both be even.  For $n>e$, $C^2$ has odd valuation $e$.
Only $n=0$ and $n=e$ remain.

If $n=0$, the third equation forces $v_p(T-\ell)>0$, so
$\ell\equiv T\equiv2\pmod p$ and $[\ell]_{\Q_p}=1$.
If $n=e$, then
\[
 M^2=(d/\ell)(1-\ell/T),
\]
and the second factor is a principal unit, hence a square.
Thus $[\ell]_{\Q_p}=[d]_{\Q_p}$.
The substitution
$(t,P,C,M,B)\mapsto(-t,C,P,B,M)$ leaves $\ell,d$ unchanged and proves
the assertion when $v_p(T)>0$.

For existence near $t=1$, choose $\ell=s+p^{2h}$ with $2h>e$.
Then $C=p^h$, and the other three radicands have square unit residues
$2,2,4$.  Hensel's lemma gives all three roots.  This realizes $[d]$;
the reflection \eqref{eq:common-reflection} realizes $1$.
The sign substitution treats the other case.
\end{proof}

\begin{lemma}[The dyadic exponent]\label{lem:common-dyadic}
If $\mathcal U_t(\Q_2)$ is nonempty, then either $v_2(t)\ge2$ or
$v_2(t)$ is negative and odd.  If additionally $d$ is a square in
$\Q_2$, then
\[
 v_2(\ell)=v_2(t)<0\quad\text{is odd}.
\]
\end{lemma}
\begin{proof}
Apply \eqref{eq:common-return}, without real sign restrictions, to obtain
coordinates $x,y$ and four distances $r_j$ in $\Q_2$.
Choose $N=2^h$ with $h\ge0$ minimal such that
$X=Nx,Y=Ny\in\Z_2$.  Then $A_j=Nr_j\in\Z_2$, since their squares
are integral.  If $h=0$, one of the four distances has square
$2\pmod4$, a contradiction.  Hence $h\ge1$, and $X,Y$ are not both
even.  They cannot both be odd, since $A_0^2=X^2+Y^2$.
After exchanging the coordinate axes, take $X$ odd and $Y$ even.
Every $A_j$ is odd.  The first and third squared distances modulo eight
now give $4\mid Y$ and $4\mid Y-N$, so $h\ge2$.

Put $U=N-2X$, $V=N-2Y$, and
$S_\pm=A_0A_3\pm A_1A_2$.  Direct expansion gives
\[
 S_+S_-=-N^2UV,\qquad t=S_+/N^2,\qquad v_2(U)=1.
\]
For $V\ne0$, we have $v_2(V)\ge2$.  Exactly one of the sum and
difference of two odd units has valuation one.  The product identity
therefore gives either $v_2(t)=1-2h$ or $v_2(t)=v_2(V)\ge2$.
If $V=0$, then $S_+\ne0$ because $t\ne0$, so $S_-=0$ and
$S_+=2A_0A_3$; again $v_2(t)=1-2h$.

If $d$ is a dyadic square, the alternative $v_2(t)\ge2$ is impossible:
it would give $d\equiv-1\pmod8$.  Write $m=v_2(t)<0$, which is odd,
and $n=v_2(\ell)$.  For $n>m$, $P^2=\ell+T$ has odd valuation $m$.
For $n<m$, $B^2=T(1+s/\ell)$ has odd valuation $m$.
Thus $n=m$.
\end{proof}

\subsection{The residue matrix}

Let $\Delta=p_1\cdots p_r>1$ be positive, odd and squarefree.  Define
$A_\Delta=(a_{ij})\in M_r(\F_2)$ by
\begin{equation}\label{eq:common-matrix}
 (-1)^{a_{ij}}=\leg{p_j}{p_i}\quad(i\ne j),\qquad
 a_{ii}=\sum_{j\ne i}a_{ij}.
\end{equation}
Thus every row sums to zero.  Quadratic reciprocity does not in general
make this matrix symmetric; the denominator $p_i$ in
\eqref{eq:common-matrix} specifies the place at which the row is imposed.

\begin{proposition}\label{prop:common-matrix}
Suppose
\begin{equation}\label{eq:common-delta-hyp}
 \Delta\equiv1\pmod8,\qquad p_i\equiv1\text{ or }7\pmod8,
 \qquad t>1,\qquad d\in\Delta\Q^{\times2}.
\end{equation}
Every rational point of $\mathcal U_t$ has
\begin{equation}\label{eq:common-kernel}
 [\ell]=\left[2\prod_{i=1}^r p_i^{e_i}\right],
 \qquad (e_1,\ldots,e_r)^\mathsf T\in\ker A_\Delta.
\end{equation}
If $\rank A_\Delta=r-1$, then
\begin{equation}\label{eq:common-two-classes}
 [\ell]\in\{[2],[2\Delta]\}.
\end{equation}
\end{proposition}
\begin{proof}
Write $d=\Delta q^2$.  An odd valuation of $d$ cannot be negative,
since $v_p(t)<0$ implies $v_p(t^2-1)=2v_p(t)$.
Thus $S(d)=\{p_1,\ldots,p_r\}$.
The support lemma leaves $[\ell]=[2^e\prod p_i^{e_i}]$.
As $\Delta\equiv1\pmod8$, $d$ is a dyadic square; the dyadic lemma
forces $e=1$.

At $p_i$, if $e_i=0$, the ramified-class lemma requires $\ell$ to be a
local square; if $e_i=1$, it requires $\ell/d$ to be a local square.
All valuation parities have already been accounted for.
Moreover $\leg2{p_i}=1$.  The remaining unit conditions are exactly
\[
 \sum_{j\ne i}a_{ij}(e_j+e_i)=0,
\]
which is row $i$ of \eqref{eq:common-kernel}.
The zero and all-one vectors belong to the kernel.  If its dimension
is one, they are its only vectors, proving \eqref{eq:common-two-classes}.
\end{proof}

\begin{lemma}[A local squareclass restriction at $37$]\label{lem:common-37}
Suppose $t\in\Z_{37}$ and
$\overline t\in\{14,-14,16,-16\}$.  Every point of
$\mathcal U_t(\Q_{37})$ has $[\ell]_{\Q_{37}}=1$.
\end{lemma}
\begin{proof}
The substitution $t\mapsto-t$ used above preserves $\ell$, so it is enough
to treat $\overline t=14,16$.  In either case $T$ is a nonsquare unit.
A positive valuation of $\ell$ would give $P^2\equiv T\pmod{37}$;
a negative valuation would give $B^2\equiv T\pmod{37}$.
Thus $\ell$ is a unit.  We prove that its reduction $\lambda$ is square.

Work in $\F_{37}$.  The elements $13,15,17$ are twice squares:
\[
 13=2\cdot5^2,\qquad15=2\cdot10^2,\qquad17=2\cdot8^2.
\]
The element $2$ is nonsquare, while $-1$ is square.
If $\lambda$ were nonsquare, the last two radicands, together with the
first two, would imply that all four elements
\[
 \lambda+\overline t+1,\quad\lambda+\overline t-1,
 \quad\lambda-\overline t+1,\quad\lambda-\overline t-1
\]
are squares, allowing zero.

Here is a polynomial test for two squares differing by $2$.
If $a^2=z+1$ and $b^2=z-1$, put $w=(a+b)^2/2$.
Since $(a+b)(a-b)=2$, we have $w\ne0$ and
\[
 z=(w+w^{-1})/2,\qquad w^{18}=-1.
\]
Let $T_j$ be the Chebyshev polynomials, characterized by
$T_j((w+w^{-1})/2)=(w^j+w^{-j})/2$.
Then $T_9(z)=0$.  Reduction of $-T_9/3$ modulo $37$ gives
\[
 F(Z)=Z^9+7Z^7+4Z^5+3Z^3-3Z.
\]
Apply this to $z=\lambda+\overline t$ and
$z=\lambda-\overline t$.  Both $F$-values vanish.  Expanding and
factoring their sum gives the identities in $\F_{37}[X]$
\begin{align*}
 F(X+14)+F(X-14)&=2X^3(X^2-14)(X^2+4)(X^2+6),\\
 F(X+16)+F(X-16)&=2X^3(X^2-14)(X^2+6)(X^2+18).
\end{align*}
The three elements $14,31=-6,19=-18$ are nonsquares: each has
$18$th power $-1$.  Since $\lambda\ne0$, the second identity is
impossible.  The first leaves only $\lambda^2=-4$, hence
$\lambda=12$ or $25$, both squares ($12=7^2$ and $25=5^2$).
This too contradicts the assumed nonsquareness.
Thus $\lambda$ is square; the square criterion for odd local fields
then gives $[\ell]_{\Q_{37}}=1$.
\end{proof}

\begin{theorem}[Residue-matrix exclusion]\label{thm:common-exclusion}\label{thm:common-squareclass}
Let $\Delta=p_1\cdots p_r>1$ be positive, odd and squarefree, and suppose
\[
 \Delta\equiv1\pmod8,\qquad
 p_i\equiv1\text{ or }7\pmod8,\qquad
 \rank_{\F_2}A_\Delta=r-1.
\]
For every rational $t>1$ satisfying
\begin{equation}\label{eq:common-excluded-t}
 t^2-1\in\Delta\Q^{\times2},\qquad t\in\Z_{37},\qquad
 \overline t\in\{14,-14,16,-16\},
\end{equation}
one has $D_t(\Q)=\varnothing$.
\end{theorem}
\begin{proof}
A rational point in the finite chart would satisfy
$[\ell]\in\{[2],[2\Delta]\}$ by Proposition~\ref{prop:common-matrix}.
The residues of $d$ at the two positive values of $\overline t$ are
$10=11^2$ and $33=12^2$.  The prime $37$ does not divide $\Delta$,
since its residue modulo eight is $5$.  From $d=\Delta q^2$ it follows
that $q$ is a $37$-adic unit and $\Delta$ is a square unit there.
Thus both $2$ and $2\Delta$ are nonsquares in $\Q_{37}$.
This contradicts Lemma~\ref{lem:common-37}.  Proposition~
\ref{prop:common-genus-boundary} excludes every rational point outside
the finite chart.
\end{proof}

The theorem constrains one global squareclass.  A different choice of
$\ell$ at each completion cannot replace the compatibility in
\eqref{eq:common-kernel}.  In particular, the theorem applies without a
bound on the height of $\ell$, on the height of $t$, or on the number of
prime factors of $\Delta$.

\subsection{An explicit infinite family}

Consider
\[
 \Delta_0=1926281441=7\cdot17\cdot23\cdot31\cdot73\cdot311.
\]
In the indicated prime order, the matrix \eqref{eq:common-matrix} is
\begin{equation}\label{eq:common-six-matrix}
 A_{\Delta_0}=\begin{pmatrix}
 0&1&0&1&1&1\\
 1&1&1&1&1&1\\
 1&1&0&0&0&0\\
 0&1&1&1&1&0\\
 1&1&0&1&1&0\\
 0&1&1&1&0&1
 \end{pmatrix}.
\end{equation}
Its entries follow from quadratic reciprocity and reduction of the six
primes modulo one another.  Its kernel has a short explicit
certificate.  Row three gives $e_1=e_2$; row five then gives
$e_4=e_5$; row four gives $e_2=e_3$; row six gives $e_4=e_6$;
and row one gives $e_2=e_4$.  Hence all six coordinates agree.
The rank is therefore five.

\begin{corollary}\label{cor:common-infinite}
Put $C_0=269297$ and $K=2^{20}\cdot7\cdot37$.  For every integer
$j\ge1$, define
\begin{equation}\label{eq:common-infinite}
 q_j=\frac1{C_0}+\frac K{1+K^2j},\qquad
 t_j=\frac{1+\Delta_0q_j^2}{1-\Delta_0q_j^2}.
\end{equation}
The parameters $t_j$ are distinct and satisfy
\[
 1<t_j<\frac{17}{16}<\frac{\sqrt5}{2},\qquad
 D_{t_j}(\Q)=\varnothing.
\]
Every $D_{t_j}$ has a real point in the strict return chamber
\eqref{eq:common-chamber}.  In addition, $t_j+1$ is neither a rational
square nor twice a rational square, and
$[t_j^2-1]=[\Delta_0]\notin\{1,[2]\}$.
\end{corollary}
\begin{proof}
For $q_0=1/C_0$, direct cancellation gives
\[
 t(q_0)=\frac{138225}{131072}=:t_0,\qquad
 1-\Delta_0q_0^2=\frac{2^{18}}{C_0},\qquad
 t_0\equiv14\pmod{37}.
\]
The perturbation $\delta_j=K/(1+K^2j)$ has $37$-adic valuation one.
Since the denominator of $t(q_0)$ is a $37$-adic unit,
$t_j\equiv14\pmod{37}$ for all $j$.  Also
\[
 t_j^2-1=\Delta_0
 \left(\frac{2q_j}{1-\Delta_0q_j^2}\right)^2.
\]
At the real place $0<\delta_j<1/K<10^{-7}$.  The integer inequality
\[
 33\Delta_0\left(\frac1{269297}+\frac1{10^7}\right)^2
 =\frac{24893364230896160241}{26929700000000000000}<1
\]
gives $0<\Delta_0q_j^2<1/33$, and hence $1<t_j<17/16$.
The sequence $q_j$ is strictly decreasing, and the function
$(1+\Delta_0q^2)/(1-\Delta_0q^2)$ is strictly increasing on
$0<q<1/\sqrt{\Delta_0}$.  Thus the $t_j$ are distinct.
Theorem~\ref{thm:common-exclusion} now applies to each one.

For the real assertion choose $\ell=\sqrt d$ and positive roots in
\eqref{eq:common-four-roots}.  Since $s<\sqrt{sT}<T$, all four
radicands are positive.  Furthermore
\[
 (PM)^2=(BC)^2=2\sqrt d.
\]
The inequality $t<\sqrt5/2$ gives $d<1/4$, so
$|PM|+|BC|=2\sqrt{2\sqrt d}<2$.
Finally $v_7(q_0)=-1$ and $v_7(\delta_j)=1$, whence
$v_7(q_j)=-1$ and $v_7(1-\Delta_0q_j^2)=-1$.
As $t_j+1=2/(1-\Delta_0q_j^2)$, we get $v_7(t_j+1)=1$,
which proves both square exclusions.
\end{proof}

The limit parameter itself also satisfies the exclusion theorem:
\begin{equation}\label{eq:common-t0}
 t_0=\frac{138225}{131072},\qquad
 t_0^2-1=\frac{\Delta_0}{2^{34}},\qquad D_{t_0}(\Q)=\varnothing.
\end{equation}
For a one-prime squareclass $\Delta=p\equiv1\pmod8$, the matrix is
$(0)$, so its rank hypothesis is automatic.  Thus
Theorem~\ref{thm:common-exclusion} also applies to every such $p$ and
every $t$ satisfying \eqref{eq:common-excluded-t}.

\subsection{The unramified genus-seventeen covers}

The support lemma also gives a uniform finite descent for each fixed
parameter.  For $t>1$ define
\[
 \mathcal K_t=\left\{2^e\prod_{p\in S(d)}p^{e_p}:
                         e,e_p\in\{0,1\}\right\}.
\]
For $\kappa\in\mathcal K_t$, let $\widetilde D_{t,\kappa}$ be the
smooth projective curve obtained by adjoining to
\eqref{eq:common-four-roots} the further equation
\begin{equation}\label{eq:common-genus17}
 Z^2=\ell/\kappa.
\end{equation}

\begin{proposition}\label{prop:common-genus17}
The morphism $\widetilde D_{t,\kappa}\to D_t$ is a geometrically
connected unramified double cover, and
$g(\widetilde D_{t,\kappa})=17$.  Every rational point of $D_t$ lifts
to one of these covers; on the affine chart the lift retains the same
coordinates $\ell,P,C,M,B$.
\end{proposition}
\begin{proof}
On $D_t$, the function $\ell$ has eight zeros and eight poles, all of order
two, since the map to the $\ell$-line has ramification index two over
$0$ and $\infty$.  Thus adjoining its square root is unramified.
To prove geometric connectedness, suppose that $\ell$ became a square in
the four-radical function field over $\overline\Q$.  In a
multiquadratic extension the base-field squareclasses that become
squares are exactly the span of its defining classes.
The valuations at $-T,s,T,-s$ force each coefficient of such an
expression for $[\ell]$ to vanish.  The remaining class $[\ell]$ is nonzero
on $\overline\Q(\ell)$, a contradiction.  Riemann--Hurwitz now gives
\[
 2g(\widetilde D_{t,\kappa})-2=2(2\cdot9-2)=32.
\]
All rational points of $D_t$ are in the affine chart by
Proposition~\ref{prop:common-genus-boundary}.  For such a point,
Lemma~\ref{lem:common-support} supplies $\kappa\in\mathcal K_t$ with
$[\ell]=[\kappa]$, giving a rational $Z$ in
\eqref{eq:common-genus17}.
\end{proof}

\begin{proposition}[Finite local compatibility]\label{prop:common-finite-places}
For fixed $t>1$, define
\[
 \mathfrak S_\ell(t)=\{\kappa\in\mathcal K_t:
    \widetilde D_{t,\kappa}(\Q_v)\ne\varnothing
    \text{ for every place }v\}.
\]
It suffices to impose these conditions at the real place, at the primes
$p\le1153$, and at the primes at which one of $t,t-1,t+1$ has nonzero
valuation.  If $\mathfrak S_\ell(t)$ is empty, then $D_t(\Q)$ is empty.
\end{proposition}
\begin{proof}
Let $p$ be an odd prime outside the stated set.  The parameters
$t,s,T,\kappa$ are units.  The six branch sections on the $\ell$-line
remain distinct modulo $p$: their differences are products of units
among $2,t,s,T$.  The normalization of the five-radical cover is
smooth over $\Z_p$.  More explicitly, away from the branch sections
it is finite \emph{\'etale}; near a branch section, after an
\'etale extension extracting square roots of unit factors, its local
models are $w^2=z$, with $z$ a parameter on the base.  These models
are smooth, and smoothness descends.  The private-valuation argument
in the preceding proposition persists on the geometric special fibre,
so the latter is connected of genus seventeen.

The Hasse--Weil bound gives
\[
 \#\widetilde D_{t,\kappa}(\F_p)\ge p+1-34\sqrt p>0
 \qquad(p\ge1163).
\]
The next prime after $1153$ is $1163$.  Smooth Hensel lifting therefore
supplies a $\Q_p$-point.  A smooth curve over a local field has
infinitely many points in a sufficiently small neighbourhood of any
rational point, so the point can also be chosen off the finite
boundary.  The final assertion follows by applying
Proposition~\ref{prop:common-genus17} to a rational point of $D_t$.
\end{proof}

Thus the common-coordinate descent produces finitely many
geometrically connected covers while retaining all four ordered roots.
The residue-matrix theorem supplies domains on which the associated
compatibility set is empty, and Corollary~\ref{cor:common-infinite}
gives an explicit infinite sequence of such domains inside the strict
real return range.

\section{Nine elliptic factors and rank constraints}\label{sec:nine-factors}

The four roots in \eqref{eq:common-four-roots} give nine elliptic
quotients.  Their rational twists can be identified explicitly, and the
images of a single strict point impose rank conditions on every factor.
Throughout this section put
\[
 s=t-1,\qquad T=t+1,\qquad d=sT,\qquad z=\ell+T.
\]
Thus the coordinate denoted by $\ell$ satisfies
\begin{equation}\label{eq:nine-roots}
 P^2=z,\qquad C^2=z-2t,\qquad
 M^2=-\frac{s(z-2T)}{z-T},\qquad
 B^2=\frac{T(z-2)}{z-T}.
\end{equation}
The smooth projective curve with this function field is $D_t$.
For a strict rational-distance point, its four roots satisfy
\[
 t>1,\qquad s<\ell<T,\qquad P,B>0,\qquad CM\ne0.
\]
Here $CM\ne0$ also follows directly from the square: $C=0$ would give
$x=y$, and $M=0$ would give $x+y=1$.  Either diagonal forces a positive
rational distance to be a nonzero rational multiple of $\sqrt2$.

\subsection{Character quotients and their rational models}

Define the elliptic curves
\begin{equation}\label{eq:nine-elliptic-models}
\begin{aligned}
 E_d &: Y^2=X(X-d)(X+d),\\
 G_t &: Y^2=X(X-d^2)(X+d),\\
 F_+ &: Y^2=X(X-2d)(X+2td),\\
 F_- &: Y^2=X(X-2d)(X-2td),\\
 H_+ &: Y^2=X(X-2T)(X-T^2),\\
 H_- &: Y^2=X(X-2s)(X-d).
\end{aligned}
\end{equation}
Each curve carries the point at infinity as its origin.  All six cubics
are nonsingular for $t\notin\{0,1,-1\}$.

\begin{theorem}\label{thm:nine-isogeny}
For every $t\in\Q\setminus\{0,1,-1\}$ there is an isogeny over $\Q$
\begin{equation}\label{eq:nine-isogeny}
 \Jac(D_t)\sim_{\Q}
 E_d^2\times G_t\times F_+^2\times F_-^2\times H_+\times H_-.
\end{equation}
The rank of the group of rational N\'eron--Severi classes of $\Jac(D_t)$
is at least twelve.
\end{theorem}
\begin{proof}
The sign changes of the four roots give a group
$V=(\Z/2\Z)^4$, in the order $(P,C,M,B)$.  On the $z$-line its six
inertia generators, over $0,2t,\infty,2T,2,T$, are
\[
 e_1,\quad e_2,\quad e_1+e_2,\quad
 e_3,\quad e_4,\quad e_3+e_4.
\]
A nonzero character supported on just one of the blocks
$\langle e_1,e_2\rangle$ and $\langle e_3,e_4\rangle$ is ramified at
two points.  Its double-cover quotient is rational.  There are six such
characters.  Each of the nine characters nonzero on both blocks is
ramified at four points and has a genus-one quotient.

The latter characters are the products of one member of
$(P,C,PC)$ and one member of $(M,B,MB)$.  Their elliptic models are
\begin{equation}\label{eq:nine-character-table}
\begin{array}{c|ccc}
 &M&B&MB\\\hline
 P&E_d&H_+&F_+\\
 C&H_-&E_d&F_-\\
 PC&F_-&F_+&G_t.
\end{array}
\end{equation}
Here are birational maps establishing all the models and their rational
twists.  The character in the first column specifies the genus-one
function field $\Q(z,\text{character})$:
\begin{equation}\label{eq:nine-map-table}
\begin{array}{c|c|c|c}
\text{character}&\text{target}&X&Y\\\hline
 PM&E_d&-s\ell&sPM\ell\\[1mm]
 BC&E_d&T\ell&TBC\ell\\[1mm]
 PB&H_+&Tz&TPB\ell\\[1mm]
 CM&H_-&s(2T-z)&sCM\ell\\[1mm]
 PCB&F_+&4tT\ell/z&4tT^2\ell PCB/z^2\\[1mm]
 CMB&F_-&4td/(T-\ell)&-4td\ell CMB/(T-\ell)^2\\[1mm]
 PCMB&G_t&dT(2t/z-1)&dtT\ell PCMB/z^2.
\end{array}
\end{equation}
For example, the $PB$ and $PCMB$ equations before transformation are
\[
 (PB\ell)^2=Tz(z-2)(z-T),\qquad
 (PCMB\ell)^2=-dz(z-2t)(z-2)(z-2T).
\]
Substitution of the corresponding rows gives the stated cubics.  The
other rows follow by multiplying their radicands in
\eqref{eq:nine-roots}.  Every displayed abscissa is an invertible
M\"obius function of $z$, and its ordinate is a nonzero rational function
of $z$ times the indicated character.  It therefore recovers that
character and defines a birational map of genus-one curves.  All maps
extend uniquely between the smooth projective models.  The inverse
image of the point at infinity is a rational point on the smooth
quotient; this point supplies the origin used for its Jacobian.
More explicitly, the origins in the seven rows lie above
$z=\infty,\infty,\infty,\infty,0,2T,0$, respectively.
Each is the unique ramified point of its genus-one double cover above
the indicated value.

The involution
\begin{equation}\label{eq:nine-iota}
 (\ell,P,C,M,B)\longmapsto(d/\ell,B,M,C,P)
\end{equation}
exchanges the $PCB$ and $PMB$ characters, and exchanges $CMB$ and
$PCM$.  Composing the two corresponding rows of
\eqref{eq:nine-map-table} with this involution gives
\[
\begin{aligned}
 PMB\longrightarrow F_+:\quad (X,Y)
 &=\left(\frac{4td}{\ell+s},
           \frac{4td\ell PMB}{(\ell+s)^2}\right),\\
 PCM\longrightarrow F_-:\quad (X,Y)
 &=\left(\frac{4ts\ell}{\ell-s},
          -\frac{4ts^2\ell PCM}{(\ell-s)^2}\right).
\end{aligned}
\]
Their origins lie above $z=2$ and $z=2t$, respectively.  This gives
the two remaining entries of \eqref{eq:nine-character-table}.

The pullback of the differential of each elliptic quotient lies in its
one-dimensional character eigenspace in $H^0(D_t,\Omega^1)$.
These nine distinct eigenspaces span that space, since $g(D_t)=9$.
The product norm map from $\Jac(D_t)$ to the nine quotient
Jacobians has the direct sum of these differential pullbacks as its
cotangent map.  It therefore has invertible differential and is an
isogeny over $\Q$.  Its dual gives the product pullback isogeny and
proves \eqref{eq:nine-isogeny}.

For an elliptic square $E^2$, the two factor divisors and the diagonal
have intersection matrix
\[
 \begin{pmatrix}0&1&1\\1&0&1\\1&1&0\end{pmatrix},
 \qquad\det=2.
\]
They give three independent rational divisor classes.  Apply this to
$E_d^2,F_+^2,F_-^2$, and take the polarization class on each of
$G_t,H_+,H_-$.  The twelve classes are independent on the product.
Pullback and the norm along a rational isogeny identify N\'eron--Severi
spaces after tensoring with $\Q$, proving the assertion.
\end{proof}

The genus-three quotient and its three Pryms also follow from the same
character decomposition.  Let $K\subset V$ be generated by the sign
changes of $(P,M)$ and of $(C,B)$.  None of the six inertia subgroups
meets $K$ nontrivially.  Thus $D_t\to D_t/K$ is an unramified cover
of degree four and its quotient has genus three.  Its invariant field
is generated by
\[
 u=\frac{PM-BC}{2},\qquad v=\frac{PM+BC}{2},
 \qquad z=\frac{u^2+v^2+2}{2}+t-\frac{uv}{t}.
\]
Consequently it is the smooth model of
\begin{equation}\label{eq:nine-quartic}
 C_t:\quad t^2(u^2+v^2)^2-4u^2v^2-4t^2d=0.
\end{equation}
The characters trivial on $K$ give
\[
 \Jac(C_t)\sim_{\Q}E_d^2\times G_t.
\]
The three intermediate double covers of $C_t$ have genus five.
Writing $\kappa_1=e_1+e_3$ and $\kappa_2=e_2+e_4$, their Pryms are
\[
\begin{aligned}
 \operatorname{Prym}(D_t/\langle\kappa_1\rangle\to C_t)
   &\sim_{\Q}F_+\times F_-,\\
 \operatorname{Prym}(D_t/\langle\kappa_2\rangle\to C_t)
   &\sim_{\Q}F_+\times F_-,\\
 \operatorname{Prym}(D_t/\langle\kappa_1+\kappa_2\rangle\to C_t)
   &\sim_{\Q}H_+\times H_-.
\end{aligned}
\]
The additional characters in these three quotients are respectively
$(PMB,PCM)$, $(PCB,CMB)$, and $(PB,CM)$.  Removing the three
characters of $C_t$ from each quotient proves the displayed isogenies.

\subsection{Two-descent and four auxiliary point sets}

We recall the elementary forms of two-descent used below; see also
\cite{cassels1991,silverman2009}.  On a split cubic
\[
 E:\quad y^2=(x-e_1)(x-e_2)(x-e_3)
\]
with distinct rational roots, the Kummer homomorphism has coordinates
$[x-e_i]$ in $\Q^\times/\Q^{\times2}$.  At $(e_i,0)$ its $i$th
coordinate is $[(e_i-e_j)(e_i-e_k)]$; at another root $(e_j,0)$ it is
$[e_j-e_i]$, and at infinity it is $1$.  The product of the three
coordinates is $1$.  The map has kernel $2E(\Q)$.
Indeed, the tangent-and-chord law proves multiplicativity, and when
all three coordinates are squares one may choose rational roots
$r_i^2=x-e_i$ with $r_1r_2r_3=-y$; the usual halving formulas give a
point with abscissa
\[
 x+r_1r_2+r_1r_3+r_2r_3.
\]
The same calculation at a root shows that $(e_i,0)$ has a rational half
if and only if both $e_i-e_j$ and $e_i-e_k$ are squares.

For the integral model $y^2=x^3+ax^2+bx$, the one-coordinate
isogeny descent sends an ordinary point to $[x]$, sends $(0,0)$ to
$[b]$, and sends infinity to $1$.  Its kernel is the image of the dual
isogeny from
\[
 E':\quad y^2=x^3-2ax^2+(a^2-4b)x.
\]
Its image consists of signed squarefree divisors $q$ of $b$ for which
\begin{equation}\label{eq:nine-isogeny-quartic}
 W^2=qU^4+aU^2V^2+(b/q)V^4,\qquad\gcd(U,V)=1,
\end{equation}
is soluble in integers, together with the stated values at zero and
infinity.  To obtain this equation write $x=q(U/V)^2$ and substitute
in the cubic.  A prime not dividing $b$ cannot occur to an odd exponent
in $x$: negative valuations follow from the leading cubic term, and
positive valuations from the constant term of $x^2+ax+b$.
These observations give the support assertion and the homogeneous
quartic.

If $E$ has full rational two-torsion, write $\alpha(E)$ and
$\alpha'(E')$ for these two images.  Then
\begin{equation}\label{eq:nine-rank-descent}
 2^{\rank E(\Q)}=\frac{|\alpha(E)|\,|\alpha'(E')|}{4}.
\end{equation}
For completeness, the dual compositions are multiplication by two.
The other rational two-torsion points of $E$ map onto the kernel of
the dual isogeny.  The index calculation therefore gives
$[E(\Q):2E(\Q)]=|\alpha(E)|\,|\alpha'(E')|$.
Mordell--Weil and $E(\Q)[2]\simeq(\Z/2\Z)^2$ identify the left side
with $2^{\rank E(\Q)+2}$, giving
\eqref{eq:nine-rank-descent}.

\begin{lemma}\label{lem:nine-small-curves}
The following are complete rational point sets:
\begin{align*}
 E_1(\Q)&=\{O,(0,0),(1,0),(-1,0)\},\\
 E_2(\Q)&=\{O,(0,0),(2,0),(-2,0)\},\\
 A(\Q)&=\{O,(0,0),(6,0),(8,0)\},
   &&A:\ y^2=x(x-6)(x-8),\\
 E_0(\Q)&=\{O,(0,0),(-1,0),(-4,0),(2,\pm6),(-2,\pm2)\},
   &&E_0:\ y^2=x(x+1)(x+4).
\end{align*}
Here $E_1$ and $E_2$ are the curves $y^2=x^3-x$ and $y^2=x^3-4x$.
\end{lemma}
\begin{proof}
For $E_1$, the first descent image is $\{1,-1\}$.
Its dual $y^2=x^3+4x$ has nonnegative rational abscissae, so its image
is contained in $\{1,2\}$.  Both classes occur, the latter at $(2,4)$.
Equation~\eqref{eq:nine-rank-descent} gives rank zero.
For $E_2$, the first image is $\{1,-1,2,-2\}$, all represented by
two-torsion.  Its dual $y^2=x^3+16x$ has only the candidates $1,2$.
The class $2$ would require
\[
 W^2=2U^4+8V^4.
\]
If $U$ is odd the right side has valuation one at $2$; otherwise $V$
is odd and it has valuation three.  Thus the second image is
$\{1\}$ and the rank is zero.

For every positive rational $d$, the torsion of $E_d$ is its rational
two-torsion.  A rational half of the largest root $d$ would require
both $d$ and $2d$ square; the other roots have negative differences.
The third-division polynomial
\[
 3x^4-6d^2x^2-d^4
\]
cannot vanish at a rational number, since it would give
$x^2/d^2=1\pm2\sqrt3/3$.  Mazur's torsion theorem
\cite{mazur1977}, in the full-two-torsion case, leaves only
$(\Z/2\Z)^2$ after orders four and three have been excluded.
This completes the point lists for $E_1,E_2$.

For $A$, the model has $a=-14,b=48$.  There are no real points with
$x<0$, so the first descent image is contained in $\{1,2,3,6\}$;
the four two-torsion points realize all four classes.  The dual is
$y^2=x^3+28x^2+4x$, with candidates $1,-1,2,-2$.
For $q=-1$, equation~\eqref{eq:nine-isogeny-quartic} is nonsquare
modulo three for every primitive pair.  For $q=2$ its right side is
\[
 2\bigl((U^2-V^2)^2+16U^2V^2\bigr);
\]
opposite parity gives valuation one at two, and both odd gives
valuation five.  For $q=-2$ the right side is
\[
 -2\bigl((U^2-V^2)^2-12U^2V^2\bigr),
\]
with valuations one and three, respectively.  Thus the dual image
is $\{1\}$, proving rank zero.

For $E_0$, the first descent image is $\{1,-1,2,-2\}$, as witnessed
by $O,(-1,0),(2,6),(-2,2)$.  Its dual is
$y^2=x(x-1)(x-9)$ and has only the candidates $1,3$.
The second class would give
\[
 W^2=3U^4-10U^2V^2+3V^4.
\]
If both variables are units at three the right side is $2$ modulo
three; if just one is a unit its valuation is one.  Thus the dual
image is $\{1\}$ and the rank is zero.

To determine the last two torsion groups, the following table gives
all cubic values at the good primes five and seven, in the order
$x=0,1,\ldots,p-1$:
\[
\begin{array}{c|c|c|c}
\text{curve}&p&f(0),\ldots,f(p-1)&\#E(\F_p)\\\hline
 A&5&0,0,3,0,2&4\\
 A&7&0,0,6,3,4,1,0&8\\
 E_0&5&0,0,1,4,0&8\\
 E_0&7&0,3,1,0,6,4,0&8.
\end{array}
\]
The nonzero squares are $1,4$ modulo five and $1,2,4$ modulo seven.
Each point count follows by adding one for infinity and counting the
two square roots of each nonzero square.  Prime-to-$p$ torsion injects
under good reduction.  Combining the two primes bounds the torsion
orders by four and eight, respectively: the five-primary part is
excluded using seven, and the seven-primary part using five.
The displayed rational points attain these bounds, completing both
point lists.
\end{proof}

\subsection{Torsion of the six factors}

We shall need a convenient parametrization of a rational flex.
\begin{lemma}\label{lem:nine-flex}
If $y^2=x(x-A)(x-B)$ is nonsingular and has a rational point of order
three, then there is $k\in\Q\setminus\{0,1,-1,-\tfrac12,-2\}$ such that
\begin{equation}\label{eq:nine-flex}
 \frac BA=\frac{1+2k}{(1-k)(1+k)^3}.
\end{equation}
There is also a nonzero rational $r$ with $A=r^2(1-k^2)$.
\end{lemma}
\begin{proof}
Write the tangent at the order-three point $(x,y)$ as
$h(X)=m(X-x)+y$.  Since it meets the cubic with multiplicity three,
\[
 X(X-A)(X-B)-h(X)^2=(X-x)^3.
\]
For each root $e\in\{0,A,B\}$, the rational number
$r_e=h(e)/(x-e)$ satisfies $r_e^2=x-e$.  These numbers are nonzero
and distinct.  The identity $h(e)=y-mr_e^2=r_e^3$ shows that they are
the three roots of $R^3+mR^2-y$.  Hence their pairwise products sum
to zero.  Put $r=r_0$ and $k=r_A/r$.  Then
$r_B/r=-k/(1+k)$, and
\[
 A=r^2(1-k^2),\qquad B=r^2\frac{1+2k}{(1+k)^2}.
\]
This gives \eqref{eq:nine-flex}.  The values $0,\pm1$ contradict
nonzero roots or $A\ne0$; $-1/2$ makes $B=0$, and $-2$ makes $A=B$.
\end{proof}

\begin{lemma}\label{lem:nine-three-torsion}
For rational $t\notin\{0,1,-1\}$, the curve $G_t$ has no rational
three-torsion.  For $t>1$ the following additional assertions hold:
\begin{enumerate}
\item if $T$ is a rational square, neither $H_+$ nor $F_-$ has
rational three-torsion;
\item if $T$ is twice a rational square, $F_+$ has no rational
three-torsion.
\end{enumerate}
\end{lemma}
\begin{proof}
Set
\[
 D=(1-k)(1+k)^3,\quad R=\frac{1+2k}{D},\quad
 N=2+4k-2k^3-k^4.
\]
Then
\begin{equation}\label{eq:nine-flex-shifts}
 R-1=\frac{k^3(k+2)}D,\qquad R+1=\frac ND.
\end{equation}
Suppose first that $H_+$ has rational three-torsion and $T$ is a
square.  Take $A=2T,B=T^2$ in Lemma~\ref{lem:nine-flex}.
It gives $T/2=R$ and
\[
 r^2=\frac{4(1+2k)}{(1-k)^2(1+k)^4}.
\]
Thus $1+2k=a^2$ and $1-k^2=2b^2$ for rational $a,b$, whence
$8b^2=3+2a^2-a^4$.  Write $a=U/V$ in lowest terms.  The rational
number $2bV^2$ is an integer, since its square multiplied by two is
integral.  Calling it $W$ gives
\[
 2W^2=3V^4+2U^2V^2-U^4.
\]
Modulo sixteen the right side is $4$ if $U,V$ are both odd, is
$7$ or $15$ if only $U$ is odd, and is $3$ or $11$ if only $V$
is odd.  None is $0,2,8$, the possible residues of $2W^2$.

For $F_-$ use $A=2d,B=2td$, so $t=R$.  If $T$ is square,
\eqref{eq:nine-flex-shifts} and $A=r^2(1-k^2)$ show that
\[
 \frac{r^2}T=
 \left(\frac{k}{(1-k)(1+k)^2}\right)^2 2k(k+2).
\]
Hence $2k(k+2)$ is square.  If $v_2(k)<0$ its valuation is
$1+2v_2(k)$, and if $v_2(k)=0$ its valuation is one.  Consequently
$v_2(k)>0$.  But then $D$ is a dyadic unit and $N\equiv2\pmod8$;
thus $v_2(T)=v_2(N/D)=1$, contradicting that $T$ is square.

For $F_+$ take $A=2d,B=-2td$, so $-t=R$.  If $T$ is twice square,
\eqref{eq:nine-flex-shifts} gives
\[
 -2k(k+2)(1-k^2)
 =2T\left(\frac{(1-k)(1+k)^2}{k}\right)^2\in\Q^{\times2}.
\]
On the quartic $w^2=-2k(k+2)(1-k^2)$, the transformation
\[
 X=4-4/k,\qquad Y=-4w/k^2
\]
gives the curve $A$ of Lemma~\ref{lem:nine-small-curves}.
The branch point $k=0$ maps to infinity.  The quartic has no rational
points at infinity, since its leading coefficient is $2$.
The complete point list on $A$ therefore leaves just
$k=0,-2,1,-1$, all excluded by Lemma~\ref{lem:nine-flex}.

Finally, rational three-torsion on $G_t$ would give
$A=d^2,B=-d$, and hence $R=-1/d$.  Since $d+1=t^2$,
\[
 t^2=\frac{R-1}R=\frac{k^3(k+2)}{1+2k}.
\]
Consequently
\[
 (X,Y)=\left(2k,\frac{2t(1+2k)}k\right)
\]
is a point on $E_0$.  Its complete list in
Lemma~\ref{lem:nine-small-curves} leaves
$k\in\{0,-\tfrac12,-2,1,-1\}$, again impossible.
\end{proof}

\subsection{Positive rank on every factor}

\begin{theorem}\label{thm:nine-positive-ranks}
Let a strict rational-distance point have product parameter
$t=r_0r_3+r_1r_2$, and let the curves be as in
\eqref{eq:nine-elliptic-models}.  Each of
$E_d,G_t,F_+,F_-,H_+,H_-$ has positive rational rank.  In particular,
\begin{equation}\label{eq:nine-rank-nine}
 \rank\Jac(D_t)(\Q)\ge9.
\end{equation}
\end{theorem}
\begin{proof}
Every ordinate in \eqref{eq:nine-map-table} is nonzero at the strict
point.  The point $(-s\ell,sPM\ell)$ on $E_d$ therefore has infinite
order, by the torsion argument in
Lemma~\ref{lem:nine-small-curves}.  If $d$ were square, $E_d$ would
be rationally isomorphic to $E_1$; if $2d$ were square it would be
isomorphic to $E_2$.  Both are rank zero.  Thus
\begin{equation}\label{eq:nine-d-exclusions}
 d\notin\Q^{\times2},\qquad 2d\notin\Q^{\times2}.
\end{equation}

We next exclude rational order four on the remaining curves.
A two-torsion point can have a rational half only at the largest
root, because any other root has a negative root difference.
For $F_+$ its half would require $2d$ square.  For $H_+$ and $H_-$
it would require $d$ square, after removing square factors.
For $G_t$ it would require $d^2+d=dt^2$ square.  These contradict
\eqref{eq:nine-d-exclusions}.

For $F_-$ the largest root is $2td$.  Its half requires both
$2td$ and $2d(t-1)$ square.  Their quotient gives
$t/(t-1)$ square, so $b^2=t(t-1)$ for a nonzero rational $b$;
the second condition also gives $a^2=2(t+1)$.
Then $(X,Y)=(a^2,2ab)$ lies on
\[
 Y^2=X(X-2)(X-4)
\]
with nonzero ordinate.  Translation $X\mapsto X-2$ identifies this
curve with $E_2$, contradicting its point list.

Mazur's theorem now implies that torsion on each of the five curves
is either $(\Z/2\Z)^2$ or $\Z/2\Z\times\Z/6\Z$.
Lemma~\ref{lem:nine-three-torsion} excludes the second case for $G_t$,
so its nonzero-ordinate image is of infinite order.
For the other four curves, odd torsion is divisible by two.  Hence its
Kummer image is trivial, and the Kummer image of all torsion is just
the image of the rational two-torsion.

The common roots yield the identities
\begin{equation}\label{eq:nine-point-kummer}
\begin{aligned}
 X_{F_+}-2d&=2(TC/P)^2,&
 X_{F_-}-2d&=2(sB/M)^2,\\
 X_{H_+}&=TP^2,&
 X_{H_-}-2s&=-sC^2.
\end{aligned}
\end{equation}
On $H_-$, the possible torsion classes of the coordinate $X-2s$
are $1$ and $[-2s]$.  The actual class $[-s]$ differs from the first
by sign and from the second by $[2]$.  This image has infinite order.
On $H_+$, the torsion classes of $X$ are $1,[2T]$.
The actual class $[T]$ differs from both unless $T$ is square.
In that case Lemma~\ref{lem:nine-three-torsion} excludes order three,
and the nonzero ordinate excludes two-torsion.

For $F_-$, the torsion classes of $X-2d$ are
\[
 1,\quad[-2d],\quad[-s],\quad[2T].
\]
The actual class $[2]$ can meet this list only when $T$ is square;
the same lemma handles that case.  For $F_+$ the list is
\[
 1,\quad[-2d],\quad[T],\quad[-2dT].
\]
The actual class $[2]$ meets it only when $T$ is twice square,
which is the remaining case of that lemma.  Thus all six ranks are
positive.  Their multiplicities in \eqref{eq:nine-isogeny} give
$2+1+2+2+1+1=9$.
\end{proof}

\subsection{Two independent images and the rank strata}

Two quotient maps can furnish independent points on a single elliptic
curve.  The following argument uses the full Kummer rows of the points
and of their sum.  We write pairs of rational numbers for their classes
in $(\Q^\times/\Q^{\times2})^2$.

On $E_d$, take the two diagonal images
\[
 Q=(-s\ell,sPM\ell),\qquad R=(T\ell,TBC\ell).
\]
For the coordinates $(X,X-d)$ their rows are
\begin{equation}\label{eq:nine-cm-rows}
 \delta Q=(-s\ell,-s),\qquad
 \delta R=(T\ell,T),\qquad
 \delta(Q+R)=(-d,-d).
\end{equation}
For example $X(Q)-d=-sP^2$ and $X(R)-d=TC^2$.
The four torsion rows are
\begin{equation}\label{eq:nine-cm-torsion}
 (1,1),\quad(-1,-d),\quad(d,2),\quad(-d,-2d).
\end{equation}
Since $d>0$ is nonsquare, the sum row in
\eqref{eq:nine-cm-rows} is absent from this list.
Comparing the individual rows, with signs retained, shows that either
one can meet a torsion row only if $T$ is square or twice square.
Thus
\begin{equation}\label{eq:nine-cm-two}
 T\notin\Q^{\times2}\cup2\Q^{\times2}
 \quad\Longrightarrow\quad \rank E_d(\Q)\ge2.
\end{equation}
Indeed the images of $Q,R$ in
$E_d(\Q)/(2E_d(\Q)+E_d(\Q)_{\mathrm{tors}})$ are then nonzero and
distinct, hence linearly independent over $\F_2$.

For $F_+$ and $F_-$, take the image in
\eqref{eq:nine-map-table} and its image after the involution
\eqref{eq:nine-iota}.  For both pairs the coordinates $(X,X-2d)$ give
\begin{equation}\label{eq:nine-f-rows}
 \delta Q_\pm=(tT\ell,2),\qquad
 \delta R_\pm=(ts\ell,2),\qquad
 \delta(Q_\pm+R_\pm)=(d,1).
\end{equation}
The first-coordinate products are, explicitly,
\[
 X(Q_+)X(R_+)=d\left(\frac{4tT}{PB}\right)^2,\qquad
 X(Q_-)X(R_-)=d\left(\frac{4ts}{MC}\right)^2.
\]
The second coordinates follow from \eqref{eq:nine-point-kummer}.
The complete torsion rows are
\begin{equation}\label{eq:nine-f-torsion}
\begin{array}{c|cccc}
 F_+&(1,1)&(-t,-2d)&(2d,T)&(-2td,-2dT)\\
 F_-&(1,1)&(t,-2d)&(2d,-s)&(2td,2T).
\end{array}
\end{equation}
For $F_+$, the sum row $(d,1)$ never occurs: the positive rows
would require $d$ or $2$ to be square.  An individual row in
\eqref{eq:nine-f-rows} can occur only if $T$ is twice square.
It follows that
\begin{equation}\label{eq:nine-fplus-two}
 T\notin2\Q^{\times2}\quad\Longrightarrow\quad\rank F_+(\Q)\ge2.
\end{equation}
For $F_-$, an individual row can meet torsion only when $T$ is
square.  The sum row can meet torsion only if both $2t$ and $2T$
are squares.  Therefore
\begin{equation}\label{eq:nine-fminus-two}
\begin{gathered}
 T\notin\Q^{\times2},\qquad
 \text{at least one of $2t,2T$ is nonsquare}
 \\
 \Longrightarrow\quad\rank F_-(\Q)\ge2.
\end{gathered}
\end{equation}

\begin{theorem}[Rank strata]\label{thm:nine-rank-strata}
For the product parameter of any strict rational-distance point, the
following lower bounds hold:
\begin{equation}\label{eq:nine-rank-table}
\begin{array}{l|ccc|c}
 \text{condition}&\rank E_d&\rank F_+&\rank F_-&\rank\Jac(D_t)\\\hline
 T\notin\Q^{\times2}\cup2\Q^{\times2}&2&2&2&15\\
 T\in\Q^{\times2}&1&2&1&11\\
 T\in2\Q^{\times2},\ 2t\notin\Q^{\times2}&1&1&2&11\\
 T\in2\Q^{\times2},\ 2t\in\Q^{\times2}&1&1&1&9.
\end{array}
\end{equation}
Every entry in the table is a lower bound.
\end{theorem}
\begin{proof}
Combine \eqref{eq:nine-cm-two}, \eqref{eq:nine-fplus-two}, and
\eqref{eq:nine-fminus-two} with
Theorem~\ref{thm:nine-positive-ranks}.  Formula~\eqref{eq:nine-isogeny}
gives the rank identity
\[
 2\rank E_d+\rank G_t+2\rank F_++2\rank F_-+
 \rank H_++\rank H_-.
\]
Substitution gives the final column.  The four rows exhaust the
positive rational values of $T$, since $2$ is nonsquare.
\end{proof}

The theorem gives a useful exclusion criterion for a fixed parameter:
a rank bound below the relevant entry in
\eqref{eq:nine-rank-table} excludes a strict point on that fibre.
It also places every fibre arising from a strict point outside the
classical Chabauty inequality $\rank\Jac(D_t)(\Q)<g(D_t)$.

\section{Descent with a common source}
\label{sec:common-source-descent}

The square conditions at different places constrain the same rational
points.  This section retains that dependence through finite quotients of
the Mordell--Weil group.  The first construction combines characters modulo
$4$ into an invariant of a complete coset modulo $2$.  A second application
uses a quotient modulo $8$ to exclude the parameter $n=15569$.

\subsection{The source coset and the formal parameter}

Write
\[
 E_n:\quad y^2=x(x-n)(x+n),\qquad
 T_0=(0,0),\quad T_+=(n,0),\quad T_-=(-n,0),
\]
where $n>1$ is odd and squarefree.  The ordered Kummer map is
\[
 \delta:E_n(\Q)/2E_n(\Q)\longrightarrow
 (\Q^*/\Q^{*2})^2,\qquad P\longmapsto([x(P)],[x(P)-n]).
\]
Its values at the two-torsion points are
\begin{equation}\label{eq:cs-torsion-rows}
\begin{array}{c|rrrr}
P&O&T_0&T_+&T_-\\ \hline
\delta(P)&(1,1)&(-1,-n)&(n,2)&(-n,-2n).
\end{array}
\end{equation}
Products and equalities in Kummer coordinates are understood modulo
rational squares.  To recall the kernel statement, let $e_1,e_2,e_3$
be the roots $0,n,-n$.  Intersecting a line with the cubic shows that,
for each $i$, the product of $x-e_i$ over the three intersection points
is a square.  This proves multiplicativity of each Kummer coordinate.
If all $x(P)-e_i$ are squares, choose $r_i^2=x(P)-e_i$ with
$r_1r_2r_3=-y(P)$, and put $s_1=r_1+r_2+r_3$,
$s_2=r_1r_2+r_1r_3+r_2r_3$.  The point
\[
 (x(P)+s_2,-y(P)-s_1s_2)
\]
doubles to $P$, as substitution and its tangent of slope $-s_1$ show.
The formulas extend at two-torsion by taking one $r_i=0$.
Thus the kernel is $2E_n(\Q)$.

For $j\geq2$, put
\[
 E_j(\Q_2)=\{O\}\cup
 \{P\in E_n(\Q_2):v_2(x(P))\leq-2j\},\qquad
 H=E_n(\Q)\cap E_2(\Q_2).
\]
Thus $E_j(\Q_2)$ is the indicated neighbourhood of the point at infinity.
For a nonzero $P\in H$, its \emph{depth} is
$j(P)=-v_2(x(P))/2$.  In the general criterion below we assume
\begin{equation}\label{eq:cs-splitting}
 E_n(\Q)=H\oplus E_n[2](\Q).
\end{equation}
This specifies the subgroup and every torsion translation used in the
argument.

Here are the formal-group estimates we shall need.  With $z=-x/y$ and
$-1/y=z^3W$, the equation of $E_n$ gives
\begin{equation}\label{eq:cs-formal-chart}
 W=1-n^2z^4W^2,\qquad
 W\in1+z^4\Z_2[[z^4]],\qquad x=z^{-2}W^{-1}.
\end{equation}
Moreover
\[
 \frac{dx}{2y}=
 \left(1+\frac{zW'}{2W}\right)dz,
 \qquad
 \ell(z)=\int_0^z\frac{dx}{2y}
       \in z+2z^5\Z_2[[z^4]].
\]
Indeed, the coefficients of $zW'/(2W)$ are even and its exponents are
positive multiples of $4$; integration divides by odd integers.  The
formal logarithm is additive and isometric on $4\Z_2$.  Consequently
$H$ is torsion-free, its nonzero depths are integers, and
\begin{equation}\label{eq:cs-local-squares}
 x(P),\quad x(P)-n,\quad x(P)+n\in\Q_2^{*2}\qquad(P\in H\setminus\{O\}).
\end{equation}
The last assertion follows from \eqref{eq:cs-formal-chart}, since all the
unit corrections are $1$ modulo $8$.

Choose $\nu\geq2$ with $H\subseteq E_\nu(\Q_2)$, and define
\begin{equation}\label{eq:cs-dyadic-character}
 \phi_0(P)=2^{-\nu}\ell(P)\pmod4,
 \qquad \eta_0(P)=\phi_0(P)\pmod2.
\end{equation}
These are homomorphisms.  Formula \eqref{eq:cs-formal-chart} also permits
$\ell(P)$ to be replaced by $z(P)$ in \eqref{eq:cs-dyadic-character}.

We consider rational points $Q=(a,b)$ and $R=(c,e)$ satisfying
\begin{equation}\label{eq:cs-source-pair}
 -n<a<0<c-n,\qquad
 \delta(Q)\delta(R)=(-n,-n),\qquad
 f(Q,R)=\frac{a+c}{2(c+n)}.
\end{equation}
In particular $f(Q,R)>0$.  The correspondence of
Theorem~\ref{thm:ell-exact-source} sends a rational-distance configuration
to such a pair with $f(Q,R)$ a nonzero rational square.

Suppose $n\equiv1$ or $7\pmod8$, and choose $W\in H$ with
\begin{equation}\label{eq:cs-source-row}
 \delta(W)=
 \begin{cases}(n,1),&n\equiv1\pmod8,\\
 (-n,-n),&n\equiv7\pmod8.
 \end{cases}
\end{equation}
Write $Q=U+\tau_Q$ and $R=V+\tau_R$ according to
\eqref{eq:cs-splitting}.  The two-torsion rows are distinct over $\Q_2$.
Equations \eqref{eq:cs-local-squares} and \eqref{eq:cs-source-pair} therefore
give $\tau_Q+\tau_R=T_0$ when $n\equiv1$, and $\tau_Q=\tau_R$ when
$n\equiv7$.  Applying \eqref{eq:cs-torsion-rows} over $\Q$ then gives
\begin{equation}\label{eq:cs-complete-coset}
 U+V\in W+2H.
\end{equation}
Thus the undetermined part of the source is one element of $2H$.

\subsection{Local conditions on the two remainders}

The following calculation extracts exactly the information needed from
the place $2$.  Put $u=x(U)$ and $v=x(V)$.  Translation by two-torsion gives
\begin{equation}\label{eq:cs-translations}
 x(P+T_0)=-\frac{n^2}{x(P)},\quad
 x(P+T_+)=n+\frac{2n^2}{x(P)-n},\quad
 x(P+T_-)=-n+\frac{2n^2}{x(P)+n}.
\end{equation}
For $n\equiv1$, the orders $(T_0,O)$ and $(O,T_0)$ have respectively
$v_2(f)=-1$ and $v_2(f)=v_2(u)-1$.  For $n\equiv7$, the orders
$(T_+,T_+)$ and $(T_-,T_-)$ have respectively $v_2(f)=-1$ and
$v_2(f)=v_2(v)-1$.  These four valuations are odd.  The other four orders
give
\begin{equation}\label{eq:cs-surviving-formulas}
\begin{array}{c|c|c}
n\pmod8&(\tau_Q,\tau_R)&f(Q,R)\\ \hline
1&(T_-,T_+)&n(u+v)/(2v(u+n))\\
1&(T_+,T_-)&(u+v)/(2(u-n))\\
7&(O,O)&(u+v)/(2(v+n))\\
7&(T_0,T_0)&-n(u+v)/(2u(v-n)).
\end{array}
\end{equation}
Each expression in the last column differs from $(u+v)/2$ by a
$\Q_2$-square.  If $U,V$ have different depths, $u+v$ has even valuation
and $(u+v)/2$ has odd valuation.  If their depths agree, write
$u=2^{-2j}r^2$, $v=2^{-2j}s^2$ with $r,s\in\Z_2^*$; then
$v_2(r^2+s^2)=1$.  We obtain
\begin{equation}\label{eq:cs-sign-match}
 v_2(f)\text{ even}\quad\Longrightarrow\quad
 j(U)=j(V),\qquad \phi_0(V)=\pm\phi_0(U).
\end{equation}
For the last assertion, at depth $\nu$ both character values are odd;
at depth $\nu+1$ both equal $2$; at larger depths both vanish.

We next specify the odd primes used in the criterion.  Let $q\equiv5\pmod8$
be a prime with $q\nmid n$ and $n\in\Q_q^{*2}$.  We allow either of the
following reduction conditions on the entire subgroup $H$:
\begin{enumerate}
\item $\operatorname{red}_q(H)\subseteq\{O,T_0\}$.  Write
 $\operatorname{red}_q(P)=\zeta_q(P)T_0$ and set
 $\phi_q=2\zeta_q:H\longrightarrow\Z/4\Z$.
\item There is $A_q\in E_n(\F_q)$ with $2A_q=T_0$ such that
 $\operatorname{red}_q(H)\subseteq\langle A_q\rangle$.  Define
 $\operatorname{red}_q(P)=\phi_q(P)A_q$ and put $i_q=x(A_q)/n$.
 Require that
 \begin{equation}\label{eq:cs-wrong-unit}
 \begin{cases}(1-i_q)/2,&n\equiv1\pmod8,\\
 (1+i_q)/2,&n\equiv7\pmod8
 \end{cases}
 \quad\text{is a nonsquare in }\F_q.
 \end{equation}
\end{enumerate}
In the second case $A_q+T_0=-A_q$, so
\eqref{eq:cs-translations} gives $i_q^2=-1$.  In either
case put $\eta_q=\phi_q\bmod2$.

\begin{lemma}\label{lem:cs-odd-row}
For such a prime, $\eta_q(W)=0$.  If a pair
\eqref{eq:cs-source-pair} has $v_2(f)$ even and $f\in\Q_q^{*2}$, then
\[
 \phi_q(U)=\phi_q(V)\in\{0,2\}.
\]
\end{lemma}
\begin{proof}
Only the cyclic case requires a proof that $\eta_q(W)=0$.  An odd
multiple of $A_q$ has first Kummer coordinate $ni_q$, a nonsquare since
$q\equiv5\pmod8$.  The first coordinate of $\delta(W)$ is $n$ or $-n$,
both local squares.  Thus $W$ has even reduction coordinate.  Equation
\eqref{eq:cs-complete-coset} now implies $\eta_q(U)=\eta_q(V)$.

If both parities are odd, their abscissae reduce to $ni_q$.
Translation by $T_0$ preserves this abscissa, whereas translation by
$T_+$ or $T_-$ changes it to $-ni_q$.  The four orders in
\eqref{eq:cs-surviving-formulas} give exactly the residues in
\eqref{eq:cs-wrong-unit}.  Their numerators and denominators are units,
so this case is excluded.

Both remainders therefore reduce to $O$ or $T_0$.  Suppose their
reductions differ.  When $n\equiv1\pmod8$, the corresponding $Q,R$
both reduce to $T_+$ or both to $T_-$.  In the first case $f$ has
residue $1/2$.  In the second case use the exact identity
\[
 (2e)^2f=2(a+c)c(c-n).
\]
The right side is a unit with residue $-8n^3$, of squareclass $2$.
The identity applies at every positive depth of approach to $T_-$,
because $e\ne0$.

When $n\equiv7\pmod8$, $Q,R$ reduce in opposite orders to $O,T_0$.
Near $O$ the abscissa has negative even valuation and square unit part;
near $T_0$ it has positive even valuation and square unit part.  These
facts follow by writing the cubic as $y^2=x^3(1-n^2/x^2)$ or as
$y^2=-n^2x(1-x^2/n^2)$, respectively.  The abscissa near $O$ dominates
$a+c$, and substitution in $f$ again gives squareclass $2$.  Since $2$
is a nonsquare at $q$, the reductions cannot differ.  Their equal labels
$O$ or $T_0$ give the asserted equal character values.
\end{proof}

\subsection{A secondary invariant of the complete source}

Take finitely many primes $q_1,\ldots,q_s$ satisfying the preceding
conditions, and abbreviate their characters to $\phi_j,\eta_j$.  Set
\[
 V_H=H/2H,\qquad
 M:V_H\longrightarrow\F_2^{s+1},\qquad
 M(\bar Z)=(\eta_0(Z),\ldots,\eta_s(Z)).
\]
For $\kappa\in\ker M$, choose a representative $W\in H$ and define
\begin{equation}\label{eq:cs-beta}
 b(W)=\left(\frac{\phi_0(W)}2,\ldots,
                  \frac{\phi_s(W)}2\right),\qquad
 \beta(\kappa)=b(W)\pmod{\operatorname{im}M}.
\end{equation}
Division by $2$ denotes the isomorphism $2(\Z/4\Z)\simeq\F_2$.
For any $Z\in H$,
\begin{equation}\label{eq:cs-compensation}
 b(W+2Z)=b(W)+M(\bar Z).
\end{equation}
Thus \eqref{eq:cs-beta} defines a linear map
$\beta:\ker M\longrightarrow\operatorname{coker}M$.

\begin{theorem}[Common-source character criterion]
\label{thm:cs-matrix}
Assume \eqref{eq:cs-splitting}, $n\equiv1$ or $7\pmod8$, and the
reduction hypotheses above.  Let $W$ satisfy
\eqref{eq:cs-source-row} and $\eta_0(W)=0$, and put $\kappa=\bar W$.
If $\beta(\kappa)\ne0$, no pair \eqref{eq:cs-source-pair} has both
$v_2(f)$ even and $f\in\Q_{q_j}^{*2}$ for every $j$.

Equivalently, an exclusion certificate is a vector
$\mathbf c=(c_0,\ldots,c_s)\in\F_2^{s+1}$ such that
\[
 \sum_{j=0}^s c_j\eta_j=0\text{ on }H,
 \qquad
 \frac12\sum_{j=0}^s c_j\phi_j(W)=1\quad\text{in }\F_2.
\]
\end{theorem}
\begin{proof}
Suppose a pair exists.  By \eqref{eq:cs-sign-match}, one of $K=U-V$
and $K=U+V$ satisfies $\phi_0(K)=0$.  By
Lemma~\ref{lem:cs-odd-row}, every other row is equal and even on $U,V$;
both their sum and their difference vanish modulo $4$.  Hence this
single choice of $K$ satisfies $\phi_j(K)=0$ for all $j$.
Both choices lie in $W+2H$, so $K=W+2Z$ for some $Z\in H$.
Dividing $0=\phi_j(W)+2\phi_j(Z)$ by $2$ gives
$b(W)=M(\bar Z)$, contrary to $\beta(\kappa)\ne0$.

Finally, a vector lies outside $\operatorname{im}M$ precisely when a
linear functional annihilating that image takes value $1$ on it.
Writing the functional as dot product with $\mathbf c$ gives the stated
certificate.  The parity relation makes
$\frac12\sum c_j\phi_j$ a homomorphism $H/2H\to\F_2$, so its value is
independent of the representative $W$.
\end{proof}

For example, if $s=1$ and $\eta_1=\eta_0$, the secondary character is
\[
 \Delta=\frac{\phi_1-\phi_0}{2}:H/2H\longrightarrow\F_2.
\]
On $\ker\eta_0$, both individual half-characters change by $\eta_0(Z)$
under $W\mapsto W+2Z$.  Their difference is unchanged.  This explains
why the same correction must be retained at both places.

\subsection{An exclusion using the common quotient modulo eight}

\begin{theorem}\label{thm:cs-15569}
Let $n=15569$.  There are no finite rational points $Q=(a,b)$,
$R=(c,e)$ satisfying \eqref{eq:cs-source-pair} for which $f(Q,R)$ is a
square in each of $\Q_2,\Q_3,\Q_{113}$.
\end{theorem}

\begin{proof}
We first determine the exact quotient of the rational group that is
needed.  The integer $n$ is prime and is $1$ modulo $8$.  One short primality
certificate is
\[
 n-1=112\cdot139,\quad 2^{112}\equiv6064,\quad
 2^{n-1}\equiv1\pmod n,\quad 803n-2062\cdot6063=1.
\]
Since $139$ is prime, every prime divisor $p$ of $n$ has
$139\mid p-1$; as $139>\sqrt n$, this proves primality.
For a rational point, both Kummer coordinates are supported on $-1,2,n$: at an odd
prime not dividing $n$, either all three cubic factors have the same
negative valuation or at most one has positive valuation, and the
equation forces the relevant valuations to be even.

Over $\Q_2$, choose $v^2=n$ and use $X=x/n$, $Y=y/(nv)$.
The curve becomes $Y^2=X(X-1)(X+1)$, with the same local Kummer
coordinates.  Its non-two-torsion points have the following rows; the
two-torsion centres have the rows in \eqref{eq:cs-torsion-rows}.
\begin{equation}\label{eq:cs-eight-local-rows}
\begin{array}{c|c}
\text{condition}&\delta\text{ in }(\Q_2^*/\Q_2^{*2})^2\\ \hline
v_2(X)\leq-4&(1,1)\\
v_2(X)=-2&(1,5)\\
v_2(X)=2&(-1,-5)\\
v_2(X)\geq4&(-1,-1)\\
X\equiv1\pmod8,\ v_2(X-1)=3&(1,10)\\
X\equiv1\pmod8,\ v_2(X-1)\geq5&(1,2)\\
X\equiv-1\pmod8,\ v_2(X+1)=3&(-1,-10)\\
X\equiv-1\pmod8,\ v_2(X+1)\geq5&(-1,-2).
\end{array}
\end{equation}
In the first four rows the nonzero valuation is even; in the last four
rows the displayed positive valuation is odd.  The table follows from
even valuation and unit part $1$ modulo $8$ in the cubic.  For example,
$X=u/4$ forces $u\equiv1\pmod8$, $X=4u$ forces
$u\equiv7\pmod8$, and near $\pm1$ the other two factors have valuations
$0,1$.  At depth $3$ the additional unit class is $5$; at greater odd
depth it is a square.  These cases exhaust the parity of the valuations
and the odd unit residues.

Since $n$ is a $\Q_2$-square, a globally supported Kummer coordinate
cannot have the extra local unit class $5$.  The four remaining rows
are precisely the distinct rows of $E_n[2]$.  Every rational point has
a unique two-torsion translate in $H$, proving
\eqref{eq:cs-splitting}.  Furthermore
\[
 \delta(H)\subseteq\{1,n\}^2,
 \qquad \ker(\delta|_H)=2H.
\]
Thus $\dim_{\F_2}H/2H\leq2$.

Two points $P,G\in H$ are specified by the following factorizations.
The three entries in each row are $x,x-n,x+n$:
\begingroup\setlength{\arraycolsep}{3pt}
\begin{equation}\label{eq:cs-rational-generators}
\begin{array}{c|c|ccc}
 &d&d^2x&d^2(x-n)&d^2(x+n)\\ \hline
P&7760&n\,7809^2&108983^2&n\,11009^2\\
G&13382920&n\,13522049^2&n\,1934751^2&2373852137^2.
\end{array}
\end{equation}
\endgroup
Choose the negative ordinate in each row:
\[
 y(P)=-\frac{n(7809)(108983)(11009)}{7760^3},\qquad
 y(G)=-\frac{n(13522049)(1934751)(2373852137)}{13382920^3}.
\]
The displayed differences are $n$, and their products are these
ordinates squared.  Hence the two rows specify rational points with
\[
 \delta(P)=(n,1),\qquad \delta(G)=(n,n).
\]
They form a basis of $H/2H$.  The depths are $j(P)=4$, $j(G)=3$.
The homomorphism $\ell/4\pmod2$ vanishes on these representatives and
therefore on $H$, proving $H\subseteq E_3(\Q_2)$.  Every element of $H$
has a representation
\begin{equation}\label{eq:cs-mod-eight}
 h=kG+mP+8h_1,\qquad 0\leq k,m<8.
\end{equation}
Indeed, successive substitution of $H=\{0,G,P,G+P\}+2H$ three times
gives this assertion.  Direct reduction of the fractions in
\eqref{eq:cs-rational-generators} gives
\begin{equation}\label{eq:cs-z-eight}
 z(G)/8\equiv7,\qquad z(P)/8\equiv6\pmod8.
\end{equation}
The formal logarithm shows that $z(h)/8\pmod8$ is a homomorphism on
$H/8H$.

Every nonzero $H$-point has positive first Kummer coordinate and hence
lies on the real identity component.  Equations
\eqref{eq:cs-source-pair} and \eqref{eq:cs-translations}, followed by
the dyadic valuation test, leave only
\begin{equation}\label{eq:cs-15569-branch}
 Q=U+T_-,\qquad R=V+T_+,\qquad U+V\in P+2H.
\end{equation}
Write
\[
 U=kG+mP+8h_1,\qquad V=k'G+m'P+8h_2.
\]
The source relation is
\begin{equation}\label{eq:cs-parities}
 k+k'\equiv0,\qquad m+m'\equiv1\pmod2.
\end{equation}

At $3$, $n\equiv2$, and $E_n(\F_3)$ consists of its four two-torsion
points.  Their local Kummer rows are
\[
 O:(1,1),\quad T_0:(2,1),\quad
 T_+:(2,2),\quad T_-:(1,2).
\]
The entries $1,2$ denote unit squareclasses in $\Q_3$.
Every root neighbourhood has an even positive offset valuation and
the indicated constant unit class; near $O$ the negative valuation is
even.  Since $-n$ is a local square, $Q,R$ have equal Kummer rows and
therefore the same reduction.  A common reduction $T_+$ gives
$f\equiv1/2$, a nonsquare.  For common reduction $T_-$ the identity
$(2e)^2f=2(a+c)c(c-n)$ has nonsquare unit residue, including every
positive pole depth.  Thus a square return requires their common
reduction to be $O$ or $T_0$.

The rational coordinates give $\operatorname{red}_3G=T_+$ and
$\operatorname{red}_3P=T_0$.  In view of
\eqref{eq:cs-15569-branch}, both $k$ and $k'$ must be odd.  It follows
from \eqref{eq:cs-z-eight} that $U,V$ have depth exactly $3$.  Put
\[
 \alpha=z(U)/8\equiv7k+6m,\qquad
 \gamma=z(V)/8\equiv7k'+6m'\pmod8.
\]
Both are odd.  Writing $u=2^{-6}r^2$ and $v=2^{-6}s^2$ in the formal
chart gives $r\equiv\alpha^{-1}$ and
$s\equiv\gamma^{-1}\pmod8$.  The square condition at $2$ is
\begin{equation}\label{eq:cs-eight-match}
 \frac{r^2+s^2}{2}\equiv1\pmod8
 \quad\Longleftrightarrow\quad
 \alpha\equiv\pm\gamma\pmod8.
\end{equation}
Here odd squares modulo $16$ are $1,9$, so their average is $1$ modulo
$8$ exactly when the two squares agree modulo $16$.

It remains to use the same coefficients at $113$.  At this prime
$n\equiv88\equiv38^2$, and the curve is $y^2=x^3+53x$.  Its complete
group structure follows from the short doubling certificate
\begin{equation}\label{eq:cs-f113-certificate}
\begin{array}{c|cccc}
 &A&2A&4A&8A\\ \hline
 &(23,39)&(81,112)&(77,79)&T_0\\[1mm]
 &B&2B&4B&8B\\ \hline
 &(2,1)&(7,6)&T_+&O.
\end{array}
\end{equation}
Thus $A$ has order $16$ and $B$ order $8$.  Their cyclic subgroups have
different points of order $2$, so intersect trivially.  They generate a
subgroup of order $128$.  Hasse's bound places the total order between
$93$ and $135$, whence
\[
 E_n(\F_{113})\simeq\Z/16\Z\oplus\Z/8\Z,
 \qquad \exp(2E_n(\F_{113}))=8.
\]
Both coordinates of $\delta(H)$ are squares in $\Q_{113}$, so every
element of $H$ is twice a local point.  Consequently
$\operatorname{red}_{113}(8H)=O$.

Set
\[
 g=\operatorname{red}_{113}G=(81,112),\qquad
 w=\operatorname{red}_{113}P=(36,58).
\]
The addition law gives
\[
 4g=T_0,\quad w=2g+T_+,\quad
 x(g)=81,\ x(3g)=16,\ x(g+T_+)=87,\ x(3g+T_+)=11.
\]
Put $\lambda=k+2m$ and $\mu=k'+2m'$ modulo $8$.  Then
$\alpha\equiv-\lambda$, $\gamma\equiv-\mu$, so $\lambda,\mu$ are odd and
\eqref{eq:cs-eight-match} gives $\mu=\pm\lambda$.  Equations
\eqref{eq:cs-parities} and \eqref{eq:cs-15569-branch} now imply
\[
 \operatorname{red}_{113}Q=(\lambda+4)g+(m+1)T_+,
 \qquad
 \operatorname{red}_{113}R=\pm\lambda g+mT_+.
\]
The possible abscissae and return residues are therefore
\begin{equation}\label{eq:cs-four-residues}
\begin{array}{cc|c|c}
\bar a&\bar c&\bar f&\bar f^{56}\\ \hline
11&81&21&-1\\
16&87&71&-1\\
81&11&29&-1\\
87&16&101&-1.
\end{array}
\end{equation}
All numerators and denominators are nonzero.  For a short check of the
last column, $21^4=8$ and $21^{28}=98$ in $\F_{113}$, while $98^2=-1$.
The other entries satisfy $71=26^2\cdot21$, $29=30^2\cdot21$, and
$101=8^2\cdot21$.  Thus each return is a nonsquare.  This contradicts
$f\in\Q_{113}^{*2}$ and completes the proof.
\end{proof}

\begin{corollary}\label{cor:cs-15569-layer}
For a strict interior rational-distance configuration, put
$t=r_{00}r_{11}+r_{10}r_{01}$.  Then $\sf(t^2-1)\ne15569$.
\end{corollary}
\begin{proof}
Theorem~\ref{thm:ell-exact-source} gives the source
\eqref{eq:cs-source-pair} with $n=\sf(t^2-1)$ and
$f(Q,R)=1/(r_{00}+r_{11})^2$.  Apply
Theorem~\ref{thm:cs-15569}.
\end{proof}

\section{A reciprocity obstruction on the complete source curve}
\label{sec:hasse-families}

We now retain the fourth square in the ordered distance construction.
For $p\in\Q$, $p>1$, put $a=p-1$ and $b=p+1$.  The complete source
is the intersection of three quadrics
\begin{equation}\label{eq:hp-full}
 C_p:\quad pV^2=2a(U^2-pZ^2),\qquad
 pW^2=2b(U^2+pZ^2),\qquad pT^2=aU^2+pbZ^2
 \quad\subset\PP^4_{[Z:U:V:W:T]}.
\end{equation}
On $Z\ne0$, the parameter $\rho=U^2/(pZ^2)$ gives the four ordered
squares
\begin{equation}\label{eq:hp-four}
 p\rho,\quad 2(p-1)(\rho-1),\quad
 2(p+1)(\rho+1),\quad (p-1)\rho+p+1.
\end{equation}
The smooth projective model has genus five: its degree-sixteen map to
$\PP^1_\rho$ has five branch values
$0,1,-1,-b/a,\infty$, with inertia of order two at each.
Thus $2g-2=-32+5\cdot8=8$.
Write $\operatorname{Loc}(p)$ for the assertion that the four expressions
in \eqref{eq:hp-four} are simultaneous nonzero squares at a finite
$\rho\notin\{0,1,-1\}$ over each completion of $\Q$.
The real solutions can be chosen in the strict distance region: take
$\rho>1$ sufficiently close to $1$.  Indeed, if
$F=a\rho+b$, $u^2=p\rho$, $v^2=2a(\rho-1)$ and
$w^2=2b(\rho+1)$, the strict margin is
$2F-4|u|-|vw|$, whose limit is $4p-4\sqrt p>0$.

The obstruction below uses every source point on \eqref{eq:hp-full}.
Its formula depends only on the valuations of $p,p-1,p+1$.
We use the additive Hilbert symbol notation $(c,d)_v\in\operatorname{Br}(\Q_v)[2]$,
and write its bit as $2\operatorname{inv}_v(c,d)_v\in\F_2$.
For an odd prime $q$, and $c=q^r c_0,d=q^s d_0$ with units $c_0,d_0$,
\begin{equation}\label{eq:hp-hilbert}
 (-1)^{2\operatorname{inv}_q(c,d)_q}
 =(-1)^{rs(q-1)/2}\leg{c_0}{q}^{s}\leg{d_0}{q}^{r}.
\end{equation}
For odd $u,v\in\Z_2^\times$, the dyadic rule is
\begin{equation}\label{eq:hp-hilbert2}
 2\operatorname{inv}_2(2^r u,2^s v)
 =\frac{(u-1)(v-1)}4+r\frac{v^2-1}8+s\frac{u^2-1}8\pmod2.
\end{equation}

\subsection{The obstruction and its arithmetic domain}

Assume
\begin{equation}\label{eq:hp-domain}
 p(p^2-1)=nk^2,\qquad n>1\text{ odd and squarefree},\qquad k\in\Q^\times.
\end{equation}
For parameters satisfying $\operatorname{Loc}(p)$ define
\begin{align}
 A(p)&=\sum_{\substack{q\text{ odd},\ (2/q)=-1\\v_q(a)>0}}\frac{v_q(a)}2,
 &B(p)&=\sum_{\substack{q\text{ odd},\ (2/q)=-1\\v_q(b)>0}}\frac{v_q(b)}2
 \quad\text{in }\F_2.\label{eq:hp-AB}
\end{align}
If $v_q(p)<0$ is odd, choose $e\in\F_q$ with $e^2=2$ and define
$\sigma_q\in\F_2$ by
\[
 (-1)^{\sigma_q}=\leg{2+e}{q},\qquad
 P(p)=\sum_{\substack{q\text{ odd}\\v_q(p)<0\text{ odd}}}\sigma_q.
\]
This definition is independent of the choice of $e$, since
$(2+e)(2-e)=2$ is a nonzero square.  The existence of $e$ and the
integrality of the half-valuations follow from the next lemma.
Put
\begin{equation}\label{eq:hp-parity}
 \mathcal T(p)=
 \begin{cases}
 B(p)+P(p),&v_2(p)\ne0,\\
 [v_2(a)=5]+B(p)+P(p),&p\equiv1\pmod {16},\\
 [v_2(b)=5]+A(p)+P(p),&p\equiv15\pmod {16}.
 \end{cases}
\end{equation}
Here brackets denote an indicator, and unit congruences are interpreted
in $\Z_2^\times$.

\begin{theorem}\label{thm:hp-uniform}
Suppose \eqref{eq:hp-domain} holds.  If $\operatorname{Loc}(p)$ holds
and $\mathcal T(p)=1$, then $C_p(\Q)=\varnothing$.
\end{theorem}

We prove the theorem by constructing a quaternion algebra on a
coordinate-omission curve and calculating its invariant at every place.
The sum is $\mathcal T(p)/2$.  A rational point would make this sum zero
by Hilbert reciprocity.

Omitting $U$ and $W$, respectively, gives the genus-one curves
\begin{align}
 D_U:&\quad v^2=2t^2-4pz^2,\qquad
 aw^2=2b(t^2-2z^2),\label{eq:hp-DU}\\
 D_W:&\quad V^2=2T^2-4pZ^2,\qquad
 aU^2=pT^2-pbZ^2.\label{eq:hp-DW}
\end{align}
These are projective equations in the four displayed coordinates.
Both omission maps are everywhere defined: if all retained coordinates
vanish, \eqref{eq:hp-full} forces the omitted coordinate to vanish too.

\begin{lemma}\label{lem:hp-norms}
Under $\operatorname{Loc}(p)$, the numbers $p,a,b$ are norms from
$\Q(\sqrt2)$.  At an odd prime where $2$ is nonsquare their valuations
are even.  Under \eqref{eq:hp-domain}, the possible dyadic branches are
\[
 v_2(p)\ge4\text{ even},\quad v_2(p)\le-4\text{ even},\quad
 p\equiv1\pmod {16},\quad p\equiv15\pmod {16}.
\]
In the last two branches, respectively, $[a]_2=[2]_2$ and
$[b]_2=[2]_2$, and the positive valuation is odd and at least five.
\end{lemma}
\begin{proof}
At a completion where $2$ is nonsquare, a point of \eqref{eq:hp-DU}
has $z,w\ne0$.  Its equations give the exact norm witnesses
\[
 \Nm\left(\frac{v+t\sqrt2}{2z(1+\sqrt2)}\right)=p,
 \qquad
 \Nm\left((2+\sqrt2)\frac b w(t+z\sqrt2)\right)=ab.
\]
At an odd inert prime the valuation of a norm is even.  At a pole of
$p$, the valuations of $p,a,b$ agree.  Away from poles, $a,b$ are
integral and cannot both have positive valuation.  The even valuation
of $ab$ therefore makes each positive valuation even.  Thus
$(2,a)_q=(2,b)_q=0$ at every odd place; their real invariants also vanish.
Reciprocity gives vanishing at $2$.  The quadratic norm criterion gives
the asserted global norms.

The norm classes at $2$ are $1,-1,2,-2$.  Since
$v_2(pab)$ is even, a nonzero $v_2(p)$ is even.  Depth $2$ is excluded
by the unit part of $a$: if $p$ has depth $2$ it is $3\pmod8$, and
if $p=u/4$ with $u\equiv1,7\pmod8$, it is $5,3\pmod8$.
For a unit $p$, its norm condition gives $p\equiv1,7\pmod8$;
the norm conditions on $b,a$ sharpen these to $1,15\pmod {16}$.
We spell out the four-square restriction at $2$.
Write $A=v_2(p)$ and $R=v_2(\rho)$.  Since $p\rho$ is square,
$A+R$ is even.  If $A,R$ are both nonzero and have the same sign,
the valuation of $2(p-1)(\rho-1)$ is odd.  If their signs are opposite
and both are even, the same conclusion holds.  If their signs are
opposite and both are odd, the fourth radicand has a unique term of
negative odd valuation: it is $(p-1)\rho$ when $A>0$, and
$p(\rho+1)$ when $A<0$.  Finally $A=R=0$ is impossible: the first
square gives $p\equiv\rho\pmod8$, so one of the middle two radicands
has valuation three.  Exactly one of $A,R$ is therefore zero and the
other is even.

For a unit $p$, the middle two squares force both $v_2(p-1)$ and
$v_2(p+1)$ odd, so $p\equiv1$ or $7\pmod8$.  If $R>0$, the fourth
square excludes $p\equiv1\pmod8$, because its valuation would be one.
Thus $[p]_2=-1$; the product of the middle two squares gives
$[p^2-1]_2=[\rho^2-1]_2=-1$.
If $R<0$, the fourth square excludes $p\equiv7\pmod8$, because
$(p-1)\rho$ would have odd valuation.  Now $[p]_2=1$ and
$[p^2-1]_2=[\rho^2-1]_2=1$.
In both cases $[pab]_2=1$.
For $p\equiv1\pmod {16}$, $[p]_2=1$ and $[b]_2=2$, whence
$[a]_2=2$; for $p\equiv15\pmod {16}$, $[p]_2=-1$ and
$[a]_2=-2$, whence $[b]_2=2$.  The corresponding valuation is odd
and at least five.
\end{proof}

\subsection{A rational anchor and the quaternion class}

Temporarily suppose an actual rational point
$Q_0=[z:v:w:t]\in D_U(\Q)$ is given.  Its $z,w,t$ are nonzero.
For instance, $z=0$ gives $v^2=2t^2$, and $w=0$ gives $t^2=2z^2$;
$t=0$ is excluded over $\R$.  Also $v\ne0$: otherwise
$2p=(t/z)^2$, contrary to the even valuation of $pab$ at $2$.
Define
\begin{equation}\label{eq:hp-anchor}
 z_1=aw-2bz,\qquad u_1=pb(w+2t+2z),\qquad
 t_1=-b(aw+2pt+2az).
\end{equation}
This is a rational point on $aU^2-pT^2+pbZ^2=0$, since
\[
 au_1^2-pt_1^2+pbz_1^2
 =2p^2b(aw^2-2bt^2+4bz^2)=0.
\]
The linear transformation $[z:t:w]\mapsto[z_1:u_1:t_1]$ has determinant
$-8p^2b^2$ and hence never produces the zero tuple.  Choose the sign of
$w$ so that $z_1t-zt_1\ne0$.  Such a choice exists: the sum for the two
signs is $4abz(t+z)$, and $t=-z$ is impossible in \eqref{eq:hp-DU}.
Fix this choice.  On $D_W$ put
\begin{equation}\label{eq:hp-polar}
 L_U=vV-2tT+4pzZ,\qquad
 L_V=au_1U-pt_1T+pbz_1Z,\qquad
 F=\frac{L_UL_V}{V^2},\qquad \alpha=(-a,F).
\end{equation}

Here is the divisor calculation underlying this construction.
The projections from $D_W$ to the two conics in \eqref{eq:hp-DW} have
degree two.  Let $E_U,E_V$ be the degree-two fibres over their respective
rational anchors $[z:v:t]$ and $[z_1:u_1:t_1]$, and let $H_V$ be the
degree-four coordinate section $V=0$.  A tangent line on a smooth
conic meets it twice at its point of tangency.  Consequently
\begin{equation}\label{eq:hp-divisor}
 \operatorname{div}(F)=2(E_U+E_V-H_V).
\end{equation}
All residues of $(-a,F)$ vanish, so it defines a class on the smooth
proper curve $D_W$.

For completeness, the associated elliptic curve and ordered classes are
\begin{equation}\label{eq:hp-elliptic}
 E:\ y^2=x(x-L)(x-R),\qquad
 L=\frac{8p}{a^2},\quad R=-\frac{8pb}{a^3},\qquad
 \chi=(2,pa,2pa),\quad\beta=(-a,-a,1).
\end{equation}
The row $\chi$ describes the actual curve $D_W$.  Indeed, with
$\xi_0=2aU,\xi_L=2V,\xi_R=2T,\zeta=a^2Z$ its equations become
\[
 2\xi_0^2-pa\xi_L^2=L\zeta^2,
 \qquad 2\xi_0^2-2pa\xi_R^2=R\zeta^2.
\]
In particular its degree-four covering map is
$x=8U^2/(a^2Z^2)$, $y=16pUVT/(a^4Z^3)$.
Write $T_0,T_L,T_R$ for the named two-torsion points, and
$\sigma_i$ for the sign change of a coordinate $i$.
The diagonal degree-four covering map is the intrinsic map
$P\mapsto[P-\sigma_Z(P)]$ to the Jacobian: its ordered root functions
are $2U/(aZ)$, $2V/(a^2Z)$ and $2T/(a^2Z)$, with the constants in
$\chi$.  It sends the boundary sections $U=0,V=0,T=0,Z=0$ to
$T_0,T_L,T_R,0$, respectively.
Choose a geometric point $O_C$ above $Z=0$ to compute Abel sums.
Then $\sigma_Z$ acts by negation and the covering map acts by
multiplication by two.  Each other reflection acts as
$P\mapsto t_i-P$, with $2t_i=0$ because it commutes with $\sigma_Z$.
At a fixed point of $\sigma_i$, the covering image is $2P=t_i$;
the boundary images therefore give
$t_U=T_0,t_V=T_L,t_T=T_R$.
Every fibre of the corresponding conic projection has Abel sum
$P+\sigma_i(P)=t_i$, including ramified fibres counted twice.
Thus the Abel sums of $E_U,E_V$ are $T_0,T_L$.
Since $H_V\sim2E_V$, we obtain
$[E_U+E_V-H_V]=T_0-T_L=T_R$.

For the cyclic algebra in \eqref{eq:hp-polar}, the divisor identity
\eqref{eq:hp-divisor} identifies the Picard cocycle as
$\gamma\mapsto\epsilon_{-a}(\gamma)T_R$:
over $\Q(\sqrt{-a})$ the norm of $E_U+E_V-H_V$ is its double,
which is $\operatorname{div}(F)$.
The Weil pairings of $T_R$ with $T_0,T_L,T_R$ are $-1,-1,1$,
respectively, so the ordered two-descent row is $\beta=(-a,-a,1)$.

We shall use only the constant-evaluation consequence of the Picard
cocycle.  Let $k=\Q_v$ and choose any point of $D_W(k)$, which exists
by $\operatorname{Loc}(p)$.  This point splits the degree map on
$\Pic(\overline{D_W})$, giving the Galois module decomposition
$\Pic(\overline{D_W})=E(\overline k)\oplus\Z$.
The coefficient $-a$ becomes a square over $\overline k$, so $\alpha$
lies in the algebraic Brauer group
$\operatorname{Br}_1(D_W)=\ker(\operatorname{Br}(D_W)\to
\operatorname{Br}(\overline{D_W}))$.
The Hochschild--Serre exact sequence
\cite[\S2, (2.1)]{creutz2020} contains
\[
 \operatorname{Br}(k)\longrightarrow\operatorname{Br}_1(D_W)
 \longrightarrow H^1(k,\Pic(\overline{D_W})).
\]
A Kummer class $\beta\in E(k)/2E(k)\subset H^1(k,E[2])$ maps to zero
in $H^1(k,E)$ by the multiplication-by-two exact sequence.  Since
$H^1(k,\Z)=0$, an algebraic Brauer class with this Picard cocycle
comes from $\operatorname{Br}(k)$ and has constant evaluation.
For a torsor the cocycle is computed modulo its torsor class; that
class becomes zero as soon as a local point is chosen.  Thus it
suffices to prove local solubility of the full cover representing
$\beta$.

\subsection{The partner cover at every place}

Put $c=-a$.  The complete partner cover is
\begin{equation}\label{eq:hp-partner}
 cA^2-cB^2=L\zeta^2,\qquad cA^2-C^2=R\zeta^2.
\end{equation}
Two useful projective points are
\begin{align}
 &[1:1:\sqrt c:0] &&\text{if $c$ is square},\label{eq:hp-partner-O}\\
 &\left[\frac{4p}{a^2}\sqrt{\frac b{2p}}:\frac{4p}{a^2}:0:1\right]
 &&\text{if $b/(2p)$ is square}.\label{eq:hp-partner-R}
\end{align}
These follow from $(R-L)/c=(4p/a^2)^2$.

At an odd prime with $v_q(p)>0$, $c$ is square; with $v_q(a)>0$,
$b/(2p)$ is square.  At a pole with $2$ square, the latter criterion
again applies.  At a pole with $2$ nonsquare, write
$p=\pi^{-2h}p_0$, $a=\pi^{-2h}a_0$, $b=\pi^{-2h}b_0$.
Replacing $(A,B,C)$ by $(\pi^{2h}A',\pi^{2h}B',\pi^hC')$ gives a
unit model whose three elliptic roots reduce to $0,8/p_0,-8/p_0$.
It is a smooth genus-one curve over $\F_q$, so the Hasse bound and
Hensel's lemma supply a point.  The same argument applies whenever
$p,a,b$ are units.

It remains to consider $v_q(b)>0$.  If $2$ is square, $c$ is square.
Otherwise write $v_q(b)=2h$ and let $B_0$ be the unit squareclass of
$\pi^{-2h}b$.  The Kummer row at $T_0$ is $(2B_0,2,B_0)$.
If $B_0=1$ this is $\beta$.  If $B_0=2$, choose a residue $X$ with
$X$ nonsquare and $X+2$ a nonzero square.  There are
$(q+\leg{-2}{q})/4>0$ such residues: expanding the two character
indicators reduces the count to
$\sum_X\leg{X(X+2)}q=-1$, which follows from
$(X+1-Y)(X+1+Y)=1$.
Lifting $x=X$ gives an actual point of $E(\Q_q)$ with row $(2,1,2)$,
since $L\equiv-2,R\equiv0$.  Adding $T_0$, whose row is now $(1,2,2)$,
gives $\beta$.  The simultaneous square roots in the Kummer row give
one point of \eqref{eq:hp-partner}.

At $2$, the branches of Lemma~\ref{lem:hp-norms} suffice.
If $v_2(p)\ge4$, $c$ is square.  If $v_2(p)=-2h\le-4$, normalize
$z=1$ in the actual rational anchor.  Its norm condition leaves
$[p]_2=1$ or $-1$.  In the first case the equation
$v^2=2t^2-4p$ forces $v_2(t)=1-h$: a smaller valuation has odd leading
depth and a larger one gives squareclass $-1$.
Then $(2b/a)t^2(1-2/t^2)$ has squareclass $2$, contradicting the
$w$-equation.  Thus $[p]_2=-1$ and $c$ is square.
If $p\equiv1\pmod {16}$, $b/(2p)\in1+8\Z_2$, so
\eqref{eq:hp-partner-R} applies.

Finally suppose $p\equiv15\pmod {16}$.  Normalizing $z=1$, the
first anchor equation forces $v_2(t)\ge2$; the second then implies
that $2b$ is square.  Choose local $r,s,h_d,h_Z$ with
\[
 r^2=2b,\quad s^2=3-p,\quad h_d^2=-a/2,\quad
 h_Z^2=\frac{p^2-4p-1}{2pa}.
\]
Here $3-p=4(1-b/4)$ and $1-b/4\in1+8\Z_2$, so the indicated $s$ exists.
The last two unit radicands are squares because
$-a/2\equiv1\pmod8$ and
$h_Z^2-1=-b^2/(2pa)\in2^8\Z_2$.
The actual point
\[
 [Z:U:V:T]=[h_Z/2:r/4:1:rs/(4h_d)]\in D_W(\Q_2)
\]
has row $\chi=(2,2,1)=\beta$ over $\Q_2$ and hence supplies a point
of \eqref{eq:hp-partner}.  Over $\R$ one may use
$[0:\sqrt{8p/a^3}:\sqrt{8pb/a^3}:1]$.
The partner is therefore everywhere locally soluble, and the evaluation
of $\alpha$ is constant at every place.

\subsection{Evaluation of the quaternion algebra}

We give the local calculations leading to \eqref{eq:hp-parity}.
At a tangent fibre the extension in \eqref{eq:hp-divisor}, rather than
the value of a vanishing rational function, is used.  When the missing
coordinate is locally a square, the actual points
\[
 P_U=[z:u_0:v:t],\quad u_0^2=\frac p a(t^2-bz^2),\qquad
 P_V=[z_1:u_1:v_1:t_1],\quad v_1^2=2t_1^2-4pz_1^2
\]
satisfy
\begin{equation}\label{eq:hp-tangent-eval}
 \alpha(P_U)=(-a,-pL_V(P_U)),\qquad
 \alpha(P_V)=(-a,-2bL_U(P_V)).
\end{equation}
To see this, the polar identities are
\[
 (zV-vZ)^2-2(zT-tZ)^2=-2zZL_U,
 \qquad
 a(z_1U-u_1Z)^2-p(z_1T-t_1Z)^2=-2z_1ZL_V.
\]
Dividing by the square of $tZ-zT$ in the first case gives the limiting
coefficient $-4p/v^2$ for $L_U$.  Dividing by the corresponding square
in the second gives $-p^2b/(2au_1^2)$ for $L_V$; removing squares and
norms from $\Q_v(\sqrt{-a})$ gives the factor $-2b$ in
\eqref{eq:hp-tangent-eval}.  If $u_1=0$, use $U$ as parameter and
$L_V/U^2\to-a/2$.  If $z_1=0$, use the chart $T=t_1$.
The choice $z_1t-zt_1\ne0$ ensures the other polar is nonzero.

For an odd prime the evaluation table is
\begin{equation}\label{eq:hp-odd-table}
\begin{array}{c|c}
 \text{parameter branch}&2\operatorname{inv}_q\alpha\\ \hline
 v_q(p),v_q(a),v_q(b)=0&0\\
 v_q(p)>0&0\\
 v_q(a)>0&2\operatorname{inv}_q(-a,-2)\\
 v_q(b)>0,\ (2/q)=1&0\\
 v_q(b)>0,\ (2/q)=-1&v_q(b)/2\\
 v_q(p)<0\text{ even}&0\\
 v_q(p)<0\text{ odd}&\text{bit of }\leg{-(2+e)}q,\quad e^2=2.
\end{array}
\end{equation}
Here and below equality of entries is in $\F_2$.
We justify the entries, including arbitrary valuation depths.
At a good unit prime, the source and its two conics have smooth proper
models.  Normalize the anchor primitively.  The anchor transformation
has unit determinant, so both tangent sections are primitive.  Their
horizontal divisors are even and neither has a vertical divisor.
The quaternion extends over the model and specializes to
$\operatorname{Br}(\F_q)=0$.  At a positive-$p$ prime, $-a$ is square.

At a positive-$a$ prime where $2$ is square, use the actual point
$[Z:U:V:T]=[1:\sqrt p:0:\pm\sqrt{2p}]$.
Normalize $\min(v_q(t),v_q(z))=0$ and choose the sign so that $L_U$
is a unit.  Direct substitution of \eqref{eq:hp-anchor} gives
$L_V/L_U\equiv-p^2b\pmod q$; all terms involving $aw$ have positive
valuation because $v_q(a)+2v_q(w)\ge0$.
Removing the square pole $V^2$ leaves the class $-b=-2$.
If $2$ is nonsquare, $v_q(a)$ is even.  In the same normalization
$z_1\equiv-4z$, $t_1\equiv-4t$, and the missing square is
$2t_1^2-4pz_1^2\equiv16v^2$.
Choose $v_1\equiv4v$; then $L_U(P_V)\equiv8v^2$ is a unit.
Both entries of the symbol in \eqref{eq:hp-tangent-eval} have unit
squareclass, so its invariant is zero, as is $(-a,-2)_q$.

At a positive-$b$ inert prime put $v_q(b)=2h$ and normalize
$\min(v_q(t),v_q(z))=0$.  The unramified norm $t^2-2z^2$ gives
$v_q(w)=h$ and $v_q(z)=0$.  The anchor transformation has
$v_q(z_1)=h$, $v_q(t_1)\ge2h$.  Hence the missing square at $P_V$
has class $-p=1$.  Choose its sign so that
$L_U(P_V)/z_1\equiv\pm2v-4z$ is nonzero.
Thus $v_q(L_U(P_V))=h$.  The first entry $-a$ has class $2$, and
\eqref{eq:hp-hilbert} gives the bit $2h+h=h$.
If $2$ is square, $-a$ is square instead.

At a pole $v_q(p)=-2h$ with $2$ nonsquare, normalize $z=1$.
The $w$-equation forbids $v_q(t)<0$, and the $v$-equation then
forces $-p$, and hence $-a$, to be square.  If $2$ is square, set
$Z=\pi^hZ'$.  Writing $p=\pi^{-2h}p_0$, $a=\pi^{-2h}a_0$,
$b=\pi^{-2h}b_0$ gives the good model
\[
 V^2-2T^2+4p_0Z'^2=0,\qquad
 a_0U^2-p_0T^2+p_0b_0Z'^2=0.
\]
Normalize $[v:t:z/\pi^h]$ primitively.  If $t$ is nonunit, $-p_0$
is square.  Otherwise the transformed second anchor, multiplied by
$\pi^{4h}$, reduces to
$[0:p_0^2(w+2t):-p_0^2(w+2t)]$, a primitive point.
The original second polar differs from its primitive-model polar by
$\pi^{-6h}$, a square factor.  The good-model argument therefore
again gives zero.

At an odd-depth pole $v_q(p)=-(2h+1)$, $2$ is square.  Normalize $t=1$.
The first anchor equation forces $v_q(z)\ge h+1$ and
$v,w\equiv\pm\sqrt2$.  On $D_W$ choose
$[Z:U:V:T]=[0:u:e:1]$, where $u^2=p/a$, $u\equiv1$ and
$e^2=2$ has sign opposite to $v$.
Then $L_U\equiv-4$; $L_V$ has valuation $-3(2h+1)$ and leading
unit $2p_0^3(w+2)$.  Thus $[F]=[(-a)(w+2)]$ and
$\alpha=(-a,-(w+2))$.  Formula \eqref{eq:hp-hilbert} gives the last
line of \eqref{eq:hp-odd-table}.

We next calculate the dyadic entries.  In either nonunit branch $-a$
is square, as proved above, and the value is zero.
For $p\equiv1\pmod {16}$ put $h=v_2(a)$, $d=(h-1)/2$ and normalize
$z=1$.  The anchor has $t=2t_0$, $v=2v_0$, with $t_0,v_0$ odd and
$v_0^2=2t_0^2-p$; moreover $v_2(w)=(3-h)/2$.
Set
\[
 C=\frac{aw}{2b},\quad g=\frac{t_1}{2z_1},\quad
 z_1=-2b(1-C),\quad v_2(C)=d.
\]
Then $g$ is odd and
\[
 g-t_0=\frac{a(t_0+1/2)+C(t_0+b/2)}{1-C},\qquad
 v_2(g-t_0)\ge d+1.
\]
Choose $e^2=2g^2-p$ close to $-v_0$, and set
$K=v_0e-2t_0g+p$.  Factoring
$e^2-v_0^2=2(g-t_0)(g+t_0)$ gives
$v_2(e+v_0)\ge d+2$, whence $[K]=[-2]_2$.
The actual point $P_V$ with $v_1=2z_1e$ has $L_U=4z_1K$.
Formula \eqref{eq:hp-tangent-eval} becomes
$\alpha_2=(-2,(1-C)K)$.
By \eqref{eq:hp-hilbert2} the $K$ term has bit one, and the
$1-C$ term has bit one exactly when $d=2$.  Therefore
\begin{equation}\label{eq:hp-dyadic-one}
 2\operatorname{inv}_2\alpha=1+[v_2(a)=5].
\end{equation}

For $p\equiv15\pmod {16}$ put $h=v_2(b)$, $d=(h-1)/2$ and again
$z=1$.  Now $t=2t_0$ with $t_0$ even, $v=2v_0$,
$v_0^2=-p+2t_0^2$, and $v_2(w)=(h+1)/2$.
The transformed anchor has
$v_2(z_1)=(h+3)/2$, $v_2(t_1)=h+2$.
Thus $g=t_1/(2z_1)$ has depth $d$, and $e^2=2g^2-p$ is a square.
Choose its sign so that $v_0e\equiv3\pmod4$.
Since $p\equiv-1\pmod {32}$, the number
$K=v_0e-2t_0g+p$ is $14$ or $2\pmod {16}$; its valuation is one
and its odd part is $7$ or $1\pmod8$.
Hence $(2,K)_2=0$ and \eqref{eq:hp-tangent-eval} gives
\begin{equation}\label{eq:hp-dyadic-minus}
 2\operatorname{inv}_2\alpha
 =2\operatorname{inv}_2(2,w/z)+[v_2(b)=5].
\end{equation}
Indeed $z_1=aw(1-2b/(aw))$, the correction has depth $d$, and
$p,a,b$ are norms from $\Q_2(\sqrt2)$.

The invariant at $\R$ is $1/2$.  To verify the sign in the chosen
rational gauge, note that
\[
 (2pt+2az)^2-a^2w^2
 =2\bigl((p^2+1)t^2+4patz+4paz^2\bigr)>0.
\]
The quadratic is positive definite, and $2pt+2az$ has the sign of $t$.
Thus $t_1$ has the opposite sign to $t$.
Cauchy--Schwarz on the two positive conics shows that $L_U$ and $L_V$
have respective signs $-\operatorname{sgn}(tT)$ and
$-\operatorname{sgn}(t_1T)$ away from their zeros.  Consequently
$F<0$.  Since $-a<0$, the real invariant is $1/2$.
The evaluation at a zero or pole is obtained by continuity of the
proper class, so this value holds on every real component.

\subsection{Reciprocity and the exclusion theorem}

The dyadic ratio in \eqref{eq:hp-dyadic-minus} has an exact global
support formula.  At an odd prime where $2$ is nonsquare,
\begin{equation}\label{eq:hp-wsupport}
 v_q(w/z)=
 \begin{cases}
 -v_q(a)/2,&v_q(a)>0,\\
 v_q(b)/2,&v_q(b)>0,\\
 0,&\text{otherwise}.
 \end{cases}
\end{equation}
For a unit or positive-$p$ branch, normalize $\min(v_q(t),v_q(z))=0$.
The unramified norm $t^2-2z^2$ forces $w$ to be a unit, and $z$
cannot be nonunit, since then $v^2$ has leading squareclass $2$.
For $v_q(a)=2h>0$ normalize $w=1$: the same norm has depth $2h$,
and the $v$-equation forces $v_q(z)=h$.
For $v_q(b)=2h>0$ the normalization above gives $v_q(w)=h$, $v_q(z)=0$.
At a pole, normalize $z=1$.  A negative valuation of $t$ would make
$w^2$ have class $2$; with $t$ integral, $t^2-2$ is a unit, so $w$
is a unit.  These arguments prove \eqref{eq:hp-wsupport} at all depths.
Hilbert reciprocity for the literal rational number $w/z$ now gives
\begin{equation}\label{eq:hp-WR}
 2\operatorname{inv}_2(2,w/z)=A(p)+B(p).
\end{equation}

Apply reciprocity also to the constant symbol $(-a,-2)$.
Its real bit is one; its odd contributions are the positive-$a$
entries of \eqref{eq:hp-odd-table} and, at odd-depth poles, the
character $\leg{-1}q$.
At $2$ its bit is one in the $p\equiv1\pmod {16}$ branch and zero
in the other branches.
Adding this zero global sum to the sum for $\alpha$ cancels the real
entry and every positive-$a$ entry.  At odd poles the remaining
character is $\leg{2+e}q$.  Hence
\[
 2\sum_v\operatorname{inv}_v\alpha
 =2\operatorname{inv}_2\alpha+
   2\operatorname{inv}_2(-a,-2)+B(p)+P(p)
 =\mathcal T(p),
\]
by \eqref{eq:hp-dyadic-one}, \eqref{eq:hp-dyadic-minus} and
\eqref{eq:hp-WR}.  This proves the claimed invariant formula.

To finish Theorem~\ref{thm:hp-uniform}, suppose a rational point of
$C_p$ exists.  Its $U$-omission supplies the single rational anchor
used throughout the construction.  A rational point on $C_p$ has
$Z,V,W,T\ne0$ by the arguments following \eqref{eq:hp-DU}, and
$U=0$ would make $pV^2=-2apZ^2<0$.  It therefore gives the
nondegenerate four squares in \eqref{eq:hp-four}.
The construction applies to this anchor and the image point on $D_W$.
Evaluation there is a quaternion algebra over $\Q$, so the sum of its
local invariants is zero.  This contradicts
$\mathcal T(p)=1$.  The theorem follows.

\subsection{A conic family with nonzero obstruction}

Let
\begin{equation}\label{eq:hp-conic}
 r^2+2s^2=8,\qquad p=3-s^2>1.
\end{equation}
Then $r^2=2(p+1)$ and
$[z:v:w:t]=[s:2(p-1):2r:r]$ is an actual rational point of $D_U$.
Every odd inert prime dividing $p-1$ would make
$s^2=2-(p-1)$ a nonsquare unit; thus $A(p)=0$.
Every pole of $p$ has even depth, so $P(p)=0$.
For rational points of this conic, $v_2(s)=1$ and $v_2(r)\ge3$.
Indeed write $r=R/D$, $s=S/D$ with $\gcd(R,S,D)=1$.
If $D$ were even, $R^2+2S^2=8D^2$ modulo $4$ would force $R,S$
even.  Thus $D$ is odd.  Reduction modulo $8$ gives $R,S$ even;
writing $R=2R_1$, $S=2S_1$ gives $R_1^2+2S_1^2=2D^2$.
Here $R_1$ is even and $S_1$ odd.  With $R_1=2R_2$, the identity
$2R_2^2+S_1^2=D^2$ modulo $8$ forces $R_2$ even.
Consequently $p\equiv-1\pmod {32}$ and
$v_2(p+1)=2v_2(r)-1$.
Formula \eqref{eq:hp-parity} reduces to
\begin{equation}\label{eq:hp-thinbit}
 \mathcal T(p)=[v_2(r)=3].
\end{equation}
Thus every locally soluble complete source in this conic family with
$v_2(r)=3$ is a genus-five failure of the Hasse principle.

\begin{theorem}\label{thm:hp-infinite}
For every integer $j\ge0$, set
\[
 N=3+154j,\quad M=N+2,\quad D=N^2+2M^2,\qquad
 r=\frac{8MN}{D},\quad s=\frac{2(N^2-2M^2)}D,\quad p=3-s^2.
\]
The curves $C_p$ form an infinite family of smooth genus-five curves
with points at every completion of $\Q$, including points in the
strict real distance region, and with no rational point.
\end{theorem}
\begin{proof}
The displayed parametrization satisfies \eqref{eq:hp-conic}.
The integers $M,N,D$ are odd, $\gcd(M,N)=1$ and $v_2(r)=3$.
The ratio $m=M/N$ decreases strictly from $5/3$ towards $1$.
On $1<m\le5/3$, the function
$s=2(1-2m^2)/(1+2m^2)$ is negative with derivative
$-16m/(1+2m^2)^2<0$; thus $p=3-s^2$ decreases as a function of $m$.
It follows that the distinct parameters $p$ increase from
$3719/3481>1$.
Their squarefree kernels in $p(p^2-1)$ are odd, since
$[p(p^2-1)]_2=1$.  Such a kernel cannot equal one:
if $p(p^2-1)=y^2$, the legs $(p^2-1)/y,2p/y$ and hypotenuse
$(p^2+1)/y$ give a rational right triangle of area one, contradicting
Lemma~\ref{lem:int:fermat}.
By \eqref{eq:hp-thinbit}, it remains to prove $\operatorname{Loc}(p)$.
We give the local proof directly.

At $2$ use $\rho=b/(3p+1)$.  Its four radicands are
\begin{equation}\label{eq:hp-rho-minus}
 \frac{r^2}{4}\frac{2p}{3p+1},\qquad
 4\frac{-p(p-1)}{3p+1},\qquad
 r^2\frac{2(2p+1)}{3p+1},\qquad
 4p\rho.
\end{equation}
The three unit factors displayed after their square factors lie in
$1+8\Z_2$, because $p\equiv-1\pmod {32}$.
At an odd prime with $v_q(p)>0$ or $v_q(p-1)>0$, the identities
$r^2=2(p+1)$ or $s^2=3-p$, respectively, make $2$ square locally.
Take $\rho=p$.  The first three radicands are then explicit squares,
and the fourth, $p^2+1$, has square residue $1$ or $2$.
If $v_q(p+1)>0$, use \eqref{eq:hp-rho-minus} again: the first and
third radicands divided by $r^2$ have residues $1/4$ and $1$;
the second divided by $4$ has residue $1$.
At a pole of $p$, the identity $s^2=3-p$ makes $-p$ a local square.
Put $h=1/p$ and $\rho=-(p+1)/(p+3)$.  The radicands become
\[
 (-p)\frac{1+h}{1+3h},\quad
 (-4p)\frac{(1-h)(1+2h)}{1+3h},\quad
 4\frac{1+h}{1+3h},\quad 4\frac{1+h}{1+3h}.
\]
Every remaining factor is a principal unit and hence square.
These witnesses treat all nonordinary odd primes and every depth.

At an ordinary prime $q$, the model has five distinct branch points
and smooth genus-five reduction.  If $q\ge101$, Hasse--Weil gives
\[
 \#C_p(\F_q)\ge q+1-10\sqrt q>0.
\]
Any such point gives a nondegenerate $\Q_q$-point.  At a finite branch
point $\rho_0$, write its one vanishing radicand as $u(\rho-\rho_0)$
with $u$ a unit, and take $\rho=\rho_0+q^2/u$.
The other three radicands remain square units.  At infinity the three
leading ratios to $p\rho$ are squares; taking
$\rho=1/(pq^2)$ gives four nonzero roots.  Thus normalized branch
points are sufficient for lifting.

Here is a short finite-field proof for the remaining primes.  It also
records the explicit arithmetic used for the small-prime part.
Let $\chi$ be the quadratic character of $\F_q$, extended by $\chi(0)=0$,
and put
\[
 H=\chi(-1),\quad T=\chi(2),\quad A=\chi(p),\quad B=\chi(p-1),
 \quad c=-\frac{p+1}{p-1}.
\]
The conic identity gives $\chi(p+1)=T$.
The five elliptic character quotients are
\begin{align*}
 E_0 &:y^2=p(p^2-1)x(x-1)(x+1),\\
 E_1 &:y^2=2p\,x(x-1)(x-c),\\
 E_2 &:y^2=2p(p^2-1)x(x+1)(x-c),\\
 E_3 &:y^2=(p+1)(x-1)(x+1)(x-c),\\
 E_4 &:y^2=p(p+1)c(x-1)(x+1)(x-c^{-1}).
\end{align*}
If $t_i=q+1-\#E_i(\F_q)$, expansion by the characters of the four
square roots gives
\begin{equation}\label{eq:hp-five-traces}
 \#C_p(\F_q)=q+1-\sum_{i=0}^4t_i.
\end{equation}
Indeed the products of three radicands give $E_0,E_1,E_2,E_3$;
the product of all four, after $x=1/\rho$, gives $E_4$.
The remaining characters have genus zero and their finite sums cancel
with their contributions at infinity.

Each $E_i$ has rational two-torsion.  On
$y^2=d\prod(x-e_i)$, the point $(e_i,0)$ is twice a rational point
if and only if $d(e_i-e_j)$ and $d(e_i-e_k)$ are both squares.
This follows by writing the tangent at a putative half-point through
$(e_i,0)$ and factoring the resulting quadratic; conversely those two
square roots give the two tangent slopes and a half-point.
Consequently the group order is $4\pmod8$ when there is no rational
point of order four, is divisible by eight when such a point exists,
and is divisible by sixteen when all four-torsion is rational.

Let $U_q,V_q,W_q$ be the largest integers at most $2\sqrt q$
congruent to $q+1\pmod8$, $q-3\pmod8$, and $q+1\pmod {16}$,
respectively.  Set
$\tau_q=-\sum_{x\in\F_q}\chi(x^3-x)$.
Then $t_0=TAB\tau_q$; for $q\equiv3\pmod4$, $\tau_q=0$ by
pairing $x$ and $-x$.
Applying the preceding halving criterion to the displayed roots gives
the following upper bounds on the sum in \eqref{eq:hp-five-traces}.
Only the sign patterns without an immediate rational branch point
need be listed:
\begin{center}
\begin{tabular}{ccl}
\toprule
$q\bmod8$&sign pattern&upper bound for $\sum t_i$\\
\midrule
$3$&all $A,B$&$3U_q+V_q$\\
$7$&$A=-1,B=1$&$2U_q+2V_q$\\
$1$&$A=B=-1$&$\tau_q+2U_q+2V_q$\\
$5$&$A=-1,B=1$&$\tau_q+4U_q$\\
$5$&$A=B$&$-\tau_q+2U_q+V_q+W_q$\\
\bottomrule
\end{tabular}
\end{center}
For example, in the first row exactly three of $E_1,\ldots,E_4$ have
order divisible by eight, and the remaining order is $4\pmod8$.
In the second and third rows two have each type; in the fourth all
four have order divisible by eight; in the last, one has full rational
four-torsion, two have an order divisible by eight, and one has order
$4\pmod8$.  These statements follow by inserting $A,B,H,T$ into
$d(e_i-e_j)$ and $d(e_i-e_k)$.
The omitted branch patterns are also explicit: when $T=1$, $\rho=1$
works for $A=1$, and $\rho=0$ works for $HB=1$; when $H=1,T=-1$,
$\rho=c$ works for $A=1,B=-1$.

The trace values needed below are
\[
\begin{array}{c|rrrrrrrrrrr}
q&5&13&17&29&37&41&53&61&73&89&97\\\hline
\tau_q&-2&6&2&-10&-2&10&14&-10&-6&10&18.
\end{array}
\]
They can be evaluated from the displayed finite sum; equivalently,
for $q=u^2+v^2$, $u$ odd, $v$ even and $u+v\equiv1\pmod4$, the
usual Gaussian evaluation of this sum is $\tau_q=2u$.
Substitution in the bounds proves positivity except for the following
cases, where the two conic conditions
$\chi(2(p+1))=1$ and $\chi(3-p)\in\{0,1\}$ are retained.
The congruence $N\equiv3$, $M\equiv5\pmod {7\cdot11}$ gives
$p\equiv1\pmod7$ and $p\equiv9\pmod {11}$.
At $17$ the remaining residues are $7,12$, and at $41$ they are
$7,35$; taking $\rho=p$ works because $p^2+1$ has respective square
residues $16,9$ and $9,37$.
At $13$ the sign case $A=-1,B=1$ would require $p=5$, but
$3-p=11$ is nonsquare.
At $53$ the remaining case is $A=-1,B=1$, and at $61$ the case
$A=B=-1$ is also settled by $\rho=-2$, whose radicands are
$-2p,-6(p-1),-r^2,s^2$; use $(3/53)=-1$ and $(3/61)=1$.
The residual cases have the following explicit nonzero square roots:
\begin{center}\small
\begin{tabular}{rrrr}
\toprule
$q$&$p\bmod q$&$\rho\bmod q$&roots in the order \eqref{eq:hp-four}\\
\midrule
11&9&4&$5,2,1,3$\\
19&2&3&$5,2,9,5$\\
19&11&5&$6,2,7,9$\\
19&13&10&$4,8,2,1$\\
19&17&4&$7,1,3,5$\\
23&17&5&$4,6,3,11$\\
29&7&22&$3,7,7,13$\\
29&25&4&$10,12,12,8$\\
43&2&3&$7,2,14,7$\\
43&6&10&$19,2,5,10$\\
43&11&9&$20,17,5,4$\\
43&21&14&$6,2,12,1$\\
43&25&35&$12,16,18,7$\\
43&29&5&$4,3,4,16$\\
43&31&10&$3,14,4,17$\\
43&32&3&$15,9,7,13$\\
43&33&12&$3,4,14,17$\\
43&36&23&$21,11,20,5$\\
61&16&3&$29,11,21,1$\\
61&42&12&$4,29,9,13$\\
61&49&12&$10,18,18,4$\\
\bottomrule
\end{tabular}
\end{center}
Every row is checked by squaring its four entries in \eqref{eq:hp-four};
all roots are nonzero and therefore lift by Hensel's lemma.
This completes the proof of finite-place solubility.
At the real place the strict open interval described after
\eqref{eq:hp-four} supplies the required points.
Now \eqref{eq:hp-thinbit} and Theorem~\ref{thm:hp-uniform} give
$C_p(\Q)=\varnothing$ for every parameter in the family.

There are also infinitely many geometric isomorphism classes among these
curves. A fixed smooth curve of genus five has finite automorphism group:
its automorphism scheme is of finite type, for example through the
tricanonical embedding, and its tangent space is $H^0(C,T_C)=0$ because
$\deg T_C=-8$. It therefore has finitely many subgroups isomorphic to
$(\Z/2)^4$. Each such subgroup gives one quotient with its unordered
branch configuration. Here that configuration is
$\{0,1,-1,c,\infty\}$, with $c=-(p+1)/(p-1)$. For a fixed configuration,
there are finitely many ways to choose and normalize three ordered branch
points, and hence finitely many values of $c$. Since $p=(c-1)/(c+1)$,
a fixed geometric curve occurs for only finitely many of our distinct
parameters.
\end{proof}

\section{Theta characteristics and a complete double cover}
\label{sec:theta}

The fixed opposite-distance difference $1/5$ gives a particularly explicit
connection between the distance equations, the geometry of a plane quartic,
and descent on its Jacobian.  In this section we determine all the odd theta
characteristics of that quartic and identify the double cover which restores
the four distances.  A rational point at infinity then gives a local
solubility theorem for every finite multiplication cover, including the
pullbacks which retain this double cover.

\subsection{The quartic quotient and its metric interpretation}

Let $C\subset\PP^4_{[P:Q:R:z:T]}$ be the curve
\begin{equation}
 \begin{split}
 P^2&=z^2+50T^2,\\
 Q^2&=z^2+14zT+51T^2,\\
 R^2&=z^2-14zT+51T^2.
 \end{split}
 \label{eq:theta-C}
\end{equation}
We use $C$ as a complete source: all three equations are imposed at the same
value of $z/T$.  On $T=1$, put
\begin{equation}
 \begin{aligned}
 x&=\frac{50-P-7z}{100},&
 y&=\frac{50-P+7z}{100},\\
 (r_0,r_1,r_2,r_3)&=
 \left(\frac{P-1}{10},\frac Q{10},\frac R{10},\frac{P+1}{10}\right).
 \end{aligned}
 \label{eq:theta-return}
\end{equation}
The strict positive locus $\mathcal U\subset C(\R)$ is defined by
\begin{equation}
 T=1,\qquad P>1,\quad Q>0,\quad R>0,\qquad P+7|z|<50.
 \label{eq:theta-U}
\end{equation}
It corresponds exactly to the strict distance configurations with
$r_3-r_0=1/5$.  For completeness, the first distance identity follows from
\[
 5000\bigl(r_0^2-x^2-y^2\bigr)
   =49(P^2-z^2-50)=0.
\]
The other three follow from
\[
 r_1^2-r_0^2=1-2x,\qquad
 r_2^2-r_0^2=1-2y,\qquad
 r_3^2-r_0^2=2-2x-2y.
\]
Conversely, the inverse formulas are
\[
 P=5(r_3+r_0),\quad Q=10r_1,\quad R=10r_2,
 \quad z=\frac{50(y-x)}7.
\]
The inequalities in \eqref{eq:theta-U} are precisely the positivity of the
chosen roots and $0<x,y<1$.

\begin{proposition}
\label{prop:theta-quotient}
The curve $C$ is smooth of genus five and has eight rational points at
$T=0$.  The fixed-point-free involution
\[
 \sigma[P:Q:R:z:T]=[-P:-Q:-R:z:T]
\]
has quotient the smooth plane quartic $D\subset\PP^2_{[X:Y:Z]}$ defined by
\begin{equation}
 F=9801(X^4+Y^4)+19598X^2Y^2
       -39592(X^2+Y^2)Z^2+39984Z^4=0.
 \label{eq:theta-quartic}
\end{equation}
The quotient map is
\[
 \pi:C\longrightarrow D,\qquad [P:Q:R:z:T]\longmapsto[Q:R:P].
\]
In particular, $\pi$ is a connected \emph{\'etale} double cover.  The points
\[
 O_C=[1:1:1:1:0],\qquad O_D=[1:1:1]
\]
are rational and satisfy $\pi(O_C)=O_D$.
\end{proposition}

\begin{proof}
On the chart $T=1$, the radicands in \eqref{eq:theta-C} have
discriminants $-200$, $-8$ and $-8$, and pairwise resultants
$9801$, $9801$ and $39984$. Their six roots are therefore distinct.  A simple zero
of each polynomial shows that its square class is independent of those of
the other two over $\overline\Q(z)$.  The map to the $[z:T]$-line has
degree eight, with six branch values and inertia order two at each.
Riemann--Hurwitz gives
\[
 2g(C)-2=8(-2)+6\cdot4=8.
\]
The simple branch points are smooth.  At $T=0$ we may set $z=1$, and the
points are exactly
\[
 [\epsilon_1:\epsilon_2:\epsilon_3:1:0],\qquad
 \epsilon_i\in\{1,-1\}.
\]
The three derivatives with respect to $P,Q,R$ are nonzero there, so these
points are smooth as well.

There is no point of $C$ with $P=Q=R=0$: subtracting the last two equations
gives $zT=0$, and the equations then force $z=T=0$.  Hence $\pi$ is a
morphism.  A projective fixed point of $\sigma$ would have either
$P=Q=R=0$ or $z=T=0$, both impossible.

On the chart $P\ne0$, set
\[
 \xi=Q/P,\qquad \eta=R/P,\qquad v=T/P,\qquad w=z/P,
 \qquad S=\xi^2+\eta^2.
\]
The three quadrics give
\begin{equation}
 v^2=\frac{S-2}{2},\qquad
 wv=\frac{\xi^2-\eta^2}{28},\qquad
 w^2=51-25S.
 \label{eq:theta-vw}
\end{equation}
Eliminating $v,w$ gives
\begin{equation}
 (\xi^2-\eta^2)^2+9800(S-2)^2-392(S-2)=0,
 \label{eq:theta-elimination}
\end{equation}
which is \eqref{eq:theta-quartic} on $Z=1$.

Here is also a direct smoothness proof for the plane equation.  Write
\[
 a=9801,\quad b=19598,\quad c=-39592,\quad d=39984,
 \qquad
 M=\begin{pmatrix}2a&b&c\\b&2a&c\\c&c&2d\end{pmatrix}.
\]
At a singular point, the rows indexed by the nonzero coordinates
annihilate the corresponding vector of coordinate squares.  The nonempty
principal minors of $M$ have values
\[
 19602,\quad79968,\quad156800,\quad6272,\quad-1229312,
\]
with repetitions.  Every such submatrix is invertible.  This proves
smoothness; a smooth plane quartic is geometrically integral and has
genus three.  The map to it has degree two, since generically $v$ determines
$w$ in \eqref{eq:theta-vw}, and the two choices are interchanged by $\sigma$.
The fixed-point calculation proves that the map is \emph{\'etale}.
\end{proof}

\begin{proposition}
\label{prop:theta-strict-lift}
A rational point $[\xi:\eta:1]\in D(\Q)$ is the image of a point of
$\mathcal U\cap C(\Q)$ if and only if
\begin{equation}
 \xi>0,\qquad\eta>0,\qquad
 v=\sqrt{\frac{\xi^2+\eta^2-2}{2}}\in\Q_{>0},\qquad
 \left|\frac{\xi^2-\eta^2}{28v}\right|<\frac7{11}.
 \label{eq:theta-strict-lift}
\end{equation}
For this choice of $v$, the unique point on the positive branch is
\begin{equation}
 P=\frac1v,\qquad Q=\frac\xi v,\qquad R=\frac\eta v,
 \qquad z=\frac wv,
 \qquad w=\frac{\xi^2-\eta^2}{28v}.
 \label{eq:theta-inverse}
\end{equation}
\end{proposition}

\begin{proof}
Equation \eqref{eq:theta-elimination} gives $w^2+50v^2=1$, and substitution
in \eqref{eq:theta-inverse} gives all three equations of $C$.
Moreover $P\ge\sqrt{50}>1$.  The remaining strict inequality is
\[
 P+7|z|<50\quad\Longleftrightarrow\quad1+7|w|<50v.
\]
Both sides on the right are positive.  Squaring and using
$50v^2=1-w^2$ reduces it to
\[
 99|w|^2+14|w|-49<0,
\]
whose positive root is $7/11$.  The inverse implication follows from the
same calculation.
\end{proof}

Every real point of $D$ lies on $Z\ne0$.  Indeed,
\[
 F(X,Y,0)=9801(X^2+Y^2)^2-4X^2Y^2>0
\]
for a nonzero real pair $(X,Y)$.  The four points where $v^2=0$ are
\begin{equation}
 E_{\epsilon,\eta}=[\epsilon:\eta:1],\qquad
 \epsilon,\eta\in\{1,-1\};
 \label{eq:theta-four-boundary}
\end{equation}
each has two rational preimages at $T=0$ on $C$.

\subsection{The complete bitangent configuration}

Put $r=\sqrt2$, $s=\sqrt{51}$, and $K=\Q(r,s)$.  Independent signs in the
following statement range over $\{1,-1\}$.

\begin{theorem}
\label{thm:theta-lines}
The twenty-eight bitangents of $D$ are precisely the following lines:
\begin{enumerate}
\item $Z=uX$ and $Z=uY$, where $u=\epsilon r/2$ or $u=5\epsilon r$;
\item $Z=(\epsilon hX+\eta hY)/14$, where $h=9$ or $h=11$;
\item $Z=s(\epsilon hX+\eta jY)/102$, where
      $(h,j)=(9,11)$ or $(11,9)$;
\item $X=tY$, where $t=\epsilon r(s+7\eta)/2$.
\end{enumerate}
Each has two distinct contact points.  The associated odd-theta
\emph{\'etale} algebra is
\begin{equation}
 L_{\mathrm{odd}}
 =\Q^8\times\Q(\sqrt2)^4\times\Q(\sqrt{51})^4
          \times\Q(\sqrt2,\sqrt{51}).
 \label{eq:theta-algebra}
\end{equation}
The full Galois action is induced by the independent sign changes of $r$
and $s$ in the displayed lines.
\end{theorem}

\begin{proof}
We first give the contact identities.  These also specify the contact
divisors, since each is the zero divisor of the displayed quadratic on its
line.  For type~1, on $Z=uX$ one has
\begin{equation}
 F|_{Z=uX}=9801(Y^2+\alpha X^2)^2,
 \quad
 \alpha=\begin{cases}
 -1/99,&u^2=1/2,\\
 -9899/99,&u^2=50.
 \end{cases}
 \label{eq:theta-contact-I}
\end{equation}
Exchange $X,Y$ for the other four lines.  For type~2,
\begin{equation}
 \begin{array}{c|c|c}
 h&k_h&q_h\\\hline
 9&-818/81&13122/49\\
 11&1222/121&29282/49
 \end{array}
 \qquad
 F|_{Z=h(\epsilon X+\eta Y)/14}
   =q_h(X^2+\epsilon\eta k_hXY+Y^2)^2.
 \label{eq:theta-contact-II}
\end{equation}
For type~3 the identity is
\begin{equation}
 F|_{Z=s(\epsilon hX+\eta jY)/102}
   =\frac{6534}{17}(X^2+\epsilon\eta k_{h,j}XY-Y^2)^2,
 \quad
 k_{9,11}=-\frac{980}{99},\quad k_{11,9}=\frac{980}{99}.
 \label{eq:theta-contact-III}
\end{equation}
Finally, $t^2=50+7\eta s$ for type~4, and
\begin{equation}
 F|_{X=tY}
  =39984\left(Z^2-\frac{39592(t^2+1)}{79968}Y^2\right)^2.
 \label{eq:theta-contact-IV}
\end{equation}
All these homogeneous quadratics have nonzero discriminant: in
\eqref{eq:theta-contact-I} the coefficient $\alpha$ is nonzero; in
\eqref{eq:theta-contact-II} one has $k_h\ne\pm2$; in
\eqref{eq:theta-contact-III} one has $k_{h,j}^2+4>0$; and in
\eqref{eq:theta-contact-IV} one has $t^2+1\ne0$.

We prove completeness by comparing coefficients of binary quartics.
A line with one vanishing coefficient is of the form $X=a_0Z$,
$Y=a_0Z$, or $X=a_0Y$.  For the first two forms the discriminant condition
for the restricted even quartic to be a square is
\[
 -3136(a_0^2-2)(50a_0^2-1)=0.
\]
For the third form it is
\[
 -6272(a_0^4-100a_0^2+1)=0.
\]
The resulting lines are exactly types~1 and~4.  The three coordinate axes
give nonzero discriminants and hence contribute no line.

Now write a line with all coefficients nonzero as $Z=\alpha X+\beta Y$,
and set $U=\alpha^2$, $V=\beta^2$.  With $a,b,c,d$ as in the proof of
Proposition~\ref{prop:theta-quotient}, its restricted quartic has
coefficients
\begin{align*}
 A_4&=a+cU+dU^2,&
 A_3&=2\alpha\beta(c+2dU),\\
 A_2&=b+c(U+V)+6dUV,&
 A_1&=2\alpha\beta(c+2dV),\\
 A_0&=a+cV+dV^2.
\end{align*}
For a square quartic, $A_4=0$ would imply $A_3=0$.  The latter forces
$U=101/204$, for which $A_4=2/51$.  Thus $A_4,A_0$ are nonzero.
If one of $A_3,A_1$ vanishes, both vanish, giving $U=V=101/204$; at that
pair
\[
 A_2^2-4A_4A_0=\frac{78368326400}{51}\ne0.
\]
In the remaining case, comparison with the square of a quadratic gives
\begin{align}
 &(c+2dU)^2A_0-(c+2dV)^2A_4\notag\\
 &\hspace{22mm}=2458624(U-V)(102U+102V-101)=0,
 \label{eq:theta-square-factor-I}\\
 E&:=A_0(c+2dV)A_2-UV(c+2dV)^3-2A_0^2(c+2dU)=0.
 \label{eq:theta-square-factor-II}
\end{align}
More explicitly, when these four coefficients are nonzero, the restricted
quartic is a square precisely when
\[
 A_4=A_0(A_3/A_1)^2,\qquad
 A_2=\frac{A_1^2}{4A_0}+\frac{2A_0A_3}{A_1}.
\]
The square is then
\[
 A_0\left(\frac{A_3}{A_1}X^2+
                   \frac{A_1}{2A_0}XY+Y^2\right)^2.
\]
Substitution gives exactly
\eqref{eq:theta-square-factor-I}--\eqref{eq:theta-square-factor-II}.
On the two lines singled out by
\eqref{eq:theta-square-factor-I}, the second expression factors as
\begin{align*}
 E|_{V=U}
  &=-1568(196U-121)(196U-81)(204U-101),\\
 E|_{V=101/102-U}
  &=\frac{76832}{17}(68U-27)(204U-121)(204U-101).
\end{align*}
After discarding $U=101/204$, already treated, these yield
\[
 (U,V)=\left(\frac{81}{196},\frac{81}{196}\right),\quad
       \left(\frac{121}{196},\frac{121}{196}\right),\quad
       \left(\frac{27}{68},\frac{121}{204}\right),\quad
       \left(\frac{121}{204},\frac{27}{68}\right).
\]
These are precisely types~2 and~3, and their sufficiency was established
by \eqref{eq:theta-contact-II}--\eqref{eq:theta-contact-III}.

Since $D$ is a smooth plane quartic, its canonical bundle is
$\OO_D(1)$.  A bitangent $\ell$ therefore cuts out $2B_\ell$, where
$B_\ell$ is an effective theta characteristic of degree two.
It has $h^0(B_\ell)=1$: a degree-two pencil would make $D$
hyperelliptic, whereas its canonical map is a plane embedding.
Conversely, an odd theta characteristic has a unique effective divisor
of degree two and gives a bitangent by taking twice that divisor.
The preceding classification is thus the full odd-theta set.

The eight type~2 lines are rational.  Type~1 consists of four quadratic
orbits over $\Q(r)$, and type~3 consists of four quadratic orbits over
$\Q(s)$.  The four type~4 lines form one orbit over $K$: the value of
$t^2$ recovers $s$ up to sign, and $t$ then recovers $r$ up to sign.
This proves \eqref{eq:theta-algebra} and the Galois assertion.
\end{proof}

\subsection{Theta fractions and the literal lifting class}

Choose the rational bitangent
\[
 \ell_0=14Z-9X-9Y.
\]
None of the twenty-eight bitangents contains $O_D$: this follows directly
from their equations (for type~3 the coefficient of $s$ is rational and
nonzero, and for type~4 $t^2=50\pm7s\ne1$).  Thus, for every line in
Theorem~\ref{thm:theta-lines}, the function
\begin{equation}
 f_\ell=\frac{\ell}{\ell_0}\frac{\ell_0(O_D)}{\ell(O_D)}
 \label{eq:theta-fractions}
\end{equation}
is defined at $O_D$ and satisfies
\[
 f_\ell(O_D)=1,\qquad
 \operatorname{div}(f_\ell)=2(B_\ell-B_{\ell_0}).
\]
The line equations and their Galois action therefore give explicitly
labelled theta fractions in each factor of \eqref{eq:theta-algebra}.
Alternatively, relative to the canonical divisor $H=(Z=0)|_D$, the
functions $\ell/Z$ have divisor $2B_\ell-H$.

\begin{proposition}
\label{prop:theta-cover-class}
Let
\[
 C_0=X^2+Y^2-2Z^2,\qquad
 N=97X^2+99Y^2-196Z^2,\qquad
 h=\frac{C_0}{2Z^2},\qquad
 \ell_\pm=Z\pm\frac X{\sqrt2}.
\]
The cover $\pi:C\to D$ is the normalization of $v^2=h$.  If
$E_0=\sum_{\epsilon,\eta}E_{\epsilon,\eta}$ is the sum of the four points
in \eqref{eq:theta-four-boundary}, its geometric covering class is
\[
 \delta=[E_0-H]\in\Jac(D)(\Q)[2]\setminus\{0\}.
\]
Over $\Q(\sqrt2)$ it is identified with
$[B_{\ell_+}-B_{\ell_-}]$ by the exact function-field identity
\begin{equation}
 h=\frac{\ell_+}{\ell_-}
       \left(\frac{N}{28\sqrt2\,Z\ell_+}\right)^2.
 \label{eq:theta-cover-identity}
\end{equation}
\end{proposition}

\begin{proof}
The polynomial identity
\begin{equation}
 N^2=F+784C_0\left(Z^2-\frac{X^2}{2}\right)
 \label{eq:theta-conic-identity}
\end{equation}
gives \eqref{eq:theta-cover-identity}, including its scalar, upon
restriction to $D$.  The function-field description in
\eqref{eq:theta-vw} identifies this cover with $C$.

On the conic $C_0=0$, equation \eqref{eq:theta-elimination} becomes
$(X^2-Y^2)^2=0$.  The conics $C_0=0$ and $X^2-Y^2=0$ meet transversely at
the four points $E_{\epsilon,\eta}$.  They have no common point on $Z=0$.
It follows that
\[
 (C_0=0)|_D=2E_0,\qquad
 \operatorname{div}(h)=2E_0-2H.
\]
Taking divisors in \eqref{eq:theta-cover-identity} now gives the asserted
identity of two-torsion classes.  The latter class is nonzero because the
two distinct contact divisors $B_{\ell_+},B_{\ell_-}$ each have
$h^0=1$; if they were linearly equivalent, they would be equal.
\end{proof}

The zeros of $h$ give useful rational charts on the normalized cover.
Near $O_D$, take $X=1+u$, $Z=1$, and solve the quartic equation for $Y$.
Its first terms are
\begin{equation}
 Y=1-u-\frac{48}{49}u^2+O(u^3),\qquad
 h=\frac{u^2}{49}+O(u^3).
 \label{eq:theta-boundary-expansion}
\end{equation}
Indeed substitution of $Y=1-u+au^2+O(u^3)$ gives the coefficient
$-784a-768$ of $u^2$ in $F$.  In the normalized chart $v=u\omega$, the
two points above $O_D$ have $\omega=\pm1/7$.  The same conclusion at the
other three points follows by the sign changes of $X,Y$.  This also
exhibits directly the eight rational boundary points of the cover.

\subsection{Local survival at every finite descent level}

The rational boundary and the real topology have a consequence independent
of the explicit bitangent calculation.  We formulate it for an arbitrary
abelian variety, so it applies both to $\Jac(C)$ and to the quartic quotient.

\begin{theorem}
\label{thm:theta-local-survival}
Let $X/\Q$ be a smooth projective geometrically integral curve, let
$B\subset X$ be a finite closed subset, and suppose $O\in B(\Q)$.
Let $A/\Q$ be an abelian variety and $a:X\to A$ a morphism with $a(O)=0$.
Let $\mathcal V\subset (X\setminus B)(\R)$ contain a point joined to $O$
by a continuous path in $X(\R)$.
For every integer $m\ge1$, form
\[
 X_m=X\times_{a,A,[m]}A.
\]
The geometrically connected component $X_m^0$ containing $(O,0)$ is
defined over $\Q$.  It has a point above $X\setminus B$ over every
$\Q_p$, and a real point above $\mathcal V$.
\end{theorem}

\begin{proof}
In characteristic zero, multiplication by $m$ on $A$ is finite
\emph{\'etale}.  Consequently $X_m\to X$ is finite \emph{\'etale},
and $X_m$ is smooth and projective.  Its geometric connected components
are disjoint.  The unique one containing the rational point $(O,0)$ is
Galois invariant and therefore descends to $\Q$.

Fix a finite prime $p$.  The $p$-adic inverse function theorem, applied
to the projection at $(O,0)$, identifies a neighborhood of this point
with a neighborhood of $O$ in $X(\Q_p)$.  A local parameter on the smooth
curve takes infinitely many values in every sufficiently small
$p$-adic neighborhood.  We may therefore choose a nearby point outside
the finite set $B$ and lift it to $X_m^0(\Q_p)$.  This argument includes
primes dividing $m$: the differential of $[m]$ is multiplication by the
nonzero element $m$ of the field $\Q_p$.

For the real assertion, let $\gamma$ be a path from $O$ to a point
$P\in\mathcal V$.  The path $a\circ\gamma$ starts at zero and lies in
$A(\R)^0$.  This identity component is a compact connected real Lie
group; its multiplication map $[m]$ is a surjective covering map.
Lift $a\circ\gamma$ starting at zero.  Together with $\gamma$ this gives
a path on $X_m(\R)$ beginning at $(O,0)$, and hence lying in $X_m^0$.
Its endpoint is the required point above $P$.
\end{proof}

\begin{corollary}
\label{cor:theta-surviving-covers}
Put $J_C=\Jac(C)$, $J_D=\Jac(D)$, and
\[
 i_C(P)=[P-O_C],\qquad i_D(S)=[S-O_D].
\]
For every $m\ge1$, the distinguished components of both covers
\begin{equation}
 C\times_{i_C,J_C,[m]}J_C,
 \qquad
 C\times_{i_D\circ\pi,J_D,[m]}J_D
 \label{eq:theta-surviving-covers}
\end{equation}
have nonboundary points at every finite place and a real point whose
image on $C$ belongs to $\mathcal U$ and gives $x\ne y$.
\end{corollary}

\begin{proof}
Take $B=C\cap(T=0)$ and choose the positive square roots in
\eqref{eq:theta-C} as $z$ decreases from $+\infty$ to $1/6$.
All three radicands are positive, since they are
$z^2+50$, $(z+7)^2+2$, and $(z-7)^2+2$.
In projective coordinates this path starts at $O_C$.  At $z=1/6$ its
coordinates are
\[
 P=\frac{\sqrt{1801}}6,\qquad
 Q=\frac{\sqrt{1921}}6,\qquad
 R=\frac{\sqrt{1753}}6.
\]
Here $1<P<8$ and $P+7z<8+7/6<50$, so it belongs to $\mathcal U$.
Formula \eqref{eq:theta-return} gives $y-x=7/300\ne0$.
Apply Theorem~\ref{thm:theta-local-survival} with $a=i_C$ and
$a=i_D\circ\pi$, respectively.
\end{proof}

There is also an elementary uniform description of a finite-prime
square sieve on this source.

\begin{proposition}
\label{prop:theta-finite-primes}
For a fixed $z\in\Q_2$, the three equations
\eqref{eq:theta-C} with $T=1$ have simultaneous roots in $\Q_2$ if and
only if $v_2(z)<0$.  The analogous assertion over $\Q_3$ holds with
$v_3(z)<0$.  For every finite set $S$ of primes, there are infinitely
many rational values of $z$ for which all three radicands are squares
in every $\Q_p$, $p\in S$, and whose positive real roots belong to
$\mathcal U$.
\end{proposition}

\begin{proof}
For integral $z$ at two, $z^2+50$ is $2,3$, or $6$ modulo eight, so
cannot be a square.  Write $z=2^{-k}u$ with $k\ge1$ and $u$ a unit.
For $k\ge2$, division by $z^2$ makes all three radicands units congruent
to one modulo eight.  For $k=1$, multiplication by four gives
\[
 u^2+200,\qquad u^2+28u+204,\qquad u^2-28u+204,
\]
which are again one modulo eight.  The complete square criterion in
$\Q_2$ proves sufficiency.

At three, the residue classes $z=0,1,-1$ respectively make the first,
third, and second radicand a nonsquare modulo three.  If $v_3(z)<0$,
division by $z^2$ gives units congruent to one modulo three.  More
generally, at any odd prime $p$, a negative valuation of $z$ gives this
same unit congruence modulo $p$, so Hensel's lemma supplies all roots.

Set $M=6\prod_{p\in S}p$ and take $z=M^{-k}$ for $k\ge1$.
All the required valuations are negative, while $0<z\le1/6$.
At the real place, the positive roots satisfy $1<P<8$ and
$P+7z<8+7/6<50$.  These distinct rational parameters have all the
asserted properties.
\end{proof}

Proposition~\ref{prop:theta-finite-primes} establishes survival of the
simultaneous square conditions at any specified finite set of places.
Corollary~\ref{cor:theta-surviving-covers} establishes the stronger
statement after imposing a fixed multiplication level on either
Jacobian: a single geometrically connected cover retains nonboundary
points at every finite place and a strict real distance configuration.
For the second cover in \eqref{eq:theta-surviving-covers}, this conclusion
includes the literal quadratic lift identified in
Proposition~\ref{prop:theta-cover-class}.  Thus these finite descent
levels have nonempty local-solubility sets with the complete source
and the strict real inequalities imposed.

\section{Symmetric return maps and quartic pencils}
\label{sec:symmetric-models}

The four distances admit coordinates in which their common squareclass
and the final rational square are separate equations.  We first give the
return map, including its normalization, and then use the two rulings of
the radius quadric to obtain pencils of plane quartics.

Put
\[
 \Phi(u)=u^4-6u^2+1,
 \qquad
 \begin{aligned}
 N_0&=1-h-k-hk,&N_1&=1-h+k+hk,\\
 N_2&=1+h-k+hk,&N_3&=1+h+k-hk.
 \end{aligned}
\]

\begin{theorem}[Symmetric return]
\label{thm:sym-return}
Let $h,k,W\in\Q$ satisfy $W^2=\Phi(h)\Phi(k)$ and set
\[
 A=(1+h^2)(1+k^2),\qquad T=A+W.
\]
If $T\ne0$, the point
\begin{equation}
 x=\frac12-\frac{2k(1-h^2)}T,
 \qquad
 y=\frac12-\frac{2h(1-k^2)}T
 \label{eq:sym-return}
\end{equation}
satisfies the four distance identities $q_j=N_j^2/T$ in the corner
order $00,10,01,11$.

Conversely, let $(x,y)$ be a real point different from the corners,
let $r_j>0$ be its four distances, and put $S=r_0+r_1+r_2+r_3$.
Then
\begin{equation}
 h=\frac{r_2+r_3-r_0-r_1}{S},\qquad
 k=\frac{r_1+r_3-r_0-r_2}{S},\qquad
 T=\frac{16}{S^2},\qquad W=T-(1+h^2)(1+k^2)
 \label{eq:sym-forward}
\end{equation}
satisfy $W^2=\Phi(h)\Phi(k)$, $N_j=4r_j/S$, and
\eqref{eq:sym-return} returns the same point.  If
$r_j\in\sqrt d\,\Q_{>0}$ for one positive squarefree $d$, all four
parameters are rational and $[T]=[d]$.

In particular, the strict rational-distance problem is equivalent to
the existence of rational $h,k,W$ for which
\begin{equation}
 \begin{gathered}
 W^2=\Phi(h)\Phi(k),\qquad
 T=(1+h^2)(1+k^2)+W\in\Q^{\times2},\\
 4|k(1-h^2)|<T,\qquad 4|h(1-k^2)|<T.
 \end{gathered}
 \label{eq:sym-complete-system}
\end{equation}
\end{theorem}

\begin{proof}
The identity
\begin{equation}
 A^2-\Phi(h)\Phi(k)
 =8\bigl(k^2(1-h^2)^2+h^2(1-k^2)^2\bigr)
 \label{eq:sym-discriminant}
\end{equation}
gives the first distance equation on substituting
$xT=T/2-2k(1-h^2)$ and $yT=T/2-2h(1-k^2)$.
Indeed, twice its residual is
$T^2-2AT+A^2-\Phi(h)\Phi(k)=W^2-\Phi(h)\Phi(k)$.
The other three follow from
\[
 N_1^2-N_0^2=4k(1-h^2),\quad
 N_2^2-N_0^2=4h(1-k^2),\quad
 N_0^2+N_3^2=N_1^2+N_2^2.
\]

For the converse, $r_0^2+r_3^2=r_1^2+r_2^2$.  Expanding the
definitions in \eqref{eq:sym-forward} using this relation gives
$S^2N_j=4r_jS$.  For example,
\[
 \frac{S^2N_0}{2}
 =r_0^2+r_1^2+r_2^2-r_3^2+2r_0(r_1+r_2+r_3)
 =2r_0S.
\]
Thus $q_j=N_j^2/T$.  The differences
$q_1-q_0=1-2x$ and $q_2-q_0=1-2y$ give the return map.
Substitution in $q_0=x^2+y^2$ and
\eqref{eq:sym-discriminant} gives $W^2=\Phi(h)\Phi(k)$.
If $r_j=\sqrt d\,\rho_j$ with positive rational $\rho_j$, then
$S=\sqrt d\,s$ for $s\in\Q_{>0}$, so all the asserted
rationalities and $[T]=[d]$ follow.

Finally, when $T=t_0^2\ne0$, the return gives positive rational
distances $|N_j/t_0|$ whenever $0<x,y<1$; an interior point has no
zero distance.  The last two inequalities in
\eqref{eq:sym-complete-system} are exactly $|1-2x|<1$ and
$|1-2y|<1$.  The forward construction proves the reverse implication.
\end{proof}

The excluded case $T=0$ is explicit.  By
\eqref{eq:sym-discriminant}, it requires
$k(1-h^2)=h(1-k^2)=0$, hence
$(h,k)=(0,0)$ or $(h,k)=(\pm1,\pm1)$, with $W=-A$.
Also, $\Phi$ has no rational root.  These observations account for
the exceptional denominators in the return map.

\subsection{The homogeneous squareclass}

For a positive squarefree $n$, consider the intersection
\[
 D_n\subset\PP^3_{[a:b:c:e]}:\qquad
 a^2=e^2+c^2,\qquad nb^2=e^2-c^2.
\]

\begin{proposition}[Normalization of the return]
\label{prop:sym-normalization}
For $n>1$, every rational point of $D_n$ has $a+e\ne0$.
The functions
\[
 h=\frac{c}{a+e},\qquad \beta=\frac{2b}{a+e}
\]
identify its rational points with $n\beta^2=\Phi(h)$.
For two such points the returned squareclass is the class of
\begin{equation}
 T(p,q)=\frac{4(aa'+nbb')}{(a+e)(a'+e')}.
 \label{eq:sym-projective-T}
\end{equation}
This expression is invariant under independent rescaling of both
projective representatives.

The curve has no rational point if $n$ is even.  If $n$ is odd and
$(a,b,c,e)$ is a primitive integral representative with $a>0$, then
\[
 a+e=
 \begin{cases}
 2s^2,&n\equiv1\pmod4,\\
 s^2,&n\equiv3\pmod4
 \end{cases}
 \qquad(s\in\Z\setminus\{0\}).
\]
Consequently, for two primitive integral representatives with
$a,a'>0$, the class in \eqref{eq:sym-projective-T} equals
$[aa'+nbb']$.
\end{proposition}

\begin{proof}
If $a+e=0$, the first quadric gives $c=0$ and the second gives
$nb^2=e^2$.  Since $n>1$ is squarefree, this would be the zero
vector.  Multiplication by $2/(a+e)$ gives the representative
\[
 (a,b,c,e)=(1+h^2,\beta,2h,1-h^2),
\]
because $(a+e)(a-e)=c^2$.  The second quadric becomes
$n\beta^2=\Phi(h)$.  For two normalized representatives,
\eqref{eq:sym-projective-T} is
$(1+h^2)(1+k^2)+n\beta\beta'$, as in
Theorem~\ref{thm:sym-return}.  Its numerator and denominator both
have bidegree $(1,1)$, proving invariance.

If a prime divides $a,c,e$, its square divides $nb^2$ and hence
the prime divides $b$, contradicting primitivity.  Thus
$(a,c,e)$ is a primitive Pythagorean triple: $a$ is odd and the
two legs have opposite parity.  In particular $e^2-c^2$ is odd,
which excludes even $n$ and makes $b$ odd.  When $c$ is even,
the coprime nonnegative integers $(a+e)/2,(a-e)/2$ have square
product $(c/2)^2$, so both are squares; here $n\equiv1\pmod4$.
When $c$ is odd, the coprime positive odd integers $a+e,a-e$
have square product $c^2$, and $n\equiv3\pmod4$.
The product $(a+e)(a'+e')$ is therefore a square.  Finally
$a^2-nb^2=2c^2>0$ at these rational points, so
$aa'>n|bb'|$ and $aa'+nbb'>0$.
\end{proof}

This normalization matters when choosing representatives.  For example,
on $D_{161}$ take
\[
 p=(17,1,8,15),\qquad q=(689,31,400,561).
\]
Then $aa'+161bb'=16704=29\cdot24^2$, whereas replacing $p$ by
$p/16704$ makes this raw value equal to $1$.  In both cases
\eqref{eq:sym-projective-T} is $1044/625=29(6/25)^2$ and
the return is $(49/348,161/696)$.  Its four distances are
$\sqrt{29}(35,115,101,149)/696$.

\subsection{Two pencils of genus-three curves}

The quadric $r_0^2+r_3^2=r_1^2+r_2^2$ also has the ruling
coordinates
\[
 (r_0:r_1:r_2:r_3)
 =(\lambda+\mu:\mu-\lambda:\lambda\mu+1:\lambda\mu-1).
\]
For positive distances with $r_2\ne r_3$, one may take
$\lambda=(r_2+r_3)/(r_0+r_1)$ and
$\mu=(r_0+r_1)/(r_2-r_3)$.  If $s=2/(r_2-r_3)$, the unit-side
condition becomes
\begin{equation}
 \Gamma_\lambda:\quad
 \bigl(2s^2-(1+\lambda^2)(1+\mu^2)\bigr)^2
 =-\Phi(\lambda)\Phi(\mu).
 \label{eq:sym-plane-quartic}
\end{equation}
For a rational point of this curve with $s\ne0$, the return is
\[
 x=\frac12+\frac{2\lambda\mu}{s^2},\qquad
 y=\frac12-\frac{(\lambda^2-1)(\mu^2-1)}{2s^2},
\]
and the four squared distances are the squares of the four displayed
bilinear coordinates divided by $s^2$.  Thus the same strict chamber
selects the original configurations on each fibre.

\begin{theorem}[Quartic pencils and elliptic factors]
\label{thm:sym-quartic-pencil}
For $\lambda\in\Q\setminus\{0,1,-1\}$, the projective
completion of \eqref{eq:sym-plane-quartic} is a smooth plane
quartic.  Put $c_\lambda=-\Phi(\lambda)$ and
\[
 E_{c_\lambda}:\quad y^2=x^3-c_\lambda^2x,
 \qquad
 F_\lambda:\quad v^2=c_\lambda(s^2-4)(s^2-4\lambda^2),
\]
where $(s,v)=(2,0)$ is chosen as origin on $F_\lambda$.
Then there is a $\Q$-isogeny
\[
 \Jac(\Gamma_\lambda)\sim_\Q
 E_{c_\lambda}\times F_\lambda^2.
\]
The other ruling gives a second pencil with the same description.
If the four distances lie in $\sqrt d\,\Q$, the two
$F_\lambda$ factors are replaced by their quadratic twists by $d$.
\end{theorem}

\begin{proof}
Set $A=1+\lambda^2$, $D=(\lambda^2-1)^2$ and
$B=\lambda^4-10\lambda^2+1$.  In the squared projective
coordinates $M=\mu^2$, $S=s^2$, $Z=z^2$, the equation is the conic
\[
 2S^2-2AS(M+Z)+D(M^2+Z^2)-2BMZ=0.
\]
Its symmetric matrix has determinant $-4\Phi(\lambda)^2$.
The discriminants on $M=0$ and $Z=0$ are
$-4\Phi(\lambda)$; on $S=0$ the discriminant is
$-64\lambda^2\Phi(\lambda)$.  They are nonzero, and the
conic contains none of the three coordinate vertices.  Pulling
back by $(\mu:s:z)\mapsto(\mu^2:s^2:z^2)$ therefore gives
a smooth degree-four plane curve.  Equivalently, its degree-four
cover of the conic has six distinct branch points, each of inertia
two, so $2g-2=-8+12=4$.

The involutions changing the signs of $\mu$ and $s$ generate
a Klein four-group.  The quotient by the sign of $s$ is
\[
 C_\lambda:\quad W^2=c_\lambda\Phi(\mu).
\]
It maps nonconstantly to
$E':Y^2=X(X^2-6c_\lambda X+c_\lambda^2)$ by
$(X,Y)=(c_\lambda\mu^2,c_\lambda\mu W)$.
The usual degree-two map
\[
 X_1=\frac{X^2-6c_\lambda X+c_\lambda^2}{X},\qquad
 Y_1=Y\left(1-\frac{c_\lambda^2}{X^2}\right)
\]
has image $Y_1^2=X_1(X_1+4c_\lambda)(X_1+8c_\lambda)$.
Setting $X_1=4(x-c_\lambda)$, $Y_1=8y$ gives
$E_{c_\lambda}$; hence $\Jac(C_\lambda)$ is isogenous to it.

The quotient by the sign of $\mu$ is $F_\lambda$.  Indeed,
with
$v=D\mu^2-B-As^2$, its equation is
$v^2=c_\lambda(s^2-4)(s^2-4\lambda^2)$.
Interchanging the projective coordinates $\mu,z$ preserves the
quartic and conjugates this involution to the third nontrivial
element of the Klein group.  The third quotient is consequently
isomorphic to $F_\lambda$ over $\Q$.
The pullbacks of the three genus-one differentials lie in the three
different nontrivial character spaces.  They span the three-dimensional
space of regular differentials, so the induced homomorphism of
Jacobians is an isogeny.

For a common twist $d$, write the side parameter as $s^2=d t^2$.
The second quotient becomes
$v^2=c_\lambda(d t^2-4)(d t^2-4\lambda^2)$.
Over $\Q(\sqrt d)$ it is identified with $F_\lambda$ by
$s=\sqrt d\,t$.  The nontrivial Galois action changes $s$ to
$-s$, which acts by negation on the Jacobian, as is also seen on
the differential $ds/v$.  This is the quadratic twist by $d$.
\end{proof}

\section{Two-adic independence of the elliptic sources}
\label{sec:source-independence}

The last square in the distance problem imposes a useful independence
condition.  The argument uses the complete rational points, rather than
chosen generators or a saturation computation.
Throughout this section $n>1$ is a positive odd squarefree integer and
\[
 E_n:\quad Y^2=X(X-n)(X+n).
\]
For a point outside $E_n[2]$, write $u=X/n$ and denote the
$2$-adic valuation by $v$.

\begin{lemma}[The dyadic level]
\label{lem:indp-level}
Each point of $E_n(\Q_2)\setminus E_n[2]$ has a unique level
$\ell\geq1$ and lies in one of the following four regions:
\[
\begin{array}{c|c|c}
\text{region}&\text{condition}&v(u-1)\\\hline
R_\infty&v(u)=-2\ell&-2\ell\\
R_0&v(u)=2\ell&0\\
R_+&v(u-1)=2\ell+1&2\ell+1\\
R_-&v(u+1)=2\ell+1&1.
\end{array}
\]
Translation by $E_n[2]$ preserves the level.  Every point is uniquely
the sum of a point of $E_n[2]$ and a point of the formal subgroup
$E_1=\{O\}\cup\{v(X)<0\}$.  This subgroup is torsion-free, and
\[
 \ell(mZ)=\ell(Z)+v(m)
 \quad(Z\in E_1\setminus\{O\},\ m\in\Z\setminus\{0\}).
\]

If $u,u'$ belong to points of different levels and
$(u-1)(u'-1)$ is a square in $\Q_2$, then
\[
 v\bigl(2(u-1)(u+u')\bigr)\quad\hbox{is odd}.
\]
\end{lemma}

\begin{proof}
The normalized equation is $n\eta^2=u(u^2-1)$.
For $v(u)<0$ or $v(u)>0$, parity of the right-hand valuation
forces respectively $v(u)=-2\ell$ or $v(u)=2\ell$.
For a unit $u\equiv1\pmod4$, one has $v(u+1)=1$, so
$v(u-1)$ is odd and at least three.  The case
$u\equiv-1\pmod4$ is symmetric.
The three nontrivial two-torsion translations act on $u$ by
\[
 -\frac1u,\qquad \frac{u+1}{u-1},\qquad
 -\frac{u-1}{u+1}.
\]
They take $R_\infty$ bijectively onto the other three regions,
preserving the indicated level.  This proves the decomposition.

The integral formal parameter $z=-X/Y$ identifies the filtration
on $E_1$ with $v(z)\geq\ell$; its successive quotients have order
two.  The formal addition law has integral coefficients and linear
term $z_1+z_2$, so odd multiplication preserves each exact level.
The duplication formula
\[
 u(2Z)=\frac{(u(Z)^2+1)^2}{4u(Z)(u(Z)^2-1)}
\]
increases the level by one.  Factoring an integer as an odd number
times a power of two proves the formula for $mZ$, and also proves
that $E_1$ has no torsion.  The formal-group assertions used here
are the usual integral expansion at the nonsingular point $O$;
see \cite[Chs.~IV and VII]{silverman2009}.

The square hypothesis in the last assertion places both points
in $R_\infty\cup R_0$, or both in $R_+\cup R_-$.
In the first case, at different levels $v(u+u')$ is even and
$v(u-1)$ is even.  In the second case, if the regions are equal
then $v(u+u')=1$.  If they are opposite, the different levels give
$v(u+u')=2\min(\ell,\ell')+1$.  In either situation
$1+v(u-1)+v(u+u')$ is odd.
\end{proof}

Use the ordered Kummer map
\[
 \delta(P)=([X(P)],[X(P)-n],[X(P)+n]).
\]
It extends at the two-torsion points to the homomorphism with values
\begin{equation}
\begin{array}{c|c}
O&(1,1,1)\\
(0,0)&(-1,-n,n)\\
(n,0)&(n,2,2n)\\
(-n,0)&(-n,-2n,2).
\end{array}
\label{eq:indp-torsion-kummer}
\end{equation}
Here every entry is taken modulo $\Q^{\times2}$.  The homomorphism
property follows from the chord-and-tangent rule: evaluating
$X(X-n)(X+n)-l(X)^2$ at each root shows that the product of the
three corresponding coordinate differences is a square.  The
values at a two-torsion point use the derivative at that root.

\begin{theorem}[Independence criterion]
\label{thm:indp-criterion}
Let $P,P'\in E_n(\Q)\setminus E_n[2]$, and put
$u=X(P)/n$, $u'=X(P')/n$.  Suppose
\[
 \delta(P')=\delta(P)(n,1,n),\qquad u+u'\ne0,
 \qquad v\bigl(2(u-1)(u+u')\bigr)\equiv0\pmod2.
\]
Then $P$ and $P'$ are linearly independent modulo torsion.
\end{theorem}

\begin{proof}
The local decomposition in Lemma~\ref{lem:indp-level} shows that
$E_n(\Q_2)_{\rm tors}=E_n[2]$.  In particular the same is true
over $\Q$.  If $P,P'$ were dependent, the subgroup they generate
together with $E_n[2]$ would have the form
$\Z M\oplus E_n[2]$, where $M$ has infinite order.  Write
$P=aM+T$ and $P'=bM+T'$ with nonzero integers $a,b$.
By \eqref{eq:indp-torsion-kummer}, the class $(n,1,n)$ does not
belong to $\delta(E_n[2])$.  The Kummer relation therefore forces
$b-a$ to be odd: otherwise its right side would be the class of
$T'-T$.  Exactly one of $a,b$ is even.

Write $M=T_M+Z$ locally, with $Z\in E_1\setminus\{O\}$.
The two levels are $\ell(Z)+v(a)$ and $\ell(Z)+v(b)$, which
are different.  The middle Kummer slot implies that
$(u-1)(u'-1)$ is a rational square.  The last assertion of
Lemma~\ref{lem:indp-level} now makes the valuation in the theorem
odd, a contradiction.
\end{proof}

\begin{corollary}[A uniform rank requirement]
\label{cor:indp-original}
Every strict rational-distance point with canonical kernel
$n=m_0m_1m_2m_3$ gives two explicit independent points of
$E_n(\Q)$.  In particular $\rank E_n(\Q)\geq2$.
\end{corollary}

\begin{proof}
Put
\[
 S=r_0+r_3,\quad D=r_3-r_0,\quad
 \Sigma=r_1+r_2,\quad \Delta=r_1-r_2,
 \qquad
 u=\frac{\Sigma^2-2}{D^2},\quad
 u'=\frac{S^2-2}{\Delta^2}.
\]
The denominators are nonzero by Lemma~\ref{lem:int:symmetry}.
The distance equations give
$S^2+D^2=\Sigma^2+\Delta^2$ and
$(S^2-2)(2-D^2)=\Sigma^2\Delta^2$.  Solving for the four
squares gives
\begin{equation}
\begin{aligned}
 S^2&=\frac{2u(u'+1)}{u+u'},&
 D^2&=\frac{2(u'-1)}{u+u'},\\
 \Sigma^2&=\frac{2u'(u+1)}{u+u'},&
 \Delta^2&=\frac{2(u-1)}{u+u'}.
\end{aligned}
\label{eq:indp-four-squares}
\end{equation}
The opposite-pair triangle inequalities make $u,u'>0$, and
\eqref{eq:indp-four-squares} then gives $u,u'>1$.
Also
\[
 t=r_0r_3+r_1r_2=\frac{uu'+1}{u+u'},\qquad
 t^2-1=\frac{(u^2-1)(u'^2-1)}{(u+u')^2}.
\]
Write $t^2-1=nq^2$.  The kernel-product result gives the positive
odd squarefree integer $n>1$.  The identities
\[
 \begin{aligned}
 u(u^2-1)&=(t^2-1)
   \left(\frac{S(u+u')}{D(u'+1)}\right)^2,\\
 u'(u'^2-1)&=(t^2-1)
   \left(\frac{\Sigma(u+u')}{\Delta(u+1)}\right)^2
 \end{aligned}
\]
produce points $P,P'$ of $E_n(\Q)$ with normalized abscissas
$u,u'$.  The scale from $n\eta^2=u(u^2-1)$ to $E_n$ is
$(X,Y)=(nu,n^2\eta)$.

By \eqref{eq:indp-four-squares}, the three numbers
\[
 (u-1)(u'-1),\quad u(u'+1)(u-1),\quad
 u'(u+1)(u-1)
\]
are rational squares.  The first gives the middle Kummer slot.
Multiplying the third by $u(u-1)(u+1)$ gives
$[uu']=[n]$; the product of the three Kummer slots is one,
so their ratio is $(n,1,n)$.
Finally,
\[
 2(u-1)(u+u')=\bigl(\Delta(u+u')\bigr)^2.
\]
Theorem~\ref{thm:indp-criterion} applies.
\end{proof}

The criterion uses the parity of one dyadic valuation.  It applies
to every even-valuation squareclass of that expression, including
some nontrivial common twists.  Its conclusion is an independence
condition on the actual source pair.

\section{Signed sums and denominator patterns}
\label{sec:denominator-patterns}

We refine the prime-power comparisons by retaining the signs of all
four radii at each denominator prime.  Use the primitive normalization
of Section~\ref{sec:introduction}:
\[
 (x,y)=(X/L,Y/L),\quad A_i=Lr_i,\quad
 G=\gcd(A_0,A_1,A_2,A_3),\quad H=L/G^2,\quad a_i=A_i/G.
\]
After exchanging coordinates, $X$ is odd, $4\mid Y$, $24\mid L$,
and the $A_i$ are odd.  Every prime of $G$ is $1$ modulo $4$,
and the $a_i$ are units at primes dividing $H$
(Proposition~\ref{int:content}).  Put
\[
 M_x=L-2X,\quad M_y=L-2Y,\qquad
 e_2=v_2(L)\geq3,\quad j=v_2(M_y)\geq3.
\]
The midline and diagonal exclusions make $M_x,M_y,X-Y,X+Y-L$
nonzero.  Also $v_2(M_x)=1$.

Let $\mathcal P=\{\pm1\}^4/\{\pm(1,1,1,1)\}$ and write
\[
 \alpha=(+,+,-,-),\quad \beta=(+,-,+,-),\quad
 \gamma=(+,-,-,+),\qquad
 \mathcal V=\{1,\alpha,\beta,\gamma\}.
\]
The elements of $\mathcal V$ have an even number of negative signs.
The other coset is represented by $o_i$, with a single negative
sign at corner $i$.  Define, up to an overall sign,
\[
 \Sigma_\epsilon=\sum_i\epsilon_i A_i,
 \qquad \sigma_\epsilon=\Sigma_\epsilon/G,
 \qquad P=A_0A_3-A_1A_2,quad Q=A_0A_3+A_1A_2.
\]
For the displayed representatives, direct expansion gives
\begin{equation}
 \begin{gathered}
 \Sigma_1\Sigma_\gamma=2P=-\Sigma_\alpha\Sigma_\beta,
 \qquad
 \Sigma_{o_0}\Sigma_{o_3}=2Q=\Sigma_{o_1}\Sigma_{o_2},\\
 PQ=-L^2M_xM_y,\\
 (Q^2-L^4)(L^4-P^2)
   =4L^4(X-Y)^2(X+Y-L)^2.
 \end{gathered}
 \label{eq:pat-product-identities}
\end{equation}

\begin{lemma}[Real bounds]
\label{lem:pat-real-bounds}
For a strict rational-distance point, all eight signed sums are nonzero,
and
\[
 \begin{aligned}
 2\sqrt2L&\leq\Sigma_1<(2+\sqrt2)L,\\
 |\Sigma_\alpha|,|\Sigma_\beta|&<\sqrt2L,
 &|\Sigma_\gamma|&<(2-\sqrt2)L,\\
 0<\Sigma_{o_i}&<(2+\sqrt2)L,
 &|P|&<L^2,\qquad L^2<Q<\frac{\sqrt5}{2}L^2.
 \end{aligned}
\]
\end{lemma}

\begin{proof}
The opposite-pair triangle inequalities give
$r_0+r_3\geq\sqrt2$ and $r_1+r_2\geq\sqrt2$.
Convexity of the sum of distances on the closed square gives
$r_0+r_1+r_2+r_3<2+\sqrt2$ in its interior.
Every adjacent-pair sum is strictly between $1$ and $1+\sqrt2$:
the lower bound is the triangle inequality and the upper bound
follows by convexity from its four vertex values.  In the
coordinates of \eqref{eq:sym-forward} this gives
$|h|,|k|<\sqrt2-1$.  With the present sign convention,
\[
 \Sigma_\alpha=-h\Sigma_1,\qquad
 \Sigma_\beta=-k\Sigma_1,\qquad
 \Sigma_\gamma=-hk\Sigma_1.
\]
These prove the stated even-sum bounds.  Each odd sum is positive:
for example $r_3<r_1+1<r_1+r_0+r_2$.
Its upper bound follows from that for the total sum.
Since $M_xM_y\ne0$, \eqref{eq:pat-product-identities} gives
$P\ne0$, and all the even sums are nonzero as well.

Put $\zeta=(x-\tfrac12)+i(y-\tfrac12)$, and write
$\zeta^2=A+iB$.  Then
\[
 \frac Q{L^2}=|\zeta^2-i/2|+|\zeta^2+i/2|\geq1.
\]
Equality lies on a diagonal, which is excluded.  Here is an
elementary upper bound.  Set
$u=(x-\tfrac12)^2$, $v=(y-\tfrac12)^2$, so
$0\leq u,v<1/4$, $A=u-v$ and $B^2=4uv$.
The separately convex polynomial
\[
 16(u-v)^2+\frac{64}{5}uv
\]
has values $0,1,1,4/5$ at the four vertices of
$[0,1/4]^2$, and is strictly less than $1$ in the stated range.
Thus $16A^2+(16/5)B^2<1$.  With
$R=A^2+B^2+1/4=(u+v)^2+1/4<1/2$, this is equivalent to
\[
 2R+2\sqrt{R^2-B^2}<\frac54.
\]
The left side is $(Q/L^2)^2$, proving the upper bound.
Finally, $PQ=-L^2M_xM_y$, $Q>L^2$ and $|M_x|,|M_y|<L$
give $|P|<L^2$.
\end{proof}

\begin{theorem}[Patterns and exact valuations]
\label{thm:pat-valuations}
For every odd prime $p\mid H$, put $e_p=v_p(H)$.  There is a
unique $\epsilon_p\in\mathcal P$ such that
\[
 a_i\equiv\epsilon_p(i)\epsilon_p(0)a_0\pmod {p^{e_p}}.
\]
For $p=2$, let $\epsilon_2(i)\epsilon_2(0)$ be the sign of
$A_i/A_0$ modulo $4$.  The same congruence then holds modulo
$2^{e_2}$.

For odd $p\mid H$, with $\epsilon=\epsilon_p$, the valuations are
\begin{equation}
\begin{aligned}
 v_p(\sigma_\epsilon)&=0,\\
 v_p(\sigma_{\epsilon\alpha})&=e_p+v_p(M_y),\\
 v_p(\sigma_{\epsilon\beta})&=e_p+v_p(M_x),\\
 v_p(\sigma_{\epsilon\gamma})&=2e_p+v_p(M_xM_y),\\
 v_p(\sigma_\eta)&=0\quad(\eta\notin\epsilon\mathcal V).
\end{aligned}
\label{eq:pat-odd-table}
\end{equation}
At $2$ the corresponding five values are, in the same order,
\begin{equation}
 2,\quad e_2+j,\quad e_2+1,\quad 2e_2+j-1,\quad 1.
 \label{eq:pat-two-table}
\end{equation}

Put $P_r=P/G^2$ and $Q_r=Q/G^2$.  If $\epsilon_p$ is even,
then for odd $p\mid H$
\[
 v_p(P_r)=2e_p+v_p(M_xM_y),\qquad v_p(Q_r)=0;
\]
if it is odd, the two values are exchanged.  At $2$ they are
$2e_2+j$ and $1$, respectively, with the same exchange rule.
For $p\mid G$ with $p\nmid H$, both $P_r$ and $Q_r$ are units.
\end{theorem}

\begin{proof}
All six squared differences $a_i^2-a_j^2$ are divisible by $H$.
For odd $p$, their unit roots therefore differ by unique signs
modulo $p^{e_p}$.  At $2$, choose signs so that the four resulting
$b_i$ are congruent to $A_0$ modulo $4$.  Then $v_2(b_i+b_j)=1$.
The horizontal squared differences have valuation $e_2+1$,
and the vertical ones have valuation $e_2+j$; hence the four
$b_i$ are congruent modulo $2^{e_2}$.

For odd $p$, put $b_i=\epsilon_p(i)\epsilon_p(0)a_i$.
All $b_i$ are congruent to $a_0$ modulo $p^{e_p}$.
The sum $S'=b_0+b_1+b_2+b_3$ and
$Q'=b_0b_3+b_1b_2$ are units.  With
$P'=b_0b_3-b_1b_2$ and $T'=b_0-b_1-b_2+b_3$, one has
\[
 S'T'=2P',\qquad P'Q'=-H^2M_xM_y.
\]
This proves the valuation of $\sigma_{\epsilon\gamma}$.
Moreover,
\[
 (b_2-b_0)+(b_3-b_1)
 =H M_y\left(\frac1{b_2+b_0}+\frac1{b_3+b_1}\right).
\]
The bracket is a unit modulo an odd $p$, since it is congruent
to $a_0^{-1}$.  This gives the $\alpha$ entry; the horizontal
version gives the $\beta$ entry.  The other coset consists,
up to sign, of $S'-2b_i$, which are units.

At $2$, the same reciprocal bracket has valuation zero: after
factoring out $1/2$, its two unit terms agree modulo $4$.
Thus the $\alpha$ and $\beta$ entries are $e_2+j$ and
$e_2+1$.  The total sum has valuation two, and the other coset
has valuation one.  The product identities then give
$v_2(P')=2e_2+j$ and $v_2(T')=2e_2+j-1$.

The normalized $P'$ is $\pm P_r$ for an even pattern and
$\pm Q_r$ for an odd one, giving the exchange rule.  Finally,
if $p\mid G$ and $p\nmid H$, then $p$ is split and $X,Y$,
$M_x,M_y$ are units at $p$.  Since
$P_rQ_r=-H^2M_xM_y$, both factors are units.
\end{proof}

\begin{proposition}[Colour determines the nonsplit patterns]
\label{prop:pat-colours}
For an odd prime $p\mid L$ with $p\equiv3\pmod4$,
\[
 \epsilon_p(i)\epsilon_p(j)=\leg{m_im_j}{p}.
\]
For $p\mid L$ with $p\equiv1\pmod4$, all six symbols
$\leg{m_im_j}{p}$ are $1$.  At $2$,
$\epsilon_2(i)\epsilon_2(j)=1$ if and only if
$m_i\equiv m_j\pmod4$.  Consequently, for nonsplit odd $p$
the pattern is even exactly when $\leg np=1$, and
$\epsilon_2$ is even exactly when $n\equiv1\pmod8$.
\end{proposition}

\begin{proof}
Use $U_i=m_iq_i^2$ and $A_i=m_i\Nm(q_i)$ from the Gaussian
normalization.  At an inert prime, the four $U_i$ have the same
nonzero image in $\F_{p^2}$.  If $\mu_i^2=m_i$, then
$\mu_iq_i=\pm\mu_0q_0$.  Frobenius sends
$\mu_i$ to $\leg{m_i}{p}\mu_i$, so taking norms gives
$\leg{m_i}{p}A_i=\leg{m_0}{p}A_0$ modulo $p$.
The sign lift is unique, proving the first assertion.
At a split prime choose the factor of $p$ at which $U_0$ is a
unit.  Reduction in its residue field $\F_p$ gives
$m_i/m_j=(q_j/q_i)^2$.
At $2$, each odd Gaussian norm $\Nm(q_i)$ is $1$ modulo $4$,
so $A_i\equiv m_i\pmod4$.
Taking the product of the four signs proves the parity claims,
using $n\equiv1$ or $7\pmod8$.
\end{proof}

Let $E_+$ be the product of $p^{e_p}$ over odd primes of $H$
with even pattern, and $E_-$ the corresponding product for odd
patterns.  Thus $H=2^{e_2}E_+E_-$.

\begin{theorem}[The two product denominators]
\label{thm:pat-product-denominators}
Write $t=r_0r_3+r_1r_2=N_t/D_t$ and
$z=r_0r_3-r_1r_2=N_z/D_z$ in lowest terms with positive
denominators.  If $n\equiv1\pmod8$, then
\[
 D_t=(2^{e_2}GE_+)^2/2,\qquad D_z=(GE_-)^2;
\]
if $n\equiv7\pmod8$, then
\[
 D_t=(GE_+)^2,\qquad D_z=(2^{e_2}GE_-)^2/2.
\]
In both cases $D_tD_z=L^2/2$ and $\gcd(D_t,D_z)=G^2$.

There are factorizations $n=n_3n_4=n_1n_2$ and positive odd
integers $u,v,u',v'$ such that
\begin{equation}
 2D_t=n_4v^2-n_3u^2,\qquad
 2D_z=n_1u'^2+n_2v'^2,
 \label{eq:pat-norm-denominators}
\end{equation}
and
\[
 n u v u'v'=|(X-Y)(X+Y-L)|,\qquad
 0<\frac{n_3u^2}{n_4v^2}<9-4\sqrt5.
\]
These two factorizations of $n$ are not required to be the
corner-colour factorizations.
\end{theorem}

\begin{proof}
We have $t=Q_r/(G^2H^2)$ and $z=P_r/(G^2H^2)$.
The valuation table cancels precisely the square factors belonging
to the opposite pattern parity.  At a prime shared by $G$ and $H$,
$M_xM_y$ is a unit, so the residual denominator contains exactly
the asserted factor from $G^2$.  At a prime of $G$ outside $H$
both numerators are units.  The dyadic entries are $1$ and
$2e_2+j$, and Proposition~\ref{prop:pat-colours} identifies
which is which.  This proves the denominator formulas, including
their product and gcd.  The numerator opposite an odd denominator
is even; the numerator opposite an even denominator is odd.

The kernel identity gives $t^2-1\in n\Q^{\times2}$.
The last identity of \eqref{eq:pat-product-identities} gives
$1-z^2\in n\Q^{\times2}$.  Both are positive by
Lemma~\ref{lem:pat-real-bounds}.  The positive odd integers
$N_t-D_t,N_t+D_t$ are coprime, as are
$D_z-N_z,D_z+N_z$.  Their products have squarefree part $n$,
so write
\[
 N_t-D_t=n_3u^2,\quad N_t+D_t=n_4v^2,\qquad
 D_z-N_z=n_1u'^2,\quad D_z+N_z=n_2v'^2.
\]
This proves \eqref{eq:pat-norm-denominators}.
Multiplying the four equations and using $D_tD_z=L^2/2$
proves the product formula for $uv u'v'$.
Finally $(t-1)/(t+1)<(\sqrt5-2)/(\sqrt5+2)=9-4\sqrt5$.
\end{proof}

\begin{corollary}[At least three denominator patterns]
\label{cor:pat-three-patterns}
If $G=1$, the primes dividing $L$ realize at least three different
patterns, and both cosets of $\mathcal V$ occur.  In particular,
if only one corner colour is nontrivial, then $L$ has a prime
divisor congruent to $1$ modulo $4$.
\end{corollary}

\begin{proof}
Put $m_p=v_p(M_xM_y)$ for odd $p\mid H$.
If all patterns are even, the table gives
\[
 |P_r|\geq 2^jH^2\prod_{p\mid H,\ p\ {
m odd}}p^{m_p},
 \qquad |P_r|<G^2H^2.
\]
If all are odd, use $Q_r$ and its upper bound
$(\sqrt5/2)G^2H^2$ instead.  Since $2^j\geq8$, $3\mid H$
and $m_3\geq1$, either alternative is impossible for $G=1$.

Suppose there are only two patterns $\epsilon,\epsilon'$ in
different cosets.  Apply the valuation table to
$\sigma_{\epsilon\gamma}\sigma_{\epsilon'\gamma}$.
Each odd prime contributes $2e_p+m_p$, and the dyadic contribution
is $2e_2+j$.  Therefore
\[
 2^jH^2\prod_{p\mid H,\ p\ {
m odd}}p^{m_p}
 \leq|\sigma_{\epsilon\gamma}\sigma_{\epsilon'\gamma}|
 <(2+\sqrt2)^2G^2H^2.
\]
For $G=1$ this contradicts $24>(2+\sqrt2)^2$.

For the last assertion, if $G>1$ a prime of $G$ suffices.
If $G=1$ and no split prime divides $L$, all patterns are
determined by the colours.  When only $m_k$ is nontrivial they
belong to $\{1,o_k\}$, contrary to the first assertion.
\end{proof}

\begin{corollary}[A sharper dyadic comparison]
\label{cor:pat-dyadic-bound}
If $m_i\equiv1\pmod4$ for every corner, and
$K=2^{v_2(L)}$, then
\[
 K^2<\frac{2-\sqrt2}{4}\,L.
\]
\end{corollary}

\begin{proof}
Here $\epsilon_2=1$.  The exact valuation of $\Sigma_\gamma$
is $2e_2+j-1\geq2e_2+2$, so
$4K^2\leq|\Sigma_\gamma|<(2-\sqrt2)L$.
\end{proof}

\section{All rational points of the one-fifth source}
\label{sec:one-fifth-points}

We now determine the rational points of the curves $C,D$ in
\eqref{eq:theta-C} and \eqref{eq:theta-quartic}.  The proof is an
explicit Dem'janenko--Manin comparison of two maps to one elliptic
curve; see also \cite{stoll2011}.  One fixed elliptic group is
certified by rational two-descent and saturation.  The height
comparison and the resulting finite endgame are given explicitly.

\begin{lemma}[The elliptic group]
\label{lem:fifth-group}
For
\[
 E:\quad y^2=x(x-81)(x-121),\qquad G=(49,336),
\]
one has
\[
 E(\Q)=\Z G\oplus E[2](\Q),\qquad
 E[2](\Q)=\{O,(0,0),(81,0),(121,0)\}.
\]
\end{lemma}

\begin{proof}
A two-descent gives a two-Selmer group with invariants $(2,2,2)$;
the torsion invariants are $(2,2)$ and $G$ has infinite order.
Thus the rank is one.  Full saturation gives
$H=(484/9,-8470/27)$, with $H+(121,0)=G$.
Here is the complete two-descent and saturation certificate, in Magma syntax:
\begin{verbatim}
E := EllipticCurve([Rationals() | 0,-202,0,9801,0]);
G := E![49,336,1];
T, tmap := TorsionSubgroup(E);
S, smap := TwoSelmerGroup(E);
assert Invariants(T) eq [2,2];
assert Invariants(S) eq [2,2,2];
assert Order(G) eq 0;
H, proved := Saturation([G] : TorsionFree := true);
assert proved and #H eq 1;
assert exists{t : t in [P @ tmap : P in T] |
             H[1] eq G+t or H[1] eq -G+t};
\end{verbatim}
The local-condition dimensions, starting with the supported
three-coordinate squareclass space, are $15,10,9,7,5,3,3$:
the successive steps are the norm condition and the places
$\infty,3,5,11,2$.  The prime-saturation bound is $7$; the
additional small-height search bound in the full saturation routine is $1379$.
This certificate uses the rational descent and full saturation
algorithms, with no analytic rank or GRH input.
\end{proof}

\begin{theorem}[The complete rational point sets]
\label{thm:fifth-complete-points}
The quartic and its genus-five double cover have rational point sets
\[
 D(\Q)=\{[\epsilon:\eta:1]:\epsilon,\eta\in\{1,-1\}\},
\]
\[
 C(\Q)=\{[\epsilon_1:\epsilon_2:\epsilon_3:1:0]:
                   \epsilon_i\in\{1,-1\}\}.
\]
Consequently no strict rational-distance point has an opposite-pair
distance difference equal to $1/5$ or $-1/5$.
\end{theorem}

\begin{proof}
There is no real point of $D$ with $Z=0$, since the homogeneous
quartic there is positive for $(X,Y)\ne(0,0)$.
Set $\xi=X/Z$, $\eta=Y/Z$ and $c=\xi^2+\eta^2-2$.
When $c=0$, the quartic gives $(\xi^2-\eta^2)^2=0$, so
$(\xi,\eta)=(\pm1,\pm1)$.
Suppose henceforth that $c\ne0$, and put
\[
 z=\frac{\xi^2-\eta^2}{14c},\qquad
 w_+=\frac{2\xi}{c},\qquad w_-=\frac{2\eta}{c}.
\]
Substitution gives
\begin{equation}
 w_+^2=(z^2+50)(z^2+14z+51),\qquad
 w_-^2=(z^2+50)(z^2-14z+51).
 \label{eq:fifth-two-maps}
\end{equation}
The first quartic is birational to $E$ through
\[
 A=w+z^2+7z+26,\qquad
 x=2A+49,\qquad y=(4z+14)A+336.
\]
Its inverse coordinate is
\begin{equation}
 \pi(P)=\frac{y(P)-7x(P)+7}{2(x(P)-49)}.
 \label{eq:fifth-pi}
\end{equation}
This function has simple poles at $O,-G$ and
$\pi(G)=-937/168$ at its removable singularity.
The involution over the $z$-line is $P\mapsto-G-P$.
Thus \eqref{eq:fifth-two-maps} produces rational points $P_+,P_-$
with $\pi(P_+)=z$ and $\pi(P_-)=-z$.

The map $P\mapsto3P+G$ commutes with this involution, and hence
descends to a degree-nine map $R$ on the $z$-line.  Write
$R=\mathcal A/\mathcal B$ with coefficients listed in
Table~\ref{tab:fifth-map}.  Substitution in the elliptic addition
law verifies
\begin{equation}
 R(\pi(P))=\pi(3P+G).
 \label{eq:fifth-triplication}
\end{equation}
Let $h_2([a:b])=\log_2\max(|a|,|b|)$ for coprime integers
$a,b$.  The B\'ezout coefficients displayed below give
\[
 U_0\mathcal A+V_0\mathcal B=D_0,
 \qquad U_\infty\mathcal A^\vee+V_\infty\mathcal B^\vee=D_0,
 \qquad D_0=55399054482258080041175062500,
\]
where $F^\vee(z)=z^9F(1/z)$ and all four B\'ezout polynomials
have degree at most eight.  Their paired coefficient norms, and
the norms of $\mathcal A,\mathcal B$, are less than $2^{124}$.
The two polynomials have no common root.

After homogenization, these identities give $D_0b^{17}$ and
$D_0a^{17}$.  The gcd of
$\mathcal A_h(a,b),\mathcal B_h(a,b)$ therefore divides $D_0$.
Use the identity corresponding to the larger of $|a|,|b|$ to
obtain the lower bound; the coefficient norms give the upper bound.
It follows that
\begin{equation}
 |h_2(R(t))-9h_2(t)|\leq124\qquad(t\in\PP^1(\Q)).
 \label{eq:fifth-height-error}
\end{equation}
The telescoping limit
$\widehat h_R(t)=\lim_{j\to\infty}9^{-j}h_2(R^j(t))$
therefore exists and differs from $h_2(t)$ by at most $124/8$.
The first four iterates of $\infty$ have maximum-coordinate
bit lengths $10,95,860,7754$.  Consequently
\begin{equation}
 h_0:=\widehat h_R(\infty)
 \geq\frac{7753-124/8}{9^4}
 =\frac{15475}{13122}>0.
 \label{eq:fifth-positive-height}
\end{equation}

Choose $A_0$ over a finite extension with $2A_0=G$.
The pole divisor $O+(-G)$ of $\pi$ is linearly equivalent to
$2(-A_0)$.  The standard comparison with the canonical elliptic
height, together with
$(3P+G)+A_0=3(P+A_0)$, gives
\[
 \widehat h_R(\pi(P))=2\widehat h_E(P+A_0).
\]
Here the heights use logarithms to base two; the bounded error in
the divisor-height comparison disappears in the limit.
Lemma~\ref{lem:fifth-group} now yields
\begin{equation}
 \widehat h_R(\pi(mG+T))=(2m+1)^2h_0
 \quad(m\in\Z,\ T\in E[2](\Q)).
 \label{eq:fifth-height-multiple}
\end{equation}
Write $P_+=mG+T$ and $P_-=nG+T'$.
Since $h_2(z)=h_2(-z)$, the height error gives
\[
 4|(m-n)(m+n+1)|h_0\leq31,
 \qquad |(m-n)(m+n+1)|<7.
\]

If the product is zero, the involution $P\mapsto-G-P$
reduces to $n=m$, and a two-torsion translation must send $z$
to $-z$.  The complete translation table is
\[
\begin{array}{c|c|c}
T''&\pi(P+T'')&\text{equation for finite }z\\\hline
O&z&z=0\\
(0,0)&(-z+700)/(14z+1)&z^2=-50\\
(81,0)&(-35z-205)/(4z+35)&z^2=205/4\\
(121,0)&(-35z-305)/(6z+35)&z^2=305/6.
\end{array}
\]
The last two equations fail at $5$, the second fails over $\R$,
and $z=0$ would give $w_+^2=2550$, which fails at $2$.

If the product is nonzero, write $a=m-n$, $b=m+n+1$.
Then $|a|+|b|\leq|ab|+1\leq7$, whence $-4\leq m\leq3$.
The involution replaces $m$ by $-m-1$ without changing $z$,
so it remains to inspect $P_+=mG+T$ with $0\leq m\leq3$.
For $z=A/B$ in lowest terms the second equation of
\eqref{eq:fifth-two-maps} requires
\[
 \mathcal N=(A^2+50B^2)(A^2-14AB+51B^2)
\]
to be an integer square.  The following table, obtained by the
rational addition law on $E$, gives $(p,v_p(\mathcal N),u)$,
where $u$ is the nonzero unit residue modulo $p$.
Every displayed $u$ is a quadratic nonresidue and every valuation
is even.  The missing entry is $z=\infty$.
\[
\begin{array}{c|cccc}
 m&T=O&T=(0,0)&T=(81,0)&T=(121,0)\\\hline
 0&\infty&(3,4,2)&(3,4,2)&(11,2,7)\\
 1&(5,0,3)&(3,4,2)&(3,4,2)&(11,2,2)\\
 2&(5,0,3)&(3,4,2)&(3,4,2)&(11,2,8)\\
 3&(11,0,2)&(3,4,2)&(3,4,2)&(11,2,8).
\end{array}
\]
This excludes every finite $z$, proving the assertion about $D$.

The quotient $C\to D$ is $[P:Q:R:z:T]\mapsto[Q:R:P]$.
Above the four points just found, $Q^2=R^2=P^2$.
The equations of $C$ give $28zT=0$ and $14zT+T^2=0$,
so $T=0$.  The eight stated boundary points are exactly these
fibres.  The strict return \eqref{eq:theta-return} uses $T=1$,
which is impossible.  The square symmetries give both opposite
pairs and both signs of the difference.
\end{proof}

\begin{center}
\small
\refstepcounter{table}\label{tab:fifth-map}
\begin{tabular}{r|r r}
$i$ & $[z^i]\mathcal A$ & $[z^i]\mathcal B$\\\hline
0 & 111607125000 & -12021749375\\
1 & 82021123125 & -6423322500\\
2 & 26108197500 & -825405000\\
3 & 5174999950 & 193235350\\
4 & 768860400 & 85090800\\
5 & 79894350 & 15797250\\
6 & 3064600 & 2089798\\
7 & -389550 & 206808\\
8 & -52500 & 12678\\
9 & -1874 & 336\\
\end{tabular}
\par\smallskip
Table~\thetable: The degree-nine map in \eqref{eq:fifth-triplication}.
\end{center}
\noindent The following two tables list the Bezout coefficients in increasing degree.
Each row gives the coefficients of the two indicated polynomials.
\begin{center}
\scriptsize
\begin{tabular}{r|r r}
$i$ & $[z^i]U_0$ & $[z^i]V_0$\\\hline
0 & -2997498492457923614300 & -27832687036826975619900\\
1 & -1095858706252693720500 & -15753573106648371231300\\
2 & -19762427635601063700 & -3841783019758516568700\\
3 & 51931145761277052296 & -636028100318545578700\\
4 & 12434000961572564136 & -83226348410177026200\\
5 & 1819310074540097136 & -5702224020804565800\\
6 & 210502428872007552 & 220117172270658448\\
7 & 15563169683956800 & 60607066424472960\\
8 & 483327205240320 & 2695699948274880\\
\end{tabular}
\end{center}
\begin{center}
\scriptsize
\begin{tabular}{r|r r}
$i$ & $[z^i]U_\infty$ & $[z^i]V_\infty$\\\hline
0 & -190236680175253051840321031250 & -1060857558315303455561667968750\\
1 & -5532338056063458586310158593750 & -20552009231971099412950192968750\\
2 & -69059947929901090666818928125000 & -41727935009628484047973450781250\\
3 & -580126053198122963653101514843750 & 2117256664470905046006824705468750\\
4 & -3869791759426723943158100156250000 & 26897617537416601772185292900390625\\
5 & -14274207303033339557638111728515625 & 200210932783826079838895782705078125\\
6 & 16848129423821193386849827861328125 & 1177195891501011517791195006591796875\\
7 & 326540805971266624719922115478515625 & 4456602395261102375145481212158203125\\
8 & 765209073668489627429760501708984375 & 7104023056216253086689948486328125000\\
\end{tabular}
\end{center}


\appendix
\section{The distance surface and its Brauer--Manin sets}
\label{sec:surface}

The four-distance equations also define a fixed surface, independent of
the square-free parameter in Section~\ref{sec:elliptic-sources}.
The surface approach is developed in \cite{berry1990}, and its
product-quotient description belongs to the general construction in
\cite{corvajaturchetzannier2025}.  We give the equations and their geometry
explicitly, and then determine the projection of the proper Brauer--Manin
set onto the positive real region.

\subsection{A canonical complete intersection}

Let $S\subset\PP^4_{[A:B:C:D:T]}$ be defined by
\begin{equation}\label{eq:surface-canonical}
 \begin{split}
 Q&=A^2+D^2-B^2-C^2=0,\\
 F&=(T^2+A^2-B^2)^2+(T^2+A^2-C^2)^2-4T^2A^2=0.
 \end{split}
\end{equation}
On $T=1$, the identification with the distance equations is
\begin{equation}\label{eq:surface-return}
 (r_0,r_1,r_2,r_3)=(A,B,C,D),\qquad
 x=\frac{1+A^2-B^2}{2},\qquad
 y=\frac{1+A^2-C^2}{2}.
\end{equation}
Indeed, $F=0$ gives $x^2+y^2=A^2$, and subtraction gives the other
three identities.  Let $X\to S$ be the minimal resolution.  Write $U$
for the open subset $T=1$, $ABCD\ne0$, viewed as an open subset of $X$,
and put
\[
 \mathcal C=\{(x,y,r_0,r_1,r_2,r_3)\in U(\R):
       0<x,y<1,\ r_j>0\text{ for all }j\}.
\]

\begin{proposition}\label{prop:surface-invariants}
The surface $S$ is geometrically integral and has precisely forty
ordinary double points.  Its resolution $X$ is a simply connected
minimal surface of general type, with
\[
 K_X^2=8,\qquad p_g(X)=5,\qquad q(X)=0,\qquad
 c_2(X)=64,\qquad b_2(X)=62.
\]
The geometric Picard group of $X$ is a finitely generated free abelian
group.
\end{proposition}

\begin{proof}
Work first over an algebraically closed field of characteristic zero.
Put $z=x+iy$, $w=x-iy$, and
\[
 (\alpha_0,\alpha_1,\alpha_2,\alpha_3)=(0,1,i,1+i).
\]
The affine distance equations become
\begin{equation}\label{eq:surface-branch-grid}
 r_j^2=(z-\alpha_j)(w-\overline\alpha_j),\qquad 0\le j\le3.
\end{equation}
Their four square classes are independent in
$\overline\Q(z,w)^*/\overline\Q(z,w)^{*2}$: the divisor
$z=\alpha_j$ occurs with odd multiplicity only in the $j$th radicand.
Thus the affine surface is geometrically integral.  The complete
intersection \eqref{eq:surface-canonical} is Cohen--Macaulay and purely
two-dimensional. Its intersection with $T=0$ has dimension one, so every
irreducible component meets the affine chart. There is consequently only
one component, and it is generically reduced. The Cohen--Macaulay property
excludes embedded components, hence $S$ is reduced and geometrically integral.

The branch grid in \eqref{eq:surface-branch-grid} determines all affine
singularities.  Away from the grid the covering is \'etale.  Along a
single branch line, its square root is a local parameter.  At an
intersection $z=\alpha_j$, $w=\overline\alpha_k$ with $j\ne k$, the two
corresponding square roots are independent local parameters.  At an
intersection with $j=k$, the local equation, after extracting square
roots of units, is $r_j^2=uv$.  This is an ordinary double point.
There are eight choices for the other three roots at each of the four
vertices, giving thirty-two affine nodes.  In distance coordinates the
three nonzero roots at a vertex are $\pm1,\pm1,\pm\sqrt2$.

At infinity, write $H_\pm=A^2-B^2\pm i(A^2-C^2)$.  The boundary is
the union of the two intersections $Q=H_\pm=0$ in $\PP^3$.  Each is
smooth: the ratios of the four diagonal coefficients of $H_\pm$ to
those of $Q$ are
\[
 1\pm i,\quad 1,\quad \pm i,\quad 0,
\]
which are pairwise distinct.  Proportional gradients would therefore
leave at most one coordinate nonzero, contrary to $Q=0$.  Their
intersections are exactly
\begin{equation}\label{eq:surface-infinity-nodes}
 [1:\epsilon_1:\epsilon_2:\epsilon_3:0],
 \qquad \epsilon_1,\epsilon_2,\epsilon_3\in\{1,-1\}.
\end{equation}
At any such point set $A=1$, $B=\epsilon_1+b$,
$C=\epsilon_2+c$, $T=t$, and eliminate $D$ using $Q$.
The quadratic term of $F$ is $4(b^2+c^2-t^2)$, which is nondegenerate.
These are eight more nodes.  The preceding smoothness calculation
also exhausts the points at infinity.  In particular $S$ is normal.

Adjunction for a complete intersection gives $\omega_S=\OO_S(1)$.
Ordinary double points are rational and their resolutions are crepant,
so
\[
 K_X=\pi^*\OO_S(1),\qquad K_X^2=\deg S=8.
\]
The Koszul resolution
\[
 0\longrightarrow\OO_{\PP^4}(-6)
 \longrightarrow\OO_{\PP^4}(-2)\oplus\OO_{\PP^4}(-4)
 \longrightarrow\OO_{\PP^4}
 \longrightarrow\OO_S\longrightarrow0
\]
gives $H^1(S,\OO_S)=0$ and $h^0(S,\OO_S(1))=5$.
Rationality of the singularities then gives $q=0$ and $p_g=5$ on $X$.
The canonical divisor is nef and big; adjunction excludes a
$(-1)$-curve.  Noether's formula gives $c_2=12(1-q+p_g)-K_X^2=64$,
and $b_1=2q=0$ gives $b_2=62$.

For simple connectedness, deform the two defining equations slightly
to a smooth complete intersection of type $(2,4)$, which is simply
connected by the Lefschetz hyperplane theorem.  Remove disjoint small
neighbourhoods of the forty nodes.  The smoothing and the resolution
have the same differentiable exterior.  Each boundary is
$\mathbf{RP}^3$.  Both the corresponding Milnor fibre and the
resolution neighbourhood are simply connected: the former retracts
onto a two-sphere, and the latter onto the exceptional projective
line.  Van Kampen's theorem therefore presents both fundamental
groups as the fundamental group of the common exterior modulo the
normal closures of the boundary-group images.  They are the same
quotient.  Thus $X(\mathbf C)$ is simply connected.

Consequently $H_1(X(\mathbf C),\Z)=0$, and the universal coefficient
theorem makes $H^2(X(\mathbf C),\Z)$ torsion-free.  The exponential
sequence and $q=0$ inject $\Pic(X_{\mathbf C})$ into this group.
Finally $\Pic^0(X)=0$ and the N\'eron--Severi theorem give finite
generation of the geometric Picard group.
\end{proof}

\subsection{A fixed product of genus-five curves}

Let $K=\Q(i)$.  Define $C_0/K$ by its function field
\begin{equation}\label{eq:surface-fixed-curve}
 K(C_0)=K(z)(u_0,u_1,u_2,u_3),\qquad u_j^2=z-\alpha_j,
\end{equation}
and let $\overline C_0$ denote its coefficient conjugate.  Both curves
are understood to be smooth and projective.

\begin{proposition}\label{prop:surface-product}
The curve $C_0$ has genus five.  The diagonal sign action of
$G=(\Z/2\Z)^4$ gives a birational identification
\[
 X_K\sim (C_0\times\overline C_0)/G.
\]
Its invariant functions include the four labelled products
$r_j=u_jv_j$, where $v_j^2=w-\overline\alpha_j$.
The Jacobian of $C_0$ is geometrically isogenous to a product of five
elliptic curves: four have $j=128$ and one has $j=1728$.
\end{proposition}

\begin{proof}
The four radicands in \eqref{eq:surface-fixed-curve} are independent
square classes by their valuations at their distinct zeros.  The
covering $C_0\to\PP^1$ has degree sixteen.  Its inertia elements are
the four individual sign changes at the finite branch points and their
sum at infinity.  Each inertia group has order two, so
\[
 2g(C_0)-2=16(-2)+5\cdot8=8.
\]
On the product, the diagonal action fixes $z,w,u_jv_j$.  The field
they generate has degree sixteen over $K(z,w)$ by the square-class
argument for \eqref{eq:surface-branch-grid}.  This is the degree of
the invariant field, since the product field has degree $256$ over
$K(z,w)$.  The invariant field is therefore precisely the distance
function field.

For a nonempty subset $J\subset\{0,1,2,3\}$, the corresponding
character quotient is
\[
 y_J^2=\prod_{j\in J}(z-\alpha_j).
\]
Only the four subsets of size three and the subset of size four
have positive genus; all five quotients have genus one.  The
pullbacks of their differentials give the five distinct character
lines in $H^0(C_0,\Omega^1)$, hence their product induces an isogeny
on Jacobians.  The branch cross-ratios can be chosen as $i$ for the
three-element quotients and $2$ for the four-element quotient.
Substitution in
\[
 j(\lambda)=256\frac{(1-\lambda+\lambda^2)^3}
                         {\lambda^2(1-\lambda)^2}
\]
gives the asserted values.
\end{proof}

\subsection{The positive real component}

We write $\operatorname{Br}$ for the cohomological Brauer group and
$\operatorname{Br}_0$ for the image of the Brauer group of the ground
field.  The rational point used below lies on the boundary
of $U$.

\begin{lemma}\label{lem:surface-real-boundary}
There is a point $P_*\in X(\Q)$ in the same real connected
component as every point of $\mathcal C$.
\end{lemma}

\begin{proof}
At the node $[1:1:1:1:0]$, use the local coordinates $b,c,t$ from
the proof of Proposition~\ref{prop:surface-invariants}.  The
exceptional conic is $b^2+c^2=t^2$.  Take the rational point
$P_*=[-1:0:1]$ on this conic.

Start with a point of $\mathcal C$, fix its $y\in(0,1)$, and let
$x$ tend to $+\infty$, choosing all four positive distance roots.
This is a path in the smooth affine locus because none of the
distances vanishes.  In the chart $A=1$ its coordinates are
$B=r_1/r_0$, $C=r_2/r_0$, $T=1/r_0$.  Rationalization gives
\[
 \frac b t=r_1-r_0=\frac{1-2x}{r_1+r_0}\longrightarrow-1,
 \qquad
 \frac c t=r_2-r_0=\frac{1-2y}{r_2+r_0}\longrightarrow0.
\]
The lifted path extends to $P_*$ on the resolution.
\end{proof}

\begin{proposition}\label{prop:surface-brauer-finite}
The group $\operatorname{Br}(X)/\operatorname{Br}_0(X)$ is finite.
\end{proposition}

\begin{proof}
Resolve the rational map $C_0\times\overline C_0\dashrightarrow X_K$
from Proposition~\ref{prop:surface-product}.  This gives a smooth
projective surface $Y$, birational to the product, and a generically
finite morphism $f:Y\to X_K$ of degree sixteen.
The product theorem of Skorobogatov--Zarhin
\cite[Theorem~A]{skorobogatovzarhin2014}, with its corrected proof
\cite{skorobogatovzarhin2023}, states that, over a field finitely
generated over $\Q$, the map
\[
 \operatorname{Br}(\overline V)^{\Gamma}\oplus
 \operatorname{Br}(\overline W)^{\Gamma}
 \longrightarrow
 \operatorname{Br}(\overline V\times\overline W)^{\Gamma}
\]
has finite cokernel for smooth projective geometrically integral
varieties $V,W$.  The geometric Brauer group of a smooth projective
curve is zero.  Applied to our two curves, this theorem and
birational invariance therefore show that
$\operatorname{Br}(\overline Y)^{\Gamma_K}$ is finite.

Restriction--corestriction on function fields shows that the kernel of
$f^*: \operatorname{Br}(\overline X)\to\operatorname{Br}(\overline Y)$
is killed by sixteen.  This kernel is finite, because the Kummer
sequence places $\operatorname{Br}(\overline X)[16]$ in a quotient of
the finite group $H^2_{\mathrm{\acute et}}(\overline X,\mu_{16})$.
Taking invariant subgroups thus proves that
$\operatorname{Br}(\overline X)^{\Gamma_K}$, and hence
$\operatorname{Br}(\overline X)^{\Gamma_\Q}$, is finite.

By Proposition~\ref{prop:surface-invariants}, $\Pic(\overline X)$
is a finitely generated free abelian group.  Its Galois action factors
through a finite group.  Inflation--restriction and the vanishing of
continuous homomorphisms from a profinite group to a discrete free
abelian group identify its first Galois cohomology with that of a
finite group acting on a lattice; this group is finite.  The
Hochschild--Serre sequence now embeds
$\operatorname{Br}_1(X)/\operatorname{Br}_0(X)$ in
$H^1(\Q,\Pic(\overline X))$, while
$\operatorname{Br}(X)/\operatorname{Br}_1(X)$ embeds in
$\operatorname{Br}(\overline X)^{\Gamma_\Q}$.  Both are finite.
\end{proof}

\subsection{Adelic points with a prescribed strict real component}

For a subgroup $B\subset\operatorname{Br}(U)$, let
$U(\mathbf A_\Q)^B$ denote the points in the usual restricted-product
adelic space orthogonal to $B$ under the Brauer--Manin pairing.  We
identify $\operatorname{Br}(X)$ with its restriction to $U$.

\begin{theorem}\label{thm:surface-brauer-projection}
Every $w\in\mathcal C$ is the real component of an adelic point in
$U(\mathbf A_\Q)^{\operatorname{Br}(X)}$.  Moreover, for every finite
subgroup $B\subset\operatorname{Br}(U)$ and every $w\in\mathcal C$,
there is a point $(P_v)\in U(\mathbf A_\Q)^B$ with $P_\infty=w$.
\end{theorem}

\begin{proof}
First let $B_0$ be a finite subgroup of $\operatorname{Br}(X)$.
Choose a finite set $\Sigma$ of primes such that $X$, $U$ and
$P_*$ extend to a smooth proper model $\mathscr X$, a smooth
geometrically integral open model $\mathscr U$, and a section over
$\Z[1/\Sigma]$, and all the classes in $B_0$ extend to
$\operatorname{Br}(\mathscr X)$.  The section is denoted by the same
symbol $P_*$ here.  Lang--Weil and Hensel's lemma, after enlarging
$\Sigma$, give a point of $\mathscr U(\Z_p)$ for every
$p\notin\Sigma$.  Every class in $B_0$ evaluates to zero at this
point and at the section, since $\operatorname{Br}(\Z_p)=0$.

At each $p\in\Sigma$, the evaluation of a Brauer class on $X(\Q_p)$
is locally constant.  Choose a sufficiently small neighbourhood of
the rational point $P_*$ on which all the evaluations in $B_0$
agree with their values at $P_*$.  This neighbourhood contains a
point of $U(\Q_p)$: in smooth analytic coordinates a proper algebraic
closed subset cannot contain an open ball.  At the real place use
the prescribed $w$.  Evaluation is constant on real connected
components, so Lemma~\ref{lem:surface-real-boundary} gives agreement
with $P_*$ there as well.

These choices form an adelic point of $U$, integral in the single
model $\mathscr U$ outside $\Sigma$.  For each $\beta\in B_0$ its
evaluations coincide place by place with those of
$\beta(P_*)\in\operatorname{Br}(\Q)$.  Global reciprocity
therefore gives orthogonality to $B_0$.

Proposition~\ref{prop:surface-brauer-finite} supplies finitely many
representatives for
$\operatorname{Br}(X)/\operatorname{Br}_0(X)$.  Because the Brauer
group is torsion, they generate a finite subgroup $B_0$.  Applying
the construction to this group gives orthogonality to all of
$\operatorname{Br}(X)$, as constants have zero total invariant.

For the last assertion, use Harari's formal lemma in the form
recalled in \cite{colliothelene_skoro2021}: if $V$ is smooth and
geometrically integral over a number field, $U\subset V$ is a
nonempty open subset, and $B\subset\operatorname{Br}(U)$ is finite,
then an adelic point of $U$ orthogonal to
$B\cap\operatorname{Br}(V)$ can be replaced by one orthogonal to
$B$, preserving its components at any prescribed finite set of
places.  Apply the lemma with $V=X$, starting with the point just
constructed, and preserve the real component.
\end{proof}

Thus the positive real region lies in the projection of the adelic
set for the entire proper Brauer group, and also in the projection
for each finite subgroup of the affine Brauer group.  In the second
assertion the finite subgroup is fixed before the adelic point is
chosen; this is the quantifier order used in the theorem.

\begin{thebibliography}{99}

\bibitem{berry1990}
T.~G.~Berry,
\emph{Points at rational distance from the corners of a unit square},
Ann. Scuola Norm. Sup. Pisa Cl. Sci. (4) \textbf{17} (1990), no.~4, 505--529.
\url{https://www.numdam.org/item/ASNSP_1990_4_17_4_505_0/}.

\bibitem{berry1992}
T.~G.~Berry,
\emph{Points at rational distance from the vertices of a triangle},
Acta Arith. \textbf{62} (1992), no.~4, 391--398.
\href{https://doi.org/10.4064/aa-62-4-391-398}{doi:10.4064/aa-62-4-391-398}.

\bibitem{bremnerulas2016}
A.~Bremner and M.~Ulas,
\emph{Points at rational distances from the vertices of certain geometric objects},
J. Number Theory \textbf{158} (2016), 104--133.
\href{https://doi.org/10.1016/j.jnt.2015.06.011}{doi:10.1016/j.jnt.2015.06.011}.

\bibitem{bruinpoonenstoll2016}
N.~Bruin, B.~Poonen, and M.~Stoll,
\emph{Generalized explicit descent and its application to curves of genus~3},
Forum Math. Sigma \textbf{4} (2016), e6.
\href{https://doi.org/10.1017/fms.2016.1}{doi:10.1017/fms.2016.1}.

\bibitem{cassels1991}
J.~W.~S.~Cassels,
\emph{Lectures on Elliptic Curves},
London Mathematical Society Student Texts, vol.~24,
Cambridge University Press, Cambridge, 1991.
\href{https://doi.org/10.1017/CBO9781139172530}{doi:10.1017/CBO9781139172530}.

\bibitem{corvajaturchetzannier2025}
P.~Corvaja, A.~Turchet, and U.~Zannier,
\emph{Rational distances from given rational points in the plane},
Geom. Dedicata \textbf{219} (2025), article~59.
\href{https://doi.org/10.1007/s10711-025-01019-0}{doi:10.1007/s10711-025-01019-0}.

\bibitem{colliothelene_skoro2021}
J.-L.~Colliot-Th\'el\`ene and A.~N.~Skorobogatov,
\emph{The Brauer--Grothendieck Group},
Ergebnisse der Mathematik und ihrer Grenzgebiete, 3.~Folge, vol.~71,
Springer, Cham, 2021.
\href{https://doi.org/10.1007/978-3-030-74248-5}{doi:10.1007/978-3-030-74248-5}.

\bibitem{faltings1983}
G.~Faltings,
\emph{Endlichkeitss\"atze f\"ur abelsche Variet\"aten \"uber Zahlk\"orpern},
Invent. Math. \textbf{73} (1983), no.~3, 349--366;
erratum, Invent. Math. \textbf{75} (1984), no.~2, 381.
\href{https://doi.org/10.1007/BF01388432}{doi:10.1007/BF01388432};
\href{https://doi.org/10.1007/BF01388572}{doi:10.1007/BF01388572}.

\bibitem{creutz2020}
B.~Creutz,
\emph{There are no transcendental Brauer--Manin obstructions on abelian varieties},
Int. Math. Res. Not. IMRN \textbf{2020} (2020), no.~9, 2684--2697.
\href{https://doi.org/10.1093/imrn/rny098}{doi:10.1093/imrn/rny098}.

\bibitem{lichtenbaum1969}
S.~Lichtenbaum,
\emph{Duality theorems for curves over $p$-adic fields},
Invent. Math. \textbf{7} (1969), no.~2, 120--136.
\href{https://doi.org/10.1007/BF01389795}{doi:10.1007/BF01389795}.

\bibitem{mazur1977}
B.~Mazur,
\emph{Modular curves and the Eisenstein ideal},
Publ. Math. Inst. Hautes \'Etudes Sci. \textbf{47} (1977), 33--186.
\href{https://doi.org/10.1007/BF02684339}{doi:10.1007/BF02684339}.

\bibitem{misme2026}
T.~Misme,
\emph{Computing Theta Characteristics and their Galois action on general curves of moderate genus},
preprint, 2026, arXiv:2610.29449v1.
\url{https://arxiv.org/abs/2610.29449v1}.

\bibitem{silverman2009}
J.~H.~Silverman,
\emph{The Arithmetic of Elliptic Curves}, second ed.,
Graduate Texts in Mathematics, vol.~106,
Springer, New York, 2009.
\href{https://doi.org/10.1007/978-0-387-09494-6}{doi:10.1007/978-0-387-09494-6}.

\bibitem{stoll2011}
M.~Stoll,
\emph{Rational points on curves},
J. Th\'eor. Nombres Bordeaux \textbf{23} (2011), no.~1, 257--277.
\href{https://doi.org/10.5802/jtnb.760}{doi:10.5802/jtnb.760}.

\bibitem{skorobogatovzarhin2014}
A.~N.~Skorobogatov and Y.~G.~Zarhin,
\emph{The Brauer group and the Brauer--Manin set of products of varieties},
J. Eur. Math. Soc. \textbf{16} (2014), no.~4, 749--768.
\href{https://doi.org/10.4171/JEMS/445}{doi:10.4171/JEMS/445}.

\bibitem{skorobogatovzarhin2023}
A.~N.~Skorobogatov and Y.~G.~Zarhin,
\emph{Corrigendum to ``The Brauer group and the Brauer--Manin set of products of varieties''},
J. Eur. Math. Soc. \textbf{25} (2023), no.~7, 2919--2925.
\href{https://doi.org/10.4171/JEMS/1364}{doi:10.4171/JEMS/1364}.

\end{thebibliography}

\end{document}
